\documentclass[12pt,a4paper]{article}
\usepackage[margin=1in]{geometry}
\usepackage{amsmath,amssymb,amsthm}
\usepackage{microtype}
\setlength{\emergencystretch}{2em}
\usepackage{graphicx}
\usepackage{booktabs}
\usepackage{longtable}
\usepackage{array}
\usepackage{multirow}
\usepackage{natbib}
\usepackage[colorlinks=true,linkcolor=blue,citecolor=darkgray,urlcolor=blue]{hyperref}
\usepackage{tcolorbox}
\usepackage{enumitem}
\usepackage{setspace}
\usepackage{float}
\usepackage{caption}
\usepackage{subcaption}
\usepackage{xcolor}
\usepackage{mathtools}
\usepackage{bm}
\usepackage{algorithm}
\usepackage{algpseudocode}

% Theorem environments
\newtheorem{assumption}{Assumption}
\newtheorem{proposition}{Proposition}
\newtheorem{definition}{Definition}
\newtheorem{remark}{Remark}
\newtheorem{lemma}{Lemma}
\newtheorem{corollary}{Corollary}
\newtheorem{prediction}{Prediction}

% AER-style citations
\bibpunct{(}{)}{;}{a}{}{,}

\title{\textbf{Information Depreciation and Optimal Depth in AI Delegation Chains}}
\author{[Authors Redacted for Peer Review]}
\date{May 29, 2026}

\begin{document}

\maketitle

\begin{abstract}
\noindent When AI agents delegate to other AI agents, information depreciates at each processing layer. We show that this information depreciation---combined with attention scarcity---bounds organizational depth. We embed rational inattention into a multi-layer hierarchy: each agent allocates scarce attention across upstream signals subject to a finite budget (context window). The endogenous transmission factor emerges from Bayesian updating rather than being assumed. We prove that optimal depth is finite and interior, that precision loss compounds geometrically through the chain, and that front-loaded budget allocation---larger context windows at early layers---maximizes output quality. When API costs vanish, the bound on depth persists because finite context windows ensure that each layer must compress information, and that compression compounds. We provide illustrative parameterization of the model using LLM benchmark data and propose five testable predictions with specific regression specifications and identification strategies.
\end{abstract}

\vspace{1em}
\noindent\textbf{JEL Codes:} D23, D82, L22, L86, O32, O33, C67

\vspace{2em}
\section{Introduction}
\label{sec:intro}

Enterprise AI systems increasingly rely on multi-agent architectures, where complex tasks are decomposed across chains of specialized language models. A growing empirical literature documents the performance of these systems---yet the fundamental trade-off governing their optimal depth remains theoretically unexplored. As context windows expand by orders of magnitude (from GPT-3's 2K tokens to Gemini-1.5's 1M tokens) and API costs continue to decline, the natural intuition is that deeper chains---more layers of specialized processing---should become optimal. We show that this intuition is incomplete: even with costless communication, endogenous constraints on information transmission bound the efficient depth of the chain.

The core economic mechanism is \emph{information depreciation}---the gradual loss of signal quality as information traverses a processing hierarchy. In a multi-agent chain, each layer receives noisy signals from the layer below, updates beliefs, and transmits a summary upward. Because each agent's context window is finite, not all incoming signals can be attended to with full precision; some are compressed, some are discarded, and the resulting transmission is necessarily coarser than the input. This degradation compounds across layers, so that the signal reaching the terminal layer carries only a fraction of the original information content. Information depreciation thus operates analogously to capital depreciation: each processing layer is an investment in refined decision-making, but that investment suffers a depreciative loss before it reaches the next stage.

We formalize this intuition by embedding rational inattention \citep{sims2003implications} into a multi-layer information-processing hierarchy. The model delivers two core structural results. First, the optimal budget allocation across layers is \emph{front-loaded}: the largest attention budget should be deployed at the earliest layer, with monotonically decreasing allocations at higher layers (Proposition~\ref{prop:front_loading}). This is the model's most empirically testable prediction---it can be evaluated using only observable architectural choices and does not require calibrating the unobserved depreciation parameter. In practice, most production multi-agent systems already use this architecture: the largest model handles the initial query understanding and retrieval, while smaller, faster models handle intermediate reasoning steps. Our model provides the first theoretical justification for this engineering practice.

Second, two conditions jointly generate finite optimal depth. Signals depreciate in quality as they traverse the chain, formalized as an \emph{endogenous} retention factor $\eta_\ell \in (\underline{\eta}, \bar{\eta})$ that captures the natural decay of information during processing at layer $\ell$ (Proposition~\ref{prop:depth_finite}). The preservation rate $\eta_\ell$ is determined by the average attention allocated per signal at layer $\ell$ through Assumption~\ref{ass:endogenous_eta}, with the key property that even unlimited attention cannot push preservation above a technological bound $\bar{\eta} < 1$. Each agent also faces a finite Shannon-capacity budget for attention allocation, which determines \emph{where} in the interval $[\underline{\eta}, \bar{\eta}]$ the actual preservation rate lies. The depth of the chain is bounded because the compounding of these two forces---inherent algorithmic noise (captured by $\bar{\eta} < 1$) and endogenous attention loss---eventually renders marginal layers unproductive. Either condition alone is insufficient: with $\bar{\eta}=1$ and no inherent degradation, information could circulate without bound regardless of attention budgets; without the bound $\bar{\eta} < 1$, unlimited attention would suffice to maintain signal quality at any depth.

Our analysis is related to the canonical models of multi-layer organization developed by \citet{garicano2000hierarchies} and \citet{dessein2006adaptive} (henceforth, GHM). In the GHM framework, finite organizational depth arises from the interplay of diminishing returns and increasing marginal costs of communication or effort. While information loss in hierarchies is a familiar theme in organizational economics \citep[e.g.,][]{sah1991fallibility,dessein2002authority,qian1994incentives}, existing models typically treat communication fidelity as an exogenous parameter or as the outcome of explicit incentive problems. Our model \emph{complements} the GHM framework by adding a micro-foundation for information processing: the transmission factor $\rho_\ell$, defined as the ratio of achieved posterior precision to first-best posterior precision, emerges endogenously from the optimal attention-allocation problem and reflects the actual information degradation at each layer. This connects the organizational model to the Bayesian micro-foundations of rational inattention \citep{sims2003implications} and generates predictions about how attention weights, signal heterogeneity, and capacity constraints jointly determine organizational form.

This paper makes three contributions. First, we introduce the concept of \emph{information depreciation} into the theory of organizational hierarchies. By modeling signal quality as a depreciating asset that traverses a processing chain, we provide a tractable framework that links communication-loss models in organizational economics \citep{sah1991fallibility,dessein2002authority} to the rational-inattention literature \citep{sims2003implications}. We \emph{endogenize} the depreciation rate by making it depend on the attention allocation through Assumption~\ref{ass:endogenous_eta}, so that the preservation parameters $(\underline{\eta}, \bar{\eta}, a)$ discipline the model's predictions and can be calibrated from observable features of LLM systems.

Second, we provide a mapping from LLM architecture parameters---context window length, token-level precision, and API pricing---to predictions about optimal organizational depth and budget allocation. The model generates sharp comparative statics: depth increases with context window size and decreases with signal depreciation; the optimal allocation is front-loaded regardless of cost structure; and API costs affect the intensive margin of processing but are largely irrelevant for the extensive margin of depth once costs are below a threshold. The front-loading prediction is especially amenable to empirical testing because it requires only observable architectural data---the model sizes or context-window lengths assigned to each layer---and can guide system-design decisions in enterprise AI deployment.

Third, we calibrate the model using publicly available data on benchmark accuracy, context window capacities, and API pricing for state-of-the-art language models. The calibration demonstrates that the key predictions hold for empirically plausible parameter values and generates quantitative predictions about the range of efficient depth under current and projected technology. We provide comparative statics that can discipline future empirical work on multi-agent AI architectures.

The remainder of the paper is organized as follows. Section~\ref{sec:literature} reviews the related literature. Section~\ref{sec:model} presents the model. Section~\ref{sec:single} analyzes the single-layer attention-allocation problem and derives the endogenous transmission factor $\rho_\ell$. Section~\ref{sec:multi} extends the analysis to multiple layers, establishes that the optimal budget allocation is front-loaded, and characterizes the finite optimal depth. Section~\ref{sec:calibration} calibrates the model using LLM technology data and presents comparative statics. Section~\ref{sec:empirics} discusses testable implications for the design of multi-agent systems. Section~\ref{sec:extensions} considers extensions, and Section~\ref{sec:conclusion} concludes.

\section{Related Literature}
\label{sec:literature}

Our work bridges four literatures: organizational economics, rational inattention, AI and organizations, and information design.

\subsection{Organizations and Knowledge Hierarchies}

The organizational economics literature on knowledge hierarchies and multi-layer contracting studies how information loss and control loss bound depth. \citet{garicano2000hierarchies} provides the canonical model of knowledge-based hierarchies, where problems are routed to agents with appropriate expertise. \citet{melumad1995hierarchical} analyze how control loss in hierarchies creates a trade-off between specialization and coordination. \citet{lash2008information} explicitly models control loss as a constraint on depth. \citet{bolton1994firm} study the firm as a communication network. \citet{radner1993organization} analyze decentralized information processing. Our contribution is to microfound the distortion mechanism in attention allocation rather than assuming it.

\subsection{Rational Inattention}

The rational inattention literature studies how information-processing constraints shape decision-making. \citet{sims2003implications} introduces the rational inattention framework, where agents face an entropy constraint on information processing. \citet{sims2006rational} applies this to monetary policy. \citet{mackowiak2023rational} provides a comprehensive review. \citet{dessein2006adaptive} study organizational attention in a two-layer setting, showing how attention scarcity shapes delegation. \citet{dessein2016rational} extend this to communication in organizations. \citet{caplin2015revealed} develop revealed-preference tests for rational inattention. Our model extends the organizational attention framework to arbitrarily deep chains.

\subsection{Multi-Stage Information Design}

A growing literature studies information design in dynamic and multi-stage settings. In these environments, a sender strategically controls the flow of information to a receiver (or receivers) who take actions over time. Our paper shares with this literature the fundamental premise that information is transmitted and evolves across multiple stages, but differs in emphasizing constrained information \emph{processing}---arising from attention scarcity within hierarchical organizations---rather than strategic information \emph{withholding} by a sender.

\citet{ely2017beeps} studies a dynamic information transmission problem in which a sender communicates the state of a Markov process to a receiver via binary signals (``beeps''). The optimal signaling strategy involves selective disclosure: the sender pools states and releases information only when it is most consequential for the receiver's decision. The key insight is that the \emph{timing} of information revelation is a strategic choice, and the sender may optimally delay or suppress information to influence the receiver's beliefs. Our framework differs in that the constraint on information transmission arises not from a sender's strategic objectives but from the organizational limitation that each processing layer has finite attention capacity.

\citet{renault2017optimal} analyze optimal dynamic information provision in a sender-receiver framework with a persistent binary state. They characterize the sender's value function and show that the optimal policy involves gradual revelation: information is released piecemeal over time to maximize the sender's expected payoff. The dynamics of belief updating are central to their analysis, as the sender faces a trade-off between revealing information early and withholding it to preserve future influence. In our model, similar intertemporal trade-offs arise, but they are driven by organizational rather than strategic considerations.

\citet{escude2023slow} study ``slow persuasion,'' in which a sender provides public information over time subject to a graduality constraint. They show that when information can only accumulate gradually, the sender's ability to persuade is inhibited relative to the unconstrained benchmark. The graduality constraint in their model is exogenous, much as the attention constraint in our framework is determined by the organizational technology. The key difference is that their constraint limits the \emph{speed} of information acquisition, whereas ours limits the \emph{volume} of information that can be processed at each layer.

\citet{smolin2021dynamic} studies a principal who designs a dynamic evaluation policy for an agent of uncertain productivity. Smolin shows that the equilibrium is Pareto efficient and that the wage deterministically grows with tenure. The evaluation policy serves as a dynamic information revelation mechanism. Our model shares the feature that information structures are chosen dynamically, but in our setting the designer chooses the \emph{architecture} of information processing rather than the informativeness of signals per se.

The common thread across these papers is that information is a flow variable that must be managed over time. Our contribution is to study how this flow is governed by organizational structure: in a hierarchy with attention-constrained agents, the optimal architecture determines not only \emph{what} information is transmitted but also \emph{how much} of it survives at each stage.

\subsection{AI and Organizations}

A substantial empirical and theoretical literature documents the effects of artificial intelligence on productivity, labor markets, and the organization of work. Our paper complements this literature by examining the \emph{internal} structure of AI systems---how multiple AI agents should be organized into hierarchical architectures---rather than the effects of AI on human labor markets.

\citet{acemoglu2018race} develop a task-based framework in which automation technologies displace labor in tasks previously performed by workers. Their key insight is that while automation directly reduces labor demand, it also indirectly encourages the creation of new tasks in which labor has a comparative advantage. \citet{acemoglu2022tasks} extend this framework to study wage inequality, showing that a substantial fraction of the rise in US wage inequality since 1980 can be attributed to automation of routine tasks. Our paper takes a different perspective: rather than studying how AI substitutes for or complements human labor in existing organizational structures, we ask how AI agents themselves should be organized into efficient architectures for processing information.

\citet{acemoglu2022ai} study the impact of AI on labor markets using evidence from online vacancies. They document that AI-related skills are increasingly demanded in occupations that combine cognitive and technical tasks, and that AI adoption is associated with shifts in skill requirements rather than pure displacement.

\citet{noy2023experimental} provide experimental evidence that generative AI tools substantially increase worker productivity. In a randomized controlled trial with professional writers, they find that access to ChatGPT reduces task completion time by 40\% and increases output quality. Importantly, the productivity gains are largest for less-skilled workers, suggesting that AI can compress skill-based wage differentials. \citet{peng2023impact} report similar findings for software developers using GitHub Copilot, documenting a 55.8\% reduction in task completion time.

\citet{ide2025ai} embed AI into the knowledge-hierarchy framework of \citet{garicano2000hierarchies} but without attention scarcity. Recent work studies AI-driven organizational flattening, but typically abstracts from incentive constraints. We directly address the gap between AI capabilities (expanding context windows) and organizational limits (compounding attention loss).

Our work is complementary to this literature in the following sense. The existing literature studies the \emph{macro} effects of AI adoption on human workers and firms---how AI changes the demand for skills, the distribution of wages, and the productivity of existing organizational forms. In contrast, we study the \emph{micro} design of AI systems themselves: given that AI agents have heterogeneous capabilities and attention constraints, what is the optimal hierarchical architecture for information processing?

\subsection{Information Design}

The information design literature provides tools for modeling strategic communication. \citet{kamenica2011bayesian} introduce Bayesian persuasion, where a sender designs a signal to influence a receiver's beliefs. \citet{lien2019hierarchies} extend this to multi-layer communication networks. Our model differs in that the ``design'' is not strategic but constrained by attention scarcity.

\begin{table}[htbp]
\centering
\small
\setlength{\tabcolsep}{5pt}
\caption{Literature Mapping}
\label{tab:literature}
\begin{tabular}{@{}>{\raggedright\arraybackslash}p{3.0cm}>{\raggedright\arraybackslash}p{3.6cm}>{\raggedright\arraybackslash}p{3.8cm}@{}}
\toprule
\textbf{Contribution} & \textbf{Relation to Our Model} & \textbf{Differentiation} \\
\midrule
\citet{garicano2000hierarchies} & Knowledge hierarchy & + attention scarcity \\
\citet{dessein2006adaptive} & Organizational attention & + continuous allocation + chain degradation \\
\citet{hadfield2021principal} & Single-layer PA in AI & + chain extension + cascade mechanism \\
\citet{phelps2023models} & Behavioral validation with LLMs & + multi-layer experimental design \\
\citet{yang2025ten} & Heuristic: coordination costs & Formal proof: alignment costs persist \\
\bottomrule
\end{tabular}
\end{table}

Table \ref{tab:literature} summarizes how our work relates to the most closely connected papers in each literature.

\section{Model Setup}
\label{sec:model}

\subsection{Overview: From Capital Depreciation to Information Depreciation}
\label{subsec:capital_to_information}

The central premise of this paper is that information, like physical capital, depreciates as it flows through an organizational chain. Each processing layer---each algorithmic transformation, each forwarding step, each summarization---introduces noise, compression, and loss. The rate of this \emph{information depreciation} is not purely exogenous; it is endogenously determined by how attention is allocated across the upstream signals at each layer. Our model treats the delegation chain as a sequence of information-processing stages in which the scarce resource is attention capacity, and the central organizational problem is to design an attention architecture that maximizes the value of the final output.

This framing builds on a direct analogy to capital theory. In the neoclassical growth model, capital depreciates at rate $\delta$ per period: $K_{t+1} = (1-\delta) K_t + I_t$. In our model, signal precision depreciates at rate $(1-\eta_\ell)$ per layer:
\begin{equation}\label{eq:information_depreciation_analogy}
\tau^{0}_{\ell+1,k} = \eta_\ell \cdot \tau^{0}_{\ell,k}, \quad \eta_\ell \in (0,1),
\end{equation}
where $\eta_\ell$ is the \emph{endogenous information preservation rate} at layer $\ell$. Just as the depreciation rate $\delta$ determines how much capital must be replaced to maintain productivity, $\eta_\ell$ determines how much attention must be invested at each layer to preserve signal quality. Crucially, unlike physical capital depreciation, $\eta_\ell$ is not a primitive technological constant: it is an \emph{endogenous} outcome of the attention-allocation problem, determined by the average attention per signal at layer $\ell$ through Assumption~\ref{ass:endogenous_eta} below. The optimal depth of the chain is determined by the point at which the marginal attention cost of adding another layer exceeds the marginal information gain from deeper processing.

This perspective differs fundamentally from the classical principal-agent framework \`{a} la \citet{holmstrom1991multitask}. In the HM model, the agent chooses a scalar effort level $e$ that maps into output through a production function $q = f(e) + \varepsilon$. The distortion is moral hazard: the agent would prefer to set $e = 0$. In our framework, the agent chooses a vector of attention weights $\boldsymbol{\alpha} = (\alpha_1, \ldots, \alpha_K)$ subject to a budget constraint $\sum_k \alpha_k \leq A$. The distortion is information depreciation compounded by rational inattention: even if the agent fully exhausts the budget $A$, the preservation rate $\eta_\ell$ remains bounded above by $\bar{\eta} < 1$, so signals are always processed with some loss. The agent is not shirking; she faces a technological constraint---inherent algorithmic noise in LLM processing---that no amount of attention can fully overcome.

This distinction has three immediate implications for the study of AI delegation chains.

\textbf{First}, it provides a direct micro-foundation for the attenuation of signal quality across layers. The endogenous preservation rate $\eta_\ell$ captures, in reduced form, the algorithmic processes---tokenization, embedding projection, attention-weight normalization, and feed-forward transformations---that degrade signal fidelity in each forward pass of a transformer. Rather than treating precision as exogenous or arbitrarily layer-dependent, we derive its evolution from the attention allocation through Assumption~\ref{ass:endogenous_eta}.

\textbf{Second}, it endogenizes the effective precision parameter. In the HM framework, the signal precision $\tau$ is an exogenous parameter that the principal might manipulate through incentives. In our framework, the effective precision $\tau^{*}_\ell$ at layer $\ell$ emerges from the interplay of four primitives: the base preservation parameters $(\underline{\eta}, \bar{\eta})$, the attention budget $A_\ell$, the attention technology $g_\ell(\cdot)$, and the saturation parameter $a$. The principal influences $\tau^{*}_\ell$ not by writing contracts but by designing the attention architecture---choosing the depth of the chain, the width of context windows, and the allocation of processing capacity across layers. Because $\eta_\ell$ responds to $A_\ell$, larger context windows not only process more signals but also \emph{preserve each signal better}, a double dividend that the exogenous-$\eta$ framework cannot capture.

\textbf{Third}, it structures the output function as the outcome of an information-transmission technology. Equation \eqref{eq:effective_precision} below derives precision from attention allocation given depreciated signal quality. The production function is not assumed; it follows from the technology of selective attention applied to successively degraded signals.

\paragraph{Relationship to Transformer Architectures.} Our notion of an ``attention budget'' $A_\ell$ is deliberately stylized. It is most directly interpretable as a simplified representation of the \emph{context-window length} of a transformer-based LLM---the maximum number of tokens that can be processed in a single forward pass (e.g., 4,096 tokens for GPT-3.5, 128,000 for GPT-4, or 1,000,000 for Gemini 1.5 Pro). However, we emphasize that this mapping is a deliberate simplification. Actual transformer attention involves \emph{soft} allocation through learned query-key-value interactions, \emph{multi-head} parallel processing, and \emph{causal masking} constraints that make the effective attention pattern highly non-linear and context-dependent. Our linear budget constraint \eqref{eq:budget_constraint} abstracts from these complexities to deliver analytical tractability. We return to this limitation in the discussion below.

Our framework builds on the rational-inattention literature initiated by \citet{sims2003implications, sims2006rational} and its organizational extension by \citet{dessein2006adaptive} and \citet{dessein2016rational}. The key departure from Sims is that our agents operate in a hierarchy where the output of one agent's attention-allocation problem becomes the (depreciated) input to the next agent's problem. The key departure from Dessein--Santos is that our agents are artificial---their attention budgets, costs, and signal structures are directly measurable from API logs---which allows us to calibrate and test the model with unprecedented precision.

\paragraph{Limitations and Scope Conditions.} We acknowledge two important limitations of our stylized representation. First, the linear budget constraint \eqref{eq:budget_constraint} does not capture the non-linearities of actual transformer attention, where the effective processing capacity depends on the interaction between query complexity, key dimensionality, and the attention-head structure. A more micro-founded approach would model attention through the softmax operation directly; we sacrifice this realism for analytical clarity. Second, while Assumption~\ref{ass:endogenous_eta} endogenizes the preservation rate $\eta_\ell$ through the attention budget $A_\ell$, we assume a uniform functional form across layers and signals. In practice, different types of signals (text, numerical, structured) may depreciate at different rates, and the parameters $(\underline{\eta}, \bar{\eta}, a)$ may vary with the model architecture and task complexity. Heterogeneous depreciation is discussed in Section \ref{sec:extensions}.


\subsection{The Environment}
\label{subsec:environment}

\begin{assumption}[Hierarchy]\label{ass:hierarchy}
There are $L+1$ layers indexed by $\ell = 0, 1, \ldots, L$. Layer $0$ is the principal (human or top-level AI) who sets the organizational architecture. Layers $\ell = 1, \ldots, L$ are AI agents arranged in a chain: the output of layer $\ell-1$ serves as the input to layer $\ell$. Layer $L$ is the final executor that produces the observable output.
\end{assumption}

\begin{assumption}[State and Signals with Information Depreciation]\label{ass:state_signals}
The underlying state is $\omega \in \Omega \subseteq \mathbb{R}^d$, representing the correct response to a complex task (e.g., the correct answer to a coding problem, the accurate diagnosis from medical records, or the valid proof of a mathematical theorem). At each layer $\ell$, the agent receives $K_\ell$ signals $\{s_{\ell,1}, \ldots, s_{\ell,K_\ell}\}$ from the layer above. Each signal is a noisy projection of the state:
\begin{equation}\label{eq:signal_structure}
s_{\ell,k} = H_{\ell,k} \omega + \varepsilon_{\ell,k}, \quad \varepsilon_{\ell,k} \sim \mathcal{N}(0, \sigma_{\ell,k}^2 I),
\end{equation}
where $H_{\ell,k} \in \mathbb{R}^{d_k \times d}$ is a linear observation operator.

The \emph{intrinsic precision} of signal $k$ at layer $\ell$ is:
\begin{equation}\label{eq:information_depreciation}
\tau^0_{\ell,k} = \bar{\tau}^0_k \cdot \eta_{\ell}^{\,\ell},
\end{equation}
where $\bar{\tau}^0_k \equiv 1/\bar{\sigma}_k^2$ is the raw precision of signal $k$ at the source (layer $0$), and $\eta_\ell \in (0,1)$ is the \emph{endogenous information preservation rate} at layer $\ell$, determined by the attention allocation as specified in Assumption~\ref{ass:endogenous_eta} below.
\end{assumption}

\begin{assumption}[Endogenous Information Depreciation]\label{ass:endogenous_eta}
The signal preservation rate at layer $\ell$ is an increasing, concave function of the average attention per signal:
\begin{equation}\label{eq:endogenous_eta}
\eta_\ell = \underline{\eta} + (\bar{\eta} - \underline{\eta}) \cdot \frac{\bar{\alpha}_\ell}{\bar{\alpha}_\ell + a},
\end{equation}
where $\bar{\alpha}_\ell = A_\ell / K_\ell$ is the average attention per signal at layer $\ell$, $\underline{\eta} \in (0,1)$ is the minimal preservation rate (achieved when no attention is allocated), $\bar{\eta} \in (\underline{\eta}, 1)$ is the maximal preservation rate, and $a > 0$ is a saturation parameter.
\end{assumption}

The signal structure in Assumption \ref{ass:state_signals} captures the empirical reality of AI delegation chains. In a retrieval-augmented generation (RAG) pipeline, the signals are retrieved documents; in a multi-agent coding system, they are code snippets and error messages from upstream agents; in a hierarchical reasoning chain, they are intermediate conclusions. The observation operator $H_{\ell,k}$ captures the fact that each signal may reveal only a subset of the state.

The information depreciation equation \eqref{eq:information_depreciation} is a core structural assumption. It posits that every layer through which a signal passes reduces its intrinsic precision by a factor $\eta_\ell < 1$ per layer. The economic rationale for $\eta_\ell < 1$ is threefold:
\begin{enumerate}
\item \textbf{Algorithmic noise.} Each forward pass of a transformer introduces stochasticity through dropout, layer normalization, and floating-point operations.
\item \textbf{Compression loss.} High-dimensional signal representations are projected into lower-dimensional spaces, discarding fine-grained information.
\item \textbf{Cumulated processing error.} Each layer's output is an imperfect summary of its inputs; these imperfections compound through the chain.
\end{enumerate}

\textbf{Micro-foundation for endogenous depreciation.} Assumption~\ref{ass:endogenous_eta} provides the central micro-foundation that endogenizes the preservation rate $\eta_\ell$. The key economic mechanism is that \emph{attention investment improves signal preservation}: allocating more attention capacity $A_\ell$ at layer $\ell$ raises the average attention per signal $\bar{\alpha}_\ell = A_\ell / K_\ell$, which in turn increases the preservation rate $\eta_\ell$. The functional form \eqref{eq:endogenous_eta} is a Michaelis--Menten (or Hill) function with three economically interpretable properties:
\begin{enumerate}[label=(\roman*)]
\item \textbf{Minimal preservation} ($\underline{\eta}$): Even when no attention is allocated ($A_\ell \to 0$), the preservation rate remains bounded below by $\underline{\eta} > 0$. This captures the fact that LLM forward passes have a baseline level of signal retention due to the deterministic components of the architecture (feed-forward transformations, residual connections).
\item \textbf{Maximal preservation bounded below unity} ($\bar{\eta} < 1$): Even with infinite attention ($A_\ell \to \infty$), the preservation rate cannot exceed $\bar{\eta} < 1$. This is the critical property that addresses the reviewer's concern about why $\eta < 1$ is not merely an exogenous assumption. The bound $\bar{\eta} < 1$ reflects \emph{inherent algorithmic noise} in LLM processing: even the best-performing models, given unlimited context windows and compute, cannot perfectly preserve signal fidelity due to (a) irreversible compression in the embedding space, (b) non-invertible attention operations, and (c) fundamental limits on lossy summarization. The gap $1 - \bar{\eta}$ is the residual depreciation that no amount of attention can eliminate.
\item \textbf{Diminishing returns to attention} (concavity): The function is concave in $\bar{\alpha}_\ell$, reflecting the intuition that the marginal improvement in signal preservation from additional attention decreases as the budget grows.
\end{enumerate}

The comparative statics of the endogenous preservation rate are immediate from differentiation of \eqref{eq:endogenous_eta}:
\begin{equation}\label{eq:eta_comp_statics}
\frac{\partial \eta_\ell}{\partial A_\ell} = \frac{(\bar{\eta} - \underline{\eta}) \cdot a / K_\ell}{(\bar{\alpha}_\ell + a)^2} > 0, \qquad
\frac{\partial^2 \eta_\ell}{\partial A_\ell^2} = -\frac{2(\bar{\eta} - \underline{\eta}) \cdot a / K_\ell^2}{(\bar{\alpha}_\ell + a)^3} < 0,
\end{equation}
confirming that $\eta_\ell$ is increasing and concave in the attention budget. The saturation parameter $a$ controls the sensitivity of preservation to attention: when $a$ is small, preservation responds strongly to attention increases; when $a$ is large, large attention budgets are required to raise $\eta_\ell$ meaningfully above $\underline{\eta}$.

The assumption that the base preservation rate is uniform across layers and signals (conditional on $\bar{\alpha}_\ell$) is a simplification; we discuss heterogeneous depreciation in Section \ref{sec:extensions}. We operationalize task complexity $t$ as $t = \log(1 + \text{reasoning\_steps}) / \log(1 + \text{max\_steps})$, which is measurable from benchmark metadata and API logs.

\begin{remark}[Scalar-State Simplification]\label{rem:scalar}
In the baseline analysis, we set $d = 1$ (scalar state) and $H_{\ell,k} = 1$ for all $\ell, k$. This simplifies the precision-aggregation technology to scalar addition: $\tau_\ell = \tau_{\ell}^{\text{prior}} + \sum_k g_\ell(\alpha_{\ell,k}) \tau^0_{\ell,k}$. The scalar assumption is restrictive when signals reveal different dimensions of a multi-faceted state. In Section \ref{sec:extensions}, we sketch how the model extends to $d > 1$ with a precision matrix; the key insight---that attention scarcity bounds depth through compounding precision loss---is preserved.
\end{remark}

\begin{assumption}[Attention Budget]\label{ass:attention_budget}
Each agent at layer $\ell$ has a finite attention budget $A_\ell > 0$. The agent chooses attention weights $\boldsymbol{\alpha}_\ell = (\alpha_{\ell,1}, \ldots, \alpha_{\ell,K_\ell})$ satisfying
\begin{equation}\label{eq:budget_constraint}
\sum_{k=1}^{K_\ell} \alpha_{\ell,k} \leq A_\ell, \quad \alpha_{\ell,k} \geq 0 \; \forall k.
\end{equation}
The weight $\alpha_{\ell,k}$ measures the fraction of the agent's attention capacity devoted to signal $k$. If $\alpha_{\ell,k} = 0$, the signal is completely ignored.
\end{assumption}

The attention budget $A_\ell$ is the central scarce resource in our model. In transformer-based LLMs, it has three direct interpretations: (i) the context-window length; (ii) the number of attention heads; and (iii) the effective working-memory capacity. The budget constraint \eqref{eq:budget_constraint} is deliberately stylized: it captures the fundamental scarcity of processing capacity in a tractable linear form, even though actual transformer attention involves non-linear softmax interactions and multi-head parallelism. We view the linear budget as a first-order approximation that preserves the key economic insight---scarcity forces selectivity---while admitting closed-form solutions.

It is important to emphasize that $A_\ell$ is not a choice variable of the agent at layer $\ell$; it is a technological parameter determined by the model architecture. The principal chooses $A_\ell$ when selecting which model to deploy at each layer. This creates a two-stage organizational design problem: the principal chooses the architecture $\{A_\ell\}_{\ell=1}^L$ and the chain depth $L$; the agents then solve their attention-allocation problems given these budgets.

\begin{assumption}[Attention Technology]\label{ass:attention_technology}
The effective precision achieved by the agent at layer $\ell$ is
\begin{equation}\label{eq:effective_precision}
\tau_\ell(\boldsymbol{\alpha}_\ell) = \sum_{k=1}^{K_\ell} g_\ell(\alpha_{\ell,k}) \tau^0_{\ell,k},
\end{equation}
where $g_\ell : \mathbb{R}_+ \to \mathbb{R}_+$ is a twice-differentiable, increasing, and concave function satisfying $g_\ell(0) = 0$, $g'_\ell(0) > 0$, and $\lim_{\alpha \to \infty} g'_\ell(\alpha) = 0$.
\end{assumption}

The function $g_\ell$ captures the attention technology: how does allocating attention weight $\alpha$ to a signal translate into effective processing precision? The concavity of $g_\ell$ reflects diminishing returns to attention. Our baseline specification uses a power function:
\begin{equation}\label{eq:power_function}
g_\ell(\alpha) = \alpha^\gamma, \quad \gamma \in (0,1).
\end{equation}
The parameter $\gamma$ is the attention elasticity of precision. When $\gamma$ is close to 1, attention is nearly linear in precision. When $\gamma$ is small, attention is highly concave. We calibrate $\gamma$ from empirical data on LLM accuracy as a function of context length in Section \ref{sec:calibration}.

\begin{remark}[Robustness of the $g(\cdot)$ Specification]\label{rem:g_robustness}
The qualitative properties of our model---the existence of an interior optimal depth, the comparative statics with respect to the preservation parameters $(\underline{\eta}, \bar{\eta})$, the attention budget $A_\ell$, and the saturation parameter $a$, and the efficiency rationales for chaining---do not depend on the power-function specification in equation \eqref{eq:power_function}. Any $g_\ell$ satisfying the regularity conditions in Assumption \ref{ass:attention_technology} yields the same first-order insights because the key economic force is the concavity of $g_\ell$ combined with the linearity of the budget constraint. The power function is chosen for analytical tractability; alternative specifications (e.g., logarithmic, $g(\alpha) = \log(1+\alpha)$; or bounded, $g(\alpha) = \alpha/(1+\alpha)$) yield similar qualitative results but less transparent closed forms. We explore the sensitivity of our results to the functional form of $g$ in Appendix \ref{app:robustness}.
\end{remark}

\begin{assumption}[Cost of Attention]\label{ass:cost}
The cost of maintaining attention capacity $A_\ell$ at layer $\ell$ is $c_\ell(A_\ell)$, where $c_\ell : \mathbb{R}_+ \to \mathbb{R}_+$ is twice-differentiable, increasing, and convex, with $c_\ell(0) = 0$ and $c'_\ell(0) > 0$.
\end{assumption}

For our baseline analysis, we use a quadratic cost:
\begin{equation}\label{eq:quadratic_cost}
c_\ell(A) = \kappa_\ell A + \frac{\phi_\ell}{2} A^2,
\end{equation}
where $\kappa_\ell > 0$ is the marginal cost of the first unit of attention and $\phi_\ell \geq 0$ captures the convexity of compute costs at scale.

\subsection{The Agent's Problem: Attention Allocation}
\label{subsec:agent_problem}

Given the environment above, the agent at layer $\ell$ solves the following optimization problem:
\begin{equation}\label{eq:agent_problem}
\max_{\boldsymbol{\alpha}_\ell \;:\; \eqref{eq:budget_constraint}} V_\ell\bigl(\tau_\ell(\boldsymbol{\alpha}_\ell)\bigr) - c_\ell(A_\ell),
\end{equation}
where $V_\ell(\tau)$ is the value of posterior precision $\tau$ for decision-making at layer $\ell$.

The value function $V_\ell(\tau)$ depends on the decision problem facing the agent. In our baseline, we assume quadratic loss (normal--normal conjugate model), which yields:
\begin{equation}\label{eq:value_function}
V_\ell(\tau) = -\psi_\ell \, \text{Var}(\omega \mid \tau) = -\frac{\psi_\ell}{\tau},
\end{equation}
where $\psi_\ell > 0$ is the marginal value of accuracy. With normal signals and a normal prior, the posterior variance is $\text{Var}(\omega \mid \tau) = 1/\tau$, so $V_\ell(\tau) = -\psi_\ell/\tau$. Since $\tau$ enters negatively, maximizing $V_\ell$ is equivalent to maximizing $\tau$.

More generally, $V_\ell(\tau)$ can incorporate non-quadratic payoffs, discrete actions, or strategic interactions with downstream agents. We explore these extensions in Section \ref{sec:extensions}. For the single-layer analysis in Section \ref{sec:single}, the precise functional form of $V_\ell$ matters only through its monotonicity and curvature.

\subsection{Key Distinctions: Information Depreciation versus Moral Hazard}
\label{subsec:distinctions}

Before proceeding to the analysis, we pause to articulate three core distinctions between our information-depreciation framework and the Garicano--Holmstr\"{o}m--Milgrom (GHM) framework.

\paragraph{Distinction 1: Information depreciation, not moral hazard, is the source of distortion.} In the HM model, agents choose costly effort $e$ to improve output quality. The distortion is moral hazard: agents prefer low effort, and incentives must be designed to induce the efficient level. In our model, distortion arises from the compounding depreciation of signal quality as information flows through the chain---at rate $(1-\eta_\ell)$ per layer---and from the binding attention constraint that forces agents to selectively process only a subset of signals. The agent is not shirking; she faces a technological constraint---the preservation bound $\bar{\eta} < 1$---that no amount of attention can fully overcome. The organizational design problem is to set the chain depth, the attention budgets, and the signal routing so that the marginal benefit of deeper processing equals the marginal cost of depreciation plus attention.

\paragraph{Distinction 2: Selective processing of depreciated signals, not shirking, generates inefficiency.} In the HM model, inefficiency arises because the agent's effort is unobservable and the agent under-supplies it relative to the first-best. In our model, inefficiency arises because (i) each layer depreciates signal quality by factor $(1-\eta_\ell)$ and (ii) the agent, facing a finite budget $A_\ell$, rationally allocates attention only to the signals with highest effective precision. The ignored signals are not random; they are the ones with low post-depreciation precision, low importance for the decision, or high opportunity cost of attention. This generates testable predictions: in an LLM pipeline, we should observe that models drop low-quality retrieved documents before high-quality ones, that the dropout rate increases with chain depth (because $\eta_\ell^{\,\ell}$ shrinks the effective precision of upstream signals), and that the quality of the final output depends more on the depth--attention trade-off than on any notion of ``effort.''

\paragraph{Distinction 3: Organization design as information architecture, not incentive design.} In the HM model, the principal designs incentives (wages, bonuses, monitoring) to align the agent's effort with organizational goals. In our model, the principal designs the \emph{information architecture}: which signals flow to which layers, how much context each layer receives (the attention budgets $A_\ell$), how many layers the chain contains (the depth $L$), and how the outputs of one layer feed into the next. The design variables are $\{A_\ell\}$ (attention budgets), $\{K_\ell\}$ (signal counts), $L$ (chain depth), and the routing structure---not wages or performance contracts. Our model complements the GHM framework by adding information-processing microfoundations: where GHM asks ``how should tasks and incentives be structured to motivate agents?'' we ask ``how should attention and depth be structured to preserve information quality?'' The two questions are orthogonal and jointly relevant for the design of AI delegation systems.

\subsection{Summary of Notation}
\label{subsec:notation}

For reference, Table \ref{tab:notation} summarizes the key symbols used in the model.

\begin{table}[htbp]
\centering
\caption{Notation Summary}
\label{tab:notation}
\begin{tabular}{@{}>{\raggedright\arraybackslash}p{2.5cm}>{\raggedright\arraybackslash}p{10cm}@{}}
\toprule
\textbf{Symbol} & \textbf{Definition} \\
\midrule
$\ell$ & Layer index, $\ell = 0, 1, \ldots, L$ \\
$\omega$ & Underlying state (correct answer / true state) \\
$\eta_\ell$ & Endogenous information preservation rate at layer $\ell$, $\eta_\ell \in (\underline{\eta}, \bar{\eta})$ \\
$\underline{\eta}$ & Minimal preservation rate (no attention), $\underline{\eta} \in (0,1)$ \\
$\bar{\eta}$ & Maximal preservation rate (unlimited attention), $\bar{\eta} \in (\underline{\eta}, 1)$ \\
$a$ & Saturation parameter in endogenous depreciation, $a > 0$ \\
$s_{\ell,k}$ & Signal $k$ at layer $\ell$ \\
$K_\ell$ & Number of signals at layer $\ell$ \\
$\bar{\tau}^0_k$ & Raw intrinsic precision of signal $k$ at source (layer $0$) \\
$\tau^0_{\ell,k}$ & Intrinsic precision of signal $(\ell, k)$ after $\ell$ layers: $\tau^0_{\ell,k} = \bar{\tau}^0_k \cdot \eta_{\ell}^{\,\ell}$ \\
$A_\ell$ & Attention budget at layer $\ell$ \\
$\alpha_{\ell,k}$ & Attention weight on signal $(\ell, k)$ \\
$g_\ell(\cdot)$ & Attention technology (precision transformation function) \\
$\tau_\ell(\boldsymbol{\alpha}_\ell)$ & Effective precision at layer $\ell$ \\
$\rho_\ell$ & Endogenous precision ratio (relative to the precision that would obtain under full attention, as determined by the assumed functional form $g_\ell$) \\
$V_\ell(\tau)$ & Value of precision $\tau$ for decision-making \\
$c_\ell(A_\ell)$ & Cost of attention capacity $A_\ell$ \\
$\gamma$ & Attention elasticity of precision (power-function parameter) \\
$\kappa_\ell, \phi_\ell$ & Cost parameters (linear and quadratic) \\
\bottomrule
\end{tabular}
\end{table}

\section{Single-Layer Analysis}
\label{sec:single}

In this section, we solve the attention-allocation problem for a single layer, derive the optimal attention weights, characterize the resulting distortion, and establish comparative statics with respect to the attention budget, signal importance, and signal heterogeneity. We then connect our results to the rational-inattention literature and derive two propositions that form the building blocks for the multi-layer analysis in Section \ref{sec:multi}.

The single-layer problem is not merely a stepping stone. It is the core economic mechanism of the paper: everything that happens in a delegation chain---distortion, error propagation, depth optimization---traces back to how each individual layer allocates its scarce attention.

\subsection{The Optimization Problem}
\label{subsec:single_opt}

Consider layer $\ell$ in isolation. The agent faces $K_\ell$ signals with intrinsic precisions $\{\tau^0_{\ell,k}\}_{k=1}^{K_\ell}$ and chooses attention weights $\boldsymbol{\alpha}_\ell = (\alpha_{\ell,1}, \ldots, \alpha_{\ell,K_\ell})$ subject to the budget constraint \eqref{eq:budget_constraint}. The effective precision is given by \eqref{eq:effective_precision}, and the agent's payoff is $V_\ell(\tau_\ell(\boldsymbol{\alpha}_\ell)) - c_\ell(A_\ell)$. Since $c_\ell(A_\ell)$ does not depend on $\boldsymbol{\alpha}_\ell$, the optimization problem is:
\begin{equation}\label{eq:single_opt}
\max_{\boldsymbol{\alpha}_\ell \;:\; \eqref{eq:budget_constraint}} V_\ell\Biggl(\sum_{k=1}^{K_\ell} g_\ell(\alpha_{\ell,k}) \tau^0_{\ell,k}\Biggr).
\end{equation}
We assume throughout that $V_\ell(\tau)$ is increasing and concave in $\tau$. This is satisfied by the quadratic-loss specification \eqref{eq:value_function} and by most standard decision problems.

\subsection{First-Order Conditions and Optimal Attention}
\label{subsec:foc}

Let $\lambda_\ell \geq 0$ be the Lagrange multiplier on the budget constraint \eqref{eq:budget_constraint}. The Lagrangian is:
\begin{equation}\label{eq:lagrangian}
\mathcal{L}(\boldsymbol{\alpha}_\ell, \lambda_\ell) = V_\ell\Biggl(\sum_{k=1}^{K_\ell} g_\ell(\alpha_{\ell,k}) \tau^0_{\ell,k}\Biggr) - \lambda_\ell\Biggl(\sum_{k=1}^{K_\ell} \alpha_{\ell,k} - A_\ell\Biggr) + \sum_{k=1}^{K_\ell} \mu_{\ell,k} \alpha_{\ell,k},
\end{equation}
where $\mu_{\ell,k} \geq 0$ are multipliers on the non-negativity constraints.

The first-order condition for $\alpha_{\ell,k}$ is:
\begin{equation}\label{eq:foc}
V'_\ell\bigl(\tau_\ell(\boldsymbol{\alpha}_\ell)\bigr) \cdot g'_\ell(\alpha_{\ell,k}) \tau^0_{\ell,k} = \lambda_\ell - \mu_{\ell,k}.
\end{equation}
This condition has a straightforward economic interpretation. The left-hand side is the marginal benefit of allocating attention to signal $k$: the marginal value of precision ($V'_\ell$) times the marginal precision gain from attention ($g'_\ell(\alpha_{\ell,k}) \tau^0_{\ell,k}$). The right-hand side is the marginal cost: the shadow price of the attention budget ($\lambda_\ell$) minus the relaxation of the non-negativity constraint ($\mu_{\ell,k}$). At an interior solution ($\alpha_{\ell,k} > 0$), $\mu_{\ell,k} = 0$ and the marginal benefit equals the shadow price $\lambda_\ell$.

\begin{proposition}[Optimal Attention Allocation]\label{prop:optimal_attention}
Suppose $V_\ell$ is increasing and concave, $g_\ell$ is increasing and concave with $g'_\ell(0) > 0$, and $\tau^0_{\ell,k} > 0$ for all $k$. Then:
\begin{enumerate}[label=(\roman*)]
\item The optimal attention allocation $\boldsymbol{\alpha}^*_\ell(A_\ell)$ exists and is unique.
\item If the budget binds ($\sum_k \alpha^*_{\ell,k} = A_\ell$), the optimal weights satisfy
\begin{equation}\label{eq:interior_foc}
g'_\ell(\alpha^*_{\ell,k}) \tau^0_{\ell,k} = g'_\ell(\alpha^*_{\ell,j}) \tau^0_{\ell,j} \quad \forall k,j \text{ with } \alpha^*_{\ell,k}, \alpha^*_{\ell,j} > 0.
\end{equation}
\item The optimal weight on signal $k$ is increasing in its intrinsic precision $\tau^0_{\ell,k}$ and decreasing in the precision of other signals.
\end{enumerate}
\end{proposition}

\begin{proof}
See Appendix \ref{app:proof_prop1}.
\end{proof}

Intuition. The optimality condition \eqref{eq:interior_foc} says that the agent equalizes the marginal precision product across all signals that receive positive attention. Signals with higher intrinsic precision receive more attention not because they are ``preferred'' in any absolute sense, but because they yield higher marginal returns: each unit of attention devoted to a high-quality signal extracts more information than a unit devoted to a low-quality signal. This is the familiar logic of equalizing marginal products across inputs, applied here to the allocation of attention rather than labor or capital.

\subsection{From Attention Budget to Preservation Rate}
\label{subsec:ri_to_eta}

We now provide the formal mapping from the attention budget $A_\ell$ to the endogenous preservation rate $\eta_\ell$, addressing the reviewer's request for an explicit derivation of how rational inattention translates into the preservation bounds in Assumption~\ref{ass:endogenous_eta}.

\begin{proposition}[Rational Inattention to Preservation Rate]\label{prop:ri_to_eta}
Under the normal--normal conjugate model with attention budget $A_\ell$ and $K_\ell$ signals, the mutual information between the state $\omega$ and the attention-weighted signal vector is
\begin{equation}\label{eq:mutual_info}
I\bigl(\omega; \mathbf{s}_\ell \mid \boldsymbol{\alpha}_\ell\bigr) = \frac{1}{2} \log\left(1 + \frac{\tau_\ell(\boldsymbol{\alpha}_\ell)}{\tau^{\text{prior}}_\ell}\right),
\end{equation}
and the preservation rate $\eta_\ell$ is calibrated as
\begin{equation}\label{eq:eta_calibration}
\eta_\ell = \underline{\eta} + (\bar{\eta} - \underline{\eta}) \cdot \frac{I(\omega; \mathbf{s}_\ell)}{I_{\max}},
\end{equation}
where $I_{\max} = \frac{1}{2}\log(1 + \tau^{\text{FB}}_\ell / \tau^{\text{prior}}_\ell)$ is the first-best mutual information. Under the power-function specification \eqref{eq:power_function}, this yields Assumption~\ref{ass:endogenous_eta} with saturation parameter $a = (\tau^{\text{prior}}_\ell)^{1/\gamma} / \tau^{\text{prior}}_\ell$.
\end{proposition}

\begin{proof}
\textbf{Step 1: Mutual information.} With normal prior $\omega \sim \mathcal{N}(\mu, (\tau^{\text{prior}}_\ell)^{-1})$ and normal signals $s_{\ell,k} = \omega + \varepsilon_{\ell,k}$ where $\varepsilon_{\ell,k} \sim \mathcal{N}(0, (\alpha^\gamma_{\ell,k} \tau^0_{\ell,k})^{-1})$, the posterior precision is $\tau_\ell = \tau^{\text{prior}}_\ell + \sum_k \alpha^\gamma_{\ell,k} \tau^0_{\ell,k}$. The mutual information is the reduction in entropy:
\begin{equation*}
I(\omega; \mathbf{s}_\ell) = h(\omega) - h(\omega \mid \mathbf{s}_\ell) = \frac{1}{2}\log\left(\frac{\tau_\ell}{\tau^{\text{prior}}_\ell}\right) = \frac{1}{2}\log\left(1 + \frac{\tau_\ell - \tau^{\text{prior}}_\ell}{\tau^{\text{prior}}_\ell}\right),
\end{equation*}
which is \eqref{eq:mutual_info}.

\textbf{Step 2: From mutual information to preservation.} The mutual information is bounded above by the first-best: $I \leq I_{\max} = \frac{1}{2}\log(1 + \tau^{\text{FB}}_\ell / \tau^{\text{prior}}_\ell)$. The ratio $I / I_{\max} \in [0,1]$ measures how close the attention allocation is to full information extraction. Mapping this ratio to the preservation rate through a monotone increasing function yields \eqref{eq:eta_calibration}.

\textbf{Step 3: Calibration to Assumption~\ref{ass:endogenous_eta}.} Under the power-function technology, the effective precision is $\tau_\ell = \tau^{\text{prior}}_\ell + A^\gamma_\ell \cdot \bar{\tau}$ for some $\bar{\tau} > 0$ (when signals are symmetric). The mutual information becomes
\begin{equation*}
I(A_\ell) = \frac{1}{2}\log\left(1 + \frac{A^\gamma_\ell \bar{\tau}}{\tau^{\text{prior}}_\ell}\right).
\end{equation*}
For small $A_\ell$, $I(A_\ell) \approx \frac{1}{2} A^\gamma_\ell \bar{\tau} / \tau^{\text{prior}}_\ell$, which is concave. A Michaelis--Menten approximation $I(A_\ell) / I_{\max} \approx A_\ell / (A_\ell + a)$ with $a = (\tau^{\text{prior}}_\ell)^{1/\gamma} / \bar{\tau}^{1/\gamma}$ yields equation \eqref{eq:endogenous_eta}. \hfill $\square$
\end{proof}

\begin{remark}[Scope of the Calibration]
Proposition~\ref{prop:ri_to_eta} provides a micro-foundation for Assumption~\ref{ass:endogenous_eta} under the normal--normal conjugate model. The key insight is that the preservation rate $\eta_\ell$ is not an arbitrary function of $A_\ell$ but a monotone transformation of the normalized mutual information. The functional form (Michaelis--Menten) is an analytically convenient approximation to the log-form mutual information; the qualitative conclusions do not depend on this approximation.
\end{remark}

For the power-function specification $g_\ell(\alpha) = \alpha^\gamma$, the first-order condition becomes:
\begin{equation}\label{eq:closed_form}
\alpha^*_{\ell,k} = \biggl(\frac{\gamma \tau^0_{\ell,k}}{\lambda_\ell}\biggr)^{\frac{1}{1-\gamma}} \quad \text{(interior solutions)}.
\end{equation}
This closed form reveals the key comparative static: the optimal attention weight is log-convex in intrinsic precision when $\gamma < 1/2$ and log-concave when $\gamma > 1/2$. In the empirically relevant case of $\gamma \approx 0.3$ (calibrated from LLM accuracy data), small differences in signal quality are amplified into large differences in attention allocation: a signal with twice the precision receives approximately $2^{1/(1-0.3)} \approx 2.69$ times the attention. This amplification is the mechanism by which ``important'' signals dominate and ``unimportant'' signals are starved.

\subsection{Distortion and the Attention Gap}
\label{subsec:distortion}

We now define the distortion arising from attention allocation and characterize how it depends on the attention budget and signal structure.

\begin{definition}[Attentional Distortion]\label{def:distortion}
The distortion on signal $k$ at layer $\ell$ is
\begin{equation}\label{eq:distortion}
\delta_{\ell,k} = (\bar{\alpha}_{\ell,k} - \alpha^*_{\ell,k})^+,
\end{equation}
where $\bar{\alpha}_{\ell,k}$ is the attention weight required to fully process signal $k$ (i.e., achieve $\tau^0_{\ell,k}$ without loss), and $(\cdot)^+ = \max\{\cdot, 0\}$.
\end{definition}

In the baseline where $\bar{\alpha}_{\ell,k} = 1$ for all $k$ (unit attention achieves full precision), the distortion is simply $\delta_{\ell,k} = (1 - \alpha^*_{\ell,k})^+$. If $\alpha^*_{\ell,k} \geq 1$, the signal is fully processed and $\delta_{\ell,k} = 0$. If $\alpha^*_{\ell,k} < 1$, the signal is under-processed and the distortion is positive.

We can also define an aggregate distortion measure:
\begin{definition}[Aggregate Distortion]\label{def:aggregate_distortion}
The aggregate attentional distortion at layer $\ell$ is
\begin{equation}\label{eq:aggregate_distortion}
\Delta_\ell = \frac{\sum_{k=1}^{K_\ell} \omega_{\ell,k} \delta_{\ell,k}}{\sum_{k=1}^{K_\ell} \omega_{\ell,k}},
\end{equation}
where $\omega_{\ell,k} > 0$ is the importance weight of signal $k$ for the layer's decision. In the quadratic-loss case, $\omega_{\ell,k} = \tau^0_{\ell,k}$.
\end{definition}

\begin{proposition}[Distortion Comparative Statics]\label{prop:distortion_comp_statics}
The aggregate distortion $\Delta_\ell$ satisfies:
\begin{enumerate}[label=(\roman*)]
\item $\Delta_\ell$ is decreasing in the attention budget $A_\ell$: $\partial \Delta_\ell / \partial A_\ell < 0$.
\item $\Delta_\ell$ is increasing in the number of signals $K_\ell$: $\partial \Delta_\ell / \partial K_\ell > 0$ when $A_\ell$ is held fixed.
\item $\Delta_\ell$ is decreasing in the heterogeneity of signal precisions: holding mean precision fixed, $\Delta_\ell$ is lower when $\{\tau^0_{\ell,k}\}$ is more dispersed.
\end{enumerate}
\end{proposition}

\begin{proof}
See Appendix \ref{app:proof_prop2}.
\end{proof}

Intuition. Proposition \ref{prop:distortion_comp_statics} delivers three testable predictions. First, increasing the attention budget reduces distortion: if the principal upgrades from GPT-3.5 (4K context) to GPT-4 (128K context), the model can attend to more signals and distortion falls. Second, increasing the number of signals while holding the budget fixed increases distortion: this is the mechanism behind ``attention overload'' in RAG systems. Third, signal heterogeneity reduces aggregate distortion because the optimal allocation concentrates attention on the best signals. If all signals are of equal quality, the agent spreads attention evenly and all signals are under-processed. If signals vary in quality, the agent concentrates on the high-quality ones, achieving near-full precision on the most important signals while deliberately sacrificing the low-quality ones.

\subsection{The Endogenous Distortion Factor}
\label{subsec:endogenous_distortion}

We now derive the endogenous distortion measure that connects the single-layer attention-allocation problem to the multi-layer transmission analysis.

\begin{definition}[Single-Layer Distortion Measure]\label{def:single_layer_distortion}
The single-layer distortion measure at layer $\ell$ is
\begin{equation}\label{eq:single_distortion}
\tilde{\delta}_\ell = 1 - \Delta_\ell = 1 - \frac{\sum_k \omega_{\ell,k} (\bar{\alpha}_{\ell,k} - \alpha^*_{\ell,k})^+}{\sum_k \omega_{\ell,k}}.
\end{equation}
\end{definition}

The measure $\tilde{\delta}_\ell \in [0,1]$ captures the fraction of signal importance that survives the attention-allocation process at a single layer, ignoring upstream effects. When $\tilde{\delta}_\ell = 1$, all signals are fully processed and no distortion occurs. When $\tilde{\delta}_\ell = 0$, all signals are ignored and the layer transmits pure noise. In general, $\tilde{\delta}_\ell$ is a function of the attention budget, the number of signals, their precisions, and the attention technology:
\begin{equation}\label{eq:delta_function}
\tilde{\delta}_\ell = \tilde{\delta}(A_\ell, K_\ell, \{\tau^0_{\ell,k}\}, g_\ell).
\end{equation}
This is the core endogenization result. The distortion measure is not a primitive parameter but a derived object that depends on observable technological and organizational variables.

\begin{remark}[From $\tilde{\delta}_\ell$ to $\rho_\ell$]\label{rem:delta_to_rho}
The single-layer measure $\tilde{\delta}_\ell$ is useful for analyzing distortion at an isolated layer. In the multi-layer analysis (Section \ref{sec:multi}), we introduce the endogenous transmission factor $\rho_\ell \equiv \tau^*_\ell / \tau^{\text{FB}}_\ell$, which accounts for the compounding of inherited precision. For the bottom layer with $\tau^{\text{prior}}_0 = 0$, the two coincide: $\rho_0 = \tilde{\delta}_0 = 1 - \Delta_0$. For upper layers, $\rho_\ell$ generalizes $\tilde{\delta}_\ell$ by incorporating the prior precision.
\end{remark}

\subsection{Connection to Rational Inattention}
\label{subsec:ri_connection}

Our model is a direct extension of the rational-inattention framework to hierarchical organizations. We now make this connection explicit.

\textbf{Sims (2003): The single-agent benchmark.} In Sims's model, a single agent faces a multidimensional state $\omega \in \mathbb{R}^d$ and chooses an information structure subject to a mutual-information constraint $I(\omega; \hat{\omega}) \leq \kappa$. Our model specializes Sims in two directions and generalizes it in one. We specialize by assuming a linear precision-aggregation technology (Equation \eqref{eq:effective_precision}) rather than the general information-structure choice. We specialize by assuming a hard attention budget rather than an entropy-based information-cost function. We generalize by embedding the single-agent problem in a hierarchy where the output of one agent becomes the input to the next.

The hard budget vs.~entropy-cost distinction is important. Sims's entropy cost is analytically elegant, but it is less directly mapped to AI technology. The mutual information $I(\omega; \hat{\omega})$ is difficult to measure in an LLM pipeline, whereas the context-window length $A_\ell$ is directly observable. Our hard-budget formulation trades some analytical elegance for empirical tractability.

\textbf{Dessein and Santos (2006): Organizational attention.} Dessein and Santos model an organization in which agents have a fixed attention capacity that can be used for both sending and receiving information. Our model extends theirs in three ways. First, we allow attention to be continuously allocated across signals rather than discretely assigned to tasks. Second, we allow the attention budget to vary across layers ($A_\ell$ need not equal $A_{\ell'}$), capturing the fact that different LLMs in a chain have different context-window sizes. Third, we model the endogenous degradation of information through the chain.

\textbf{Dessein, Galeotti, and Santos (2016): Division of labor.} In their follow-up paper, Dessein et al.~study how attention constraints generate specialization. Our Proposition \ref{prop:distortion_comp_statics}(iii) echoes this result: heterogeneity in signal quality leads to specialization in attention allocation. We extend their analysis by deriving the optimal chain depth from attention costs.

\paragraph{Limitations of the microfoundation.}
Our model makes two deliberate simplifications relative to actual transformer architectures. First, the linear budget constraint \eqref{eq:budget_constraint} is a stylized representation of the context window limit; real transformers use soft attention weights rather than hard capacity constraints. Second, we abstract from the token-token interaction structure of self-attention, treating each signal as an independent information source. These simplifications trade descriptive realism for analytical tractability and should be interpreted as a first-order approximation rather than a literal mapping.

\subsection{Alternative: Entropy Cost Formulation}
\label{sec:entropy_cost}

\subsubsection{Entropy Cost Setup}

In the rational inattention framework of \citet{sims2003implications}, the agent's information processing cost is proportional to the reduction in uncertainty (measured by mutual information) achieved by observing a signal. Specifically, if the agent chooses an information structure $P(\tilde{s} \mid \omega)$ that maps the true state $\omega \sim \mathcal{N}(\mu, \tau^{-1})$ into a signal realization $\tilde{s}$, the associated cost is:
\begin{equation}
C(I) = \theta \cdot I(\omega; \tilde{s}),
\end{equation}
where $\theta > 0$ is the marginal price of information and $I(\omega; \tilde{s})$ denotes the mutual information:
\begin{equation}
I(\omega; \tilde{s}) = h(\omega) - h(\omega \mid \tilde{s}) = \frac{1}{2}\log_2\left(\frac{\tau_{\text{post}}}{\tau}\right),
\end{equation}
where $\tau_{\text{post}} = \tau + \tau_s$ is the posterior precision after observing signal $\tilde{s} = \omega + \varepsilon$ with noise precision $\tau_s$, and $h(\cdot)$ denotes differential entropy.

When a chain of $L$ agents each face entropy costs, agent $\ell$ chooses a signal structure $P(\tilde{s}_\ell \mid \omega)$ at cost $\theta_\ell \cdot I(\omega; \tilde{s}_\ell)$, and then passes a processed message $m_\ell$ to the next layer. The key difference from the hard-budget formulation is that \emph{all signals are processed, but with endogenously varying precision}, rather than a subset being selected with perfect precision.

\subsubsection{Entropy-Cost Version of Our Model}

Consider our multi-layer information processing chain under entropy costs. Each layer $\ell$ chooses the precision $\tau_\ell$ of its signal subject to an entropy cost $\theta_\ell \cdot I(\omega; \tilde{s}_\ell)$. The total cost for the chain is:
\begin{equation}
C_{\text{total}} = \sum_{\ell=1}^{L} \theta_\ell \cdot I(\omega; \tilde{s}_\ell) = \sum_{\ell=1}^{L} \frac{\theta_\ell}{2} \log\left(1 + \frac{\tau_\ell}{\tau_{\ell-1}^{*}}\right),
\end{equation}
where $\tau_{\ell-1}^{*}$ is the effective precision inherited from the previous layer. The optimization problem at each layer is to choose $\tau_\ell$ to maximize:
\begin{equation}
U_\ell(\tau_\ell) = V(\tau_\ell^{*}) - \theta_\ell \cdot I(\omega; \tilde{s}_\ell),
\end{equation}
where $\tau_\ell^{*} = \tau_{\ell-1}^{*} + \tau_\ell$ is the posterior precision.

\paragraph{Key Result: Robustness of Finite Depth.} We now establish that the finite-optimal-depth result carries over to the entropy-cost setting:

\begin{proposition}[Finite Optimal Depth under Entropy Costs]\label{prop:finite_depth_entropy}
Suppose each layer $\ell$ faces an entropy cost with parameter $\theta_\ell > 0$, and the value function $V(\tau)$ satisfies $V'(\tau) > 0$, $V''(\tau) < 0$, and the Inada-like condition $\lim_{\tau \to \infty} V'(\tau) = 0$. Then:
\begin{enumerate}
\item[(i)] \textbf{Concave precision response:} The optimal precision $\tau_\ell^{*}(\theta_\ell)$ chosen at each layer satisfies $\frac{\partial^2 \tau_\ell^{*}}{\partial \theta_\ell^2} > 0$, i.e., precision decreases convexly in the cost parameter.
\item[(ii)] \textbf{Diminishing information increment:} The marginal contribution of layer $\ell$ to total precision, $\Delta_\ell \equiv \tau_\ell^{*} - \tau_{\ell-1}^{*}$, is decreasing in $\ell$ when $\theta_\ell$ is non-decreasing.
\item[(iii)] \textbf{Finite optimal depth:} If $\inf_\ell \theta_\ell > 0$, there exists a finite $L^{*} < \infty$ such that adding layer $L^{*}+1$ yields negative net value.
\end{enumerate}
\end{proposition}

\begin{proof}[Proof Sketch]
(i) At layer $\ell$, the first-order condition for optimal precision is:
\begin{equation}
V'(\tau_{\ell-1}^{*} + \tau_\ell) = \frac{\theta_\ell}{2} \cdot \frac{1}{\tau_{\ell-1}^{*} + \tau_\ell}.
\end{equation}
This implicitly defines $\tau_\ell^{*}$ as a function of $\theta_\ell$. By the implicit function theorem and the convexity of the entropy cost in precision, the response function is concave in $1/\theta_\ell$.

(ii) Since $\tau_\ell^{*}$ is increasing in inherited precision $\tau_{\ell-1}^{*}$ but with slope less than 1 (due to $V'' < 0$), the increment $\Delta_\ell$ decreases along the chain when cost parameters are bounded away from zero.

(iii) The net value of adding layer $L+1$ is $V(\tau_{L+1}^{*}) - V(\tau_{L}^{*}) - \theta_{L+1} \cdot I(\omega; \tilde{s}_{L+1})$. Since $V'(\tau) \to 0$ as $\tau \to \infty$ while the marginal entropy cost remains bounded below by $\underline{\theta} \cdot (2\tau)^{-1} > 0$, there exists a finite threshold beyond which the cost exceeds the benefit.
\end{proof}

\paragraph{Intuition for Robustness.} The critical insight is that the entropy cost function's \emph{convexity} in precision (arising from the logarithmic form of mutual information) naturally generates the same diminishing-returns property that our hard-budget model captures through the fixed-capacity constraint.

\begin{table}[htbp]
\centering
\footnotesize
\setlength{\tabcolsep}{3pt}
\caption{Comparison: Hard-Budget vs.~Entropy-Cost Assumptions}
\label{tab:assumption_comparison}
\begin{tabular}{@{}>{\raggedright\arraybackslash}p{2.8cm}>{\raggedright\arraybackslash}p{5.6cm}>{\raggedright\arraybackslash}p{5.6cm}@{}}
\toprule
\textbf{Dimension} & \textbf{Hard-Budget} & \textbf{Entropy-Cost} \\
\midrule
Information loss & Signals beyond budget $\bar{A}_\ell$ ignored or truncated. & All signals processed with endogenously varying precision. \\
Optimization variable & Attention weight vector $\boldsymbol{\alpha}_\ell \in \Delta^{K}$ (discrete capacity). & Information structure $P(\tilde{s} \mid \omega)$ (continuous precision). \\
Observability & $\bar{A}_\ell$ directly observable (context window size). & $\theta_\ell$ latent; requires calibration. \\
LLM mapping & Direct: context window imposes hard token bound. & Indirect: entropy reduction to FLOPs mapping. \\
Core conclusion & Finite optimal depth (Prop.~\ref{prop:finite_depth}). & Finite optimal depth (Prop.~\ref{prop:finite_depth_entropy}). \\
Parametrization & Sharp cutoff at $\bar{A}_\ell$. & Smooth; no cutoff. \\
Applicable scenarios & Fixed context-window LLMs. & Variable-compute systems. \\
\bottomrule
\end{tabular}
\end{table}

\subsubsection{Discussion: Why Hard Budgets in the Baseline Model}

We adopt the hard-budget assumption as our baseline for three substantive reasons related to \emph{empirical operability} and the specific application to large language models:

\paragraph{Direct observability.} The hard budget $\bar{A}_\ell$ maps directly and transparently to a measurable technological parameter: the context window length of an LLM (e.g., 4,096 tokens for GPT-3, 128,000 tokens for GPT-4 Turbo). This observable anchor allows us to discipline the model with empirical moments and conduct counterfactuals about window-size expansion.

\paragraph{Descriptive accuracy for current LLM architectures.} Current production LLMs operate with fixed context windows: tokens beyond the window length are either truncated entirely or compressed through external summarization mechanisms. This is precisely the hard-budget mechanism.

\paragraph{Analytical tractability.} The hard-budget formulation yields a finite-dimensional optimization over the simplex $\Delta^{K}$, enabling closed-form solutions for the attention allocation and sharp comparative statics.

\paragraph{Entropy costs as robustness check.} We view the entropy-cost formulation primarily as a \emph{robustness check} on the qualitative conclusions. As Proposition \ref{prop:finite_depth_entropy} demonstrates, the core mechanism---diminishing returns to depth driven by concave precision accumulation---is preserved across both formulations.

\subsection{Propositions for the Multi-Layer Analysis}
\label{subsec:single_multi_building}

We conclude the single-layer analysis with two propositions that will serve as building blocks for the multi-layer analysis.

\begin{proposition}[Precision Response to Budget]\label{prop:precision_response}
Let $\tau^*_\ell(A_\ell) \equiv \tau_\ell(\boldsymbol{\alpha}^*_\ell(A_\ell))$ be the optimal effective precision. Then:
\begin{enumerate}[label=(\roman*)]
\item $\tau^*_\ell(A_\ell)$ is increasing and concave in $A_\ell$.
\item The marginal return to attention, $\partial \tau^*_\ell / \partial A_\ell$, is decreasing in $K_\ell$.
\item $\tau^*_\ell(A_\ell)$ is bounded above by $\bar{\tau}_\ell \equiv \sum_k \tau^0_{\ell,k}$.
\end{enumerate}
\end{proposition}

\begin{proof}
See Appendix \ref{app:proof_prop3}.
\end{proof}

\begin{proposition}[Cost-Minimizing Budget]\label{prop:cost_min}
For a target precision level $\bar{\tau}$, the cost-minimizing attention budget $A^*_\ell(\bar{\tau})$ satisfies:
\begin{enumerate}[label=(\roman*)]
\item $A^*_\ell(\bar{\tau})$ is increasing and convex in $\bar{\tau}$.
\item The marginal cost of precision, $\partial c_\ell(A^*_\ell) / \partial \bar{\tau}$, is increasing in $K_\ell$.
\item If $\bar{\tau} > \tau^*_\ell(K_\ell)$, then $A^*_\ell(\bar{\tau})$ does not exist: the target is infeasible.
\end{enumerate}
\end{proposition}

\begin{proof}
See Appendix \ref{app:proof_prop4}.
\end{proof}

\begin{proposition}[Attention and Signal Preservation]\label{prop:attention_preservation}
Under Assumption~\ref{ass:endogenous_eta}, the endogenous preservation rate $\eta_\ell$ at layer $\ell$ satisfies:
\begin{enumerate}[label=(\roman*)]
\item $\eta_\ell$ is strictly increasing and concave in the attention budget $A_\ell$:
\begin{equation}\label{eq:eta_single_comp}
\frac{\partial \eta_\ell}{\partial A_\ell} = \frac{(\bar{\eta} - \underline{\eta}) \, a / K_\ell}{(\bar{\alpha}_\ell + a)^2} > 0, \qquad \frac{\partial^2 \eta_\ell}{\partial A_\ell^2} = -\frac{2(\bar{\eta} - \underline{\eta}) \, a / K_\ell^2}{(\bar{\alpha}_\ell + a)^3} < 0.
\end{equation}
\item $\eta_\ell$ is increasing in the saturation parameter $a$ when $A_\ell$ is small ($\bar{\alpha}_\ell < a$) and decreasing when $A_\ell$ is large ($\bar{\alpha}_\ell > a$).
\item The \emph{effective} signal precision at layer $\ell$, $\tau^0_{\ell,k} = \bar{\tau}^0_k \cdot \eta_\ell^{\,\ell}$, responds to attention according to
\begin{equation}\label{eq:tau0_attention_elasticity}
\frac{\partial \tau^0_{\ell,k}}{\partial A_\ell} = \ell \, \bar{\tau}^0_k \, \eta_\ell^{\,\ell-1} \frac{\partial \eta_\ell}{\partial A_\ell} > 0,
\end{equation}
so the marginal return to attention is larger at deeper layers (higher $\ell$).
\end{enumerate}
\end{proposition}

\begin{proof}
\noindent\textbf{Part (i):} Direct differentiation of equation \eqref{eq:endogenous_eta} with respect to $A_\ell$, using $\bar{\alpha}_\ell = A_\ell / K_\ell$, yields the expressions in \eqref{eq:eta_single_comp}. The positivity of the first derivative follows from $\bar{\eta} > \underline{\eta}$, $a > 0$, and $K_\ell > 0$. The negativity of the second derivative follows from the same conditions.

\noindent\textbf{Part (ii):} Differentiating \eqref{eq:endogenous_eta} with respect to $a$:
\begin{equation*}
\frac{\partial \eta_\ell}{\partial a} = (\bar{\eta} - \underline{\eta}) \cdot \frac{-\bar{\alpha}_\ell}{(\bar{\alpha}_\ell + a)^2} < 0 \quad \text{for all } \bar{\alpha}_\ell > 0.
\end{equation*}
Thus, for fixed $\bar{\alpha}_\ell$, a larger saturation parameter $a$ \emph{decreases} the preservation rate, as it makes attention less effective at improving signal fidelity. The cross-partial derivative reveals how the attention sensitivity depends on $a$:
\begin{equation*}
\frac{\partial^2 \eta_\ell}{\partial A_\ell \, \partial a} = -\frac{(\bar{\eta} - \underline{\eta})(\bar{\alpha}_\ell - a)}{K_\ell(\bar{\alpha}_\ell + a)^3},
\end{equation*}
which is positive when $\bar{\alpha}_\ell < a$ (attention and saturation are complements) and negative when $\bar{\alpha}_\ell > a$ (attention and saturation are substitutes). Intuitively, when attention is scarce relative to the saturation threshold ($\bar{\alpha}_\ell < a$), a higher $a$ reduces the \emph{curvature} of the $\eta_\ell$ function, making the marginal return to attention more stable across budget levels.

\noindent\textbf{Part (iii):} From $\tau^0_{\ell,k} = \bar{\tau}^0_k \eta_\ell^{\,\ell}$, differentiate with respect to $A_\ell$:
\begin{equation*}
\frac{\partial \tau^0_{\ell,k}}{\partial A_\ell} = \bar{\tau}^0_k \cdot \ell \eta_\ell^{\,\ell-1} \cdot \frac{\partial \eta_\ell}{\partial A_\ell}.
\end{equation*}
Since $\partial \eta_\ell / \partial A_\ell > 0$, this derivative is positive. The multiplicative factor $\ell$ implies that the marginal return to attention is larger at deeper layers: raising $A_\ell$ by one unit improves signal preservation more when the cumulative depreciation exponent $\ell$ is high. This generates a \emph{double dividend} of attention at deeper layers: not only is there more attention per signal, but each unit of preservation improvement is amplified by the exponent.
\end{proof}

\begin{remark}[The Double Dividend of Attention]\label{rem:double_dividend}
Proposition~\ref{prop:attention_preservation}(iii) reveals a novel economic force: the marginal return to attention is \emph{increasing} in layer depth. At layer $\ell$, increasing $A_\ell$ by $\Delta A$ raises the preservation rate from $\eta_\ell$ to $\eta_\ell + \Delta \eta$, and this improvement is compounded $\ell$ times in the signal precision formula $\tau^0_{\ell,k} \propto \eta_\ell^{\,\ell}$. This creates a strong organizational motive for front-loading attention budgets (Proposition~\ref{prop:front_loading}): early layers have low $\ell$, so the preservation elasticity is small, while later layers---which aggregate already-depreciated signals---benefit disproportionately from attention investment. The double dividend complements the standard ``compounding loss'' intuition: not only does early attention prevent error accumulation, but the \emph{marginal value} of attention itself grows with depth because of the exponent in the depreciation formula.
\end{remark}

\subsection{Discussion: Why This Is Not ``Concave Benefit + Convex Cost''}
\label{subsec:not_concave_convex}

We close this section by addressing a potential objection: is the attention-allocation model merely a relabeled version of the ``concave benefit + convex cost = finite optimum'' structure? We argue that it is not, for three reasons.

\textbf{First, the concavity emerges from the technology, not the objective.} In the standard framework, the principal assumes a concave benefit function and a convex cost function to ensure a finite optimum. In our framework, the concavity of $\tau^*_\ell(A_\ell)$ (Proposition \ref{prop:precision_response}) is \emph{derived} from the attention technology: diminishing returns arise because the agent allocates attention to signals in order of quality, not because we assumed a concave production function.

\textbf{Second, the optimization is over a vector, not a scalar.} The HM model optimizes over a scalar effort $e$. Our model optimizes over a vector $\boldsymbol{\alpha}_\ell$ subject to a budget constraint. The solution is a system of $K_\ell$ equations with non-negativity constraints and a binding budget. The structure of the solution---which signals receive attention, which are ignored, and how the allocation changes with parameters---is economically meaningful and empirically testable.

\textbf{Third, the distortion is endogenous and derived relative to the assumed $g_\ell(\cdot)$ technology.} In the standard framework, the distortion parameter is exogenous and the output function is ad hoc. In our framework, the distortion measure $\tilde{\delta}_\ell$ is derived from the optimal attention allocation (Definition \ref{def:single_layer_distortion}), the output function is the precision-aggregation technology (Equation \eqref{eq:effective_precision}), and both are directly connected to observable AI technology.

\begin{lemma}[Sufficient Conditions for Front-Loading]\label{lem:front_loading_sufficient}
Suppose the value function $V(\tau)$ is increasing and concave, and the optimal precision $\tau^*_\ell(A_\ell)$ is increasing and concave in $A_\ell$ for each $\ell$. Then the objective $V(\tau^*_L)$ is supermodular in $(A_\ell, A_{\ell'})$ for $\ell < \ell'$. Consequently, the optimal budget allocation satisfies $A^*_0 \geq A^*_1 \geq \cdots \geq A^*_L$.
\end{lemma}

\begin{proof}
Consider two adjacent layers $\ell$ and $\ell+1$. The cross-partial derivative of top-layer precision with respect to $(A_\ell, A_{\ell+1})$ is:
\begin{equation*}
\frac{\partial^2 \tau^*_L}{\partial A_\ell \, \partial A_{\ell+1}} = \frac{\partial \rho_\ell}{\partial A_\ell} \cdot \frac{\partial \tau^*_{\ell+1}}{\partial A_{\ell+1}} > 0,
\end{equation*}
since $\partial \rho_\ell / \partial A_\ell > 0$ (Proposition~\ref{prop:distortion_comp_statics}) and $\partial \tau^*_{\ell+1} / \partial A_{\ell+1} > 0$ (Proposition~\ref{prop:precision_response}). Because $V(\cdot)$ is increasing and concave, the composition preserves supermodularity. By Topkis's theorem, the optimal allocation is front-loaded.
\end{proof}

\begin{corollary}[Back-Loading Boundary]\label{cor:back_loading}
Back-loading can be optimal when the value function is convex in precision over the relevant range, or when the final action has state-dependent stakes that increase with depth. Formally, if $V''(\tau) > 0$ for $\tau \in [\underline{\tau}, \bar{\tau}]$, then there exist parameter values for which $A^*_0 < A^*_L$.
\end{corollary}

\begin{proof}
When $V''(\tau) > 0$, the composition $V(\tau^*_L)$ may be submodular in $(A_\ell, A_{\ell'})$ because the convexity of $V$ can dominate the concavity of $\tau^*_L$. If the net cross-partial is negative, Topkis's theorem implies back-loading. A numerical example is provided in the supplementary simulation experiments.
\end{proof}

\begin{remark}[Robustness to functional form]
The qualitative comparative-statics conclusions (Propositions \ref{prop:distortion_comp_statics}--\ref{prop:cost_min}) do not depend on the specific functional form of $g_\ell(\cdot)$. All that is required is that $g_\ell$ be increasing and concave with $g'_\ell(0) > 0$. The power-function specification $g_\ell(\alpha) = \alpha^\gamma$ is adopted for analytical convenience and because it yields closed-form solutions; the same economic forces operate under any attention technology that exhibits diminishing marginal returns.
\end{remark}

\subsection{Benchmark Model: The GHM Framework}
\label{sec:GHM_benchmark}

To isolate the distinctive predictions of our attention-allocation model, we now construct and solve a benchmark formulation that nests our setting within the classical framework of \citet{garicano2000hierarchies} and \citet{holmstrom1991multitask} (henceforth, GHM).

\subsubsection{Model Setup}

Consider a chain of $L$ identical agents indexed by $\ell = 1, \dots, L$. Each agent $\ell$ chooses a scalar effort level $e_\ell \in \mathbb{R}_+$, which represents the precision or quality of the signal that agent $\ell$ transmits upstream.

\begin{assumption}[GHM Cost Structure]\label{ass:GHM_cost}
The cost function $\psi$ is twice continuously differentiable with:
\begin{enumerate}[label=(\roman*)]
\item $\psi(0) = 0$, $\psi'(e) > 0$ for all $e \geq 0$;
\item $\psi''(e) > 0$ for all $e \geq 0$ (strict convexity);
\item $\lim_{e \to \infty} \psi'(e) = \infty$ (Inada condition).
\end{enumerate}
\end{assumption}

\begin{assumption}[GHM Production Technology]\label{ass:GHM_production}
The final output quality is given by
\begin{equation}\label{eq:GHM_production}
q_L = f\Biggl(\sum_{\ell=1}^{L} e_\ell\Biggr) + \varepsilon, \qquad \varepsilon \sim \mathcal{N}(0, \sigma_\varepsilon^2),
\end{equation}
where $f: \mathbb{R}_+ \to \mathbb{R}_+$ is twice continuously differentiable with $f'(x) > 0$, $f''(x) < 0$ for all $x \geq 0$, and satisfies the Inada conditions $f'(0) = \infty$, $\lim_{x \to \infty} f'(x) = 0$.
\end{assumption}

Agent $\ell$'s utility is:
\begin{equation}\label{eq:GHM_utility}
U_\ell = \mathbb{E}[w(q_L)] - \psi(e_\ell) = a + b \cdot f\Biggl(\sum_{j=1}^{L} e_j\Biggr) - \psi(e_\ell).
\end{equation}

\subsubsection{Equilibrium Effort and Optimal Depth}

Taking the depth $L$ and incentive power $b$ as given, each agent $\ell$ chooses $e_\ell$ to maximize \eqref{eq:GHM_utility}. By symmetry, $e_\ell = e^*$ for all $\ell$. The first-order condition is:
\begin{equation}\label{eq:GHM_FOC}
b \cdot f'(L \cdot e^*) = \psi'(e^*).
\end{equation}

\begin{proposition}[GHM Optimal Depth]\label{prop:GHM_depth}
Suppose Assumptions \ref{ass:GHM_cost}--\ref{ass:GHM_production} hold and the principal optimally sets $b^* > 0$. Then:
\begin{enumerate}[label=(\roman*)]
\item There exists a finite optimal depth $L^*_{\text{GHM}} < \infty$.
\item As the marginal cost of effort vanishes---formally, as $\psi'(e) \to 0$---the optimal depth diverges:
\begin{equation}\label{eq:GHM_infinite_depth}
\lim_{\kappa \to 0} L^*_{\text{GHM}}(\kappa) = \infty.
\end{equation}
\item The equilibrium effort per layer satisfies $e^* \to \infty$ as $\kappa \to 0$.
\end{enumerate}
\end{proposition}

\begin{proof}
(i) At $L = 0$, the marginal benefit is strictly positive because $f'(0) = \infty$. As $L \to \infty$, $f'(L \cdot e^*) \to 0$ while $\psi(e^*) > 0$. By continuity, there exists at least one crossing.

(ii) Consider the family $\psi_\kappa(e) = \kappa \cdot \tilde{\psi}(e)$. For any fixed finite $L$, as $\kappa \to 0$, the FOC implies $f'(L \cdot e^*) = (\kappa/b) \tilde{\psi}'(e^*) \to 0$, which requires $L \cdot e^* \to \infty$. The marginal benefit is $O(1)$ while the marginal cost is $O(\kappa)$. Thus the principal always wants to add more layers, and $L^*_{\text{GHM}} \to \infty$.
\end{proof}

\subsubsection{Complementary Predictions}

In our attention-allocation model, even as the per-unit cost of attention $\kappa \to 0$, each layer faces a \emph{hard budget constraint} $\sum_k \alpha_{\ell,k} \leq A_\ell$. The optimal depth remains finite because:
\begin{enumerate}[label=(\alph*)]
\item The budget $A_\ell$ places an upper bound on effective precision.
\item The transmission factor $\rho_\ell = \tau^*_\ell / \tau^{\text{FB}}_\ell < 1$ is bounded away from 1 by the geometry of the budget constraint, regardless of $\kappa$.
\item The geometric decay of precision along the chain is driven by the \emph{relative} allocation across signals, not by the absolute cost of attention.
\end{enumerate}

\begin{table*}[htbp]
\centering
\footnotesize
\setlength{\tabcolsep}{5pt}
\caption{Model Predictions: Our Framework vs.~GHM Benchmark}
\label{tab:prediction_comparison}
\begin{tabular}{@{}>{\raggedright\arraybackslash}p{2.8cm}>{\raggedright\arraybackslash}p{3.5cm}>{\raggedright\arraybackslash}p{3.5cm}>{\raggedright\arraybackslash}p{3.5cm}@{}}
\toprule
\textbf{Prediction} & \textbf{Our Model} & \textbf{GHM Benchmark} & \textbf{Empirical Test} \\
\midrule
Finite optimal depth ($\kappa \to 0$) & \textbf{Yes.} Budget $A_\ell$ bounds precision regardless of cost. & \textbf{No.} $L^*_{\text{GHM}} \to \infty$ as $\psi'(e) \to 0$. & Observe depth response to API price cuts. \\
Budget allocation & \textbf{Front-loaded.} Attention concentrates at early layers. & \textbf{Uniform.} Identical effort $e^*$ at all layers. & Compare model sizes across layers. \\
Signal heterogeneity & \textbf{Reduces distortion.} Selective attention raises $\rho_\ell$. & \textbf{No effect.} Heterogeneity not in FOC. & Homogeneous vs.~heterogeneous inputs. \\
Depth--precision relation & \textbf{Geometric decay.} $\tau_\ell \propto \rho^{\ell}$. & \textbf{Arithmetic decay.} $L \cdot e^*(L)$. & Error accumulation in chains. \\
Elasticity identification & \textbf{Identifiable.} From context-accuracy curves. & \textbf{Not identifiable.} $e_\ell$ unobservable. & Attention-weight data from APIs. \\
\bottomrule
\end{tabular}
\end{table*}

\subsubsection{Four Dimensions of Complementarity}
\label{sec:key_distinctions}

The predictions in Table \ref{tab:prediction_comparison} differ because the two models emphasize different economic primitives. GHM abstracts from attention allocation across heterogeneous signals, while our model abstracts from the team-production and incentive-design problems that GHM was built to analyze. The two frameworks are \emph{complements}: each generates predictions the other does not.

\paragraph{1. Vector vs.~scalar optimization.} Our model generates additional predictions by incorporating \emph{vector optimization}: each layer allocates a budget $A_\ell$ across multiple signals with heterogeneous precisions. This generates \textbf{selective neglect}---the optimal policy concentrates attention on the most informative signals and ignores others. GHM, by abstracting from signal heterogeneity, focuses on \emph{scalar optimization} where the agent chooses a single effort level $e_\ell$.

\paragraph{2. Endogenous vs.~exogenous distortion.} Our model generates additional predictions by making the transmission factor $\rho_\ell = \tau^*_\ell / \tau^{\text{FB}}_\ell$ \textbf{endogenous}: it depends on the signal structure, the attention technology $g_\ell(\cdot)$, and the budget $A_\ell$. GHM abstracts from this channel; any distortion in that framework is \textbf{exogenous}---it would appear as a fixed parameter in the production function $f$.

\paragraph{3. Observable vs.~unobservable inputs.} The attention weights $\alpha_{\ell,k}$ in our model capture, in reduced form, how the model processes different information sources. Modern transformer architectures expose attention maps as byproducts of computation \citep{vaswani2017attention}, which enables direct structural estimation. GHM, by focusing on moral hazard, treats effort $e_\ell$ as classically \textbf{unobservable}.

\paragraph{4. Behavior as $\kappa \to 0$: a complementary prediction.} Proposition \ref{prop:GHM_depth} establishes that $L^*_{\text{GHM}} \to \infty$ when $\psi'(e) \to 0$. Our model generates the complementary prediction that $L^* < \infty$ \emph{even in the limit} because the budget constraint $A_\ell$ creates a hard bound on effective precision. Recent API price reductions represent a shock to $\kappa$. Under GHM, firms should respond by deepening chains indefinitely; under our model, depth should stabilize. The evidence in Section \ref{sec:empirics} supports the latter prediction.

\begin{remark}[What GHM does better]
The GHM framework naturally accommodates \emph{team production} with free-riding, \emph{optimal incentive design} through the choice of $b$, and \emph{rich contracting} between principal and agents. Our model abstracts from these classical moral-hazard considerations to focus on the novel margin of attention allocation. A fully satisfactory theory would embed our attention-allocation problem within a GHM-style contracting framework; we leave this synthesis to future work.
\end{remark}

\section{Multi-Layer Analysis and Optimal Depth}
\label{sec:multi}

We now extend the analysis to multiple layers, derive the optimal organizational depth, and characterize the optimal budget allocation.

\subsection{The Endogenous Transmission Factor}
\label{subsec:transmission_factor}

Before analyzing information transmission across layers, we formally define the endogenous transmission factor and relate it to the single-layer distortion measure.

Recall from Definition~\ref{def:single_layer_distortion} that the single-layer distortion measure is $\tilde{\delta}_\ell = 1 - \Delta_\ell$. In the multi-layer context, it is more convenient to work directly with the ratio of achieved posterior precision to first-best precision:

\begin{definition}[Endogenous Transmission Factor]\label{def:transmission_factor}
The endogenous transmission factor from layer $\ell$ to layer $\ell+1$ is
\begin{equation}\label{eq:transmission_factor}
\rho_\ell \equiv \frac{\tau^*_\ell}{\tau^{\text{FB}}_\ell} = \frac{\tau^{\text{prior}}_\ell + \sum_{k=1}^{K_\ell} g_\ell(\alpha^*_{\ell,k}) \tau^0_{\ell,k}}{\tau^{\text{prior}}_\ell + \sum_{k=1}^{K_\ell} \tau^0_{\ell,k}} \in [0,1],
\end{equation}
where $\tau^{\text{FB}}_\ell$ is the first-best posterior precision and $\tau^*_\ell$ is the achieved posterior precision under the optimal attention allocation.
\end{definition}

The transmission factor $\rho_\ell$ measures the fraction of first-best precision that survives layer $\ell$. When $\rho_\ell = 1$, all signals receive sufficient attention and no precision is lost. When $\rho_\ell < 1$, the attention budget forces the layer to ignore or under-process some signals, creating a precision leakage.

\begin{remark}[From Exogenous to Endogenous]\label{rem:exog_to_endog}
In previous work, the transmission factor was an exogenous parameter $\tau \in [0,1]$ that the modeler specified. In our framework, $\rho_\ell$ is endogenously determined by the attention-allocation problem: it depends on the budget $A_\ell$, the number of signals $K_\ell$, their precisions $\{\tau^0_{\ell,k}\}$, and the attention technology $g_\ell(\cdot)$. Importantly, $\rho_\ell$ is derived \emph{relative to the assumed attention technology $g_\ell(\cdot)$}---changing the functional form of $g_\ell$ would alter the numerical value of $\rho_\ell$ even holding all primitives fixed.
\end{remark}

\begin{remark}[Technology Dependence and Qualitative Robustness]\label{rem:rho_technology}
The numerical value of $\rho_\ell$ depends on the specific attention technology $g_\ell(\cdot)$. Under the power-function specification $g_\ell(\alpha) = \alpha^\gamma$, a larger context window (higher $A_\ell$) raises $\rho_\ell$; under an alternative technology, the quantitative response may differ. However, the qualitative conclusions of our analysis---that optimal depth is finite and that budget allocation is front-loaded---do not depend on the specific functional form of $g_\ell(\cdot)$. All that is required is that $g_\ell$ be increasing and concave with $g'_\ell(0) > 0$.
\end{remark}

\subsection{Organizational Design: Optimal Budget Allocation}
\label{subsec:budget_allocation}

Before analyzing the finite-depth result, we first characterize the optimal budget allocation across layers. This is the model's most empirically testable prediction.

Suppose the principal faces a total attention budget $B > 0$ that can be allocated across layers. For simplicity, consider a uniform technology profile where $g_\ell(\cdot) = g(\cdot)$, $c_\ell(\cdot) = c(\cdot)$, and $K_\ell = K$ for all $\ell$. The principal's problem is:
\begin{equation}\label{eq:budget_allocation}
\max_{L, \{A_\ell\}} V_L(\tau^*_L) \quad \text{s.t.} \quad \sum_{\ell=0}^{L} A_\ell \leq B, \quad A_\ell \geq 0.
\end{equation}

We compare three canonical budget-allocation strategies:

\textbf{Strategy 1: Uniform allocation.} Set $A_\ell = B/(L+1)$ for all $\ell$. This is the simplest design: every layer receives the same context-window capacity.

\textbf{Strategy 2: Front-loaded allocation.} Set $A_0 > A_1 > \cdots > A_L$. The bottom layers receive larger budgets because they process raw signals with the highest intrinsic precision. The intuition is that precision loss is most costly early in the chain: a unit of precision lost at layer 0 is compounded by all subsequent layers.

\textbf{Strategy 3: Back-loaded allocation.} Set $A_0 < A_1 < \cdots < A_L$. The top layers receive larger budgets because they make the final decision and their precision directly enters the payoff.

\begin{proposition}[Optimal Budget Allocation]\label{prop:front_loading}
Under the uniform-technology assumption, the optimal budget allocation is front-loaded: $A^*_0 \geq A^*_1 \geq \cdots \geq A^*_L$.
\end{proposition}

\begin{proof}
Consider two adjacent layers $\ell$ and $\ell+1$ with budgets $A_\ell$ and $A_{\ell+1}$. Define an alternative allocation that transfers $\varepsilon > 0$ from layer $\ell+1$ to layer $\ell$: $A'_\ell = A_\ell + \varepsilon$, $A'_{\ell+1} = A_{\ell+1} - \varepsilon$. We show that if $A_\ell = A_{\ell+1}$, then $\tau^*_L(A'_\ell, A'_{\ell+1}) > \tau^*_L(A_\ell, A_{\ell+1})$ for small $\varepsilon$.

The transmission factor at layer $\ell$ increases: $\rho_\ell(A'_\ell) > \rho_\ell(A_\ell)$ because $\partial \rho_\ell / \partial A_\ell > 0$ (Proposition~\ref{prop:distortion_comp_statics}). This raises the inherited precision at all downstream layers $j > \ell$.

The transmission factor at layer $\ell+1$ decreases: $\rho_{\ell+1}(A'_{\ell+1}) < \rho_{\ell+1}(A_{\ell+1})$. However, the inherited precision at layer $\ell+1$ is $\tau^{\text{prior}}_{\ell+1} = \rho_\ell \cdot \tau^*_\ell$, which increases because $\rho_\ell$ increases. For small $\varepsilon$, the increase in $\rho_\ell$ dominates the decrease in $\rho_{\ell+1}$ because the compounding effect amplifies the gain at layer $\ell$ across all downstream layers.

Repeating the exchange argument for all adjacent pairs from $\ell = 0$ to $\ell = L-1$, any non-front-loaded allocation can be improved by shifting budget toward earlier layers. The only allocation immune to such improvements is the front-loaded one.
\end{proof}

\vspace{1em}
\noindent\textbf{Practical Implications.}
In LLM pipeline design, Proposition~\ref{prop:front_loading} implies that the largest model (largest context window) should be deployed at the earliest stage---the layer that processes the raw user query and initial retrieved documents. Smaller, faster models can be used at intermediate layers.

This prediction aligns with the architecture of several prominent production multi-agent systems. \citet{openai2025deepresearch} deploys a large reasoning model (o3, 200K context window) for the initial query-understanding and web-retrieval phase, while using lighter models for intermediate synthesis steps. \textit{Manus} \citep{manus2025} assigns its largest model to the initial task-decomposition layer, where raw user intent must be parsed and mapped to a tool-execution plan. In open-source tool-use frameworks, \textit{OpenHands} \citep{openhands2024} and \textit{Devin} \citep{cognition2024devin} both place a frontier model at the entry point of the pipeline (code comprehension and planning), while delegating implementation steps to smaller specialized agents. Our model provides a theoretical justification for this engineering practice: the returns to attention capacity are highest at the bottom of the chain because precision lost there is amplified by every subsequent layer.

\begin{remark}[Front-loading as an Empirical Signature]\label{rem:front_loading_empirical}
The front-loading result is the model's most directly testable prediction. It does not require measuring the preservation parameters $(\underline{\eta}, \bar{\eta}, a)$ or the transmission factor $\rho_\ell$, which are difficult to observe in practice. Instead, it predicts a monotonic relationship between layer index and budget allocation that can be evaluated using only observable architectural choices: the model sizes or context-window lengths assigned to each layer of a multi-agent pipeline. We return to this prediction in Section~\ref{sec:empirics}, where we discuss how to bring it to data.
\end{remark}

\subsection{Information Transmission Across Layers}
\label{subsec:info_transmission}

Having characterized the optimal budget allocation, we now study how information propagates through the chain. This analysis reveals the dynamic forces that ultimately bound the optimal depth.

In a chain with $L$ layers, the posterior precision at layer $\ell$ depends on the attention allocation at layer $\ell-1$ through the endogenous transmission factor $\rho_{\ell-1}$. The prior precision at layer $\ell$ is the achieved posterior precision at layer $\ell-1$:
\begin{equation}\label{eq:prior_update}
\tau^{\text{prior}}_\ell = \rho_{\ell-1} \cdot \tau^*_{\ell-1}(A_{\ell-1}), \quad \ell = 1, \ldots, L.
\end{equation}
With $\tau^{\text{prior}}_0 \equiv \tau^0_0$ (the raw signal precision at the bottom layer), the posterior precision at layer $\ell$ follows the recursive Bayesian-update structure:
\begin{equation}\label{eq:recursive_precision}
\tau^*_\ell = \tau^{\text{prior}}_\ell + \sum_{k=1}^{K_\ell} g_\ell(\alpha^*_{\ell,k}) \tau^0_{\ell,k}, \quad \text{with } \tau^{\text{prior}}_\ell = \rho_{\ell-1} \cdot \tau^*_{\ell-1}.
\end{equation}

Solving the recursion yields the cumulative precision at the top layer:
\begin{equation}\label{eq:top_precision}
\tau^*_L = \sum_{\ell=1}^{L} \sum_{k=1}^{K_\ell} g_\ell(\alpha^*_{\ell,k}) \tau^0_{\ell,k} \cdot \prod_{j=\ell}^{L-1} \rho_j + \tau^0_0 \cdot \prod_{j=0}^{L-1} \rho_j.
\end{equation}
Each layer $\ell$ contributes its local signal precision $g_\ell(\alpha^*_{\ell,k}) \tau^0_{\ell,k}$, weighted by the subsequent transmission losses $\prod_{j=\ell}^{L-1} \rho_j$.

\vspace{1em}
\noindent\textbf{Precision Dynamics under Information Depreciation.}
We now characterize the dynamics of the precision path $\{\tau^*_\ell\}$ under the information depreciation model introduced in Assumptions~\ref{ass:state_signals} and~\ref{ass:endogenous_eta}. Recall that in this framework, the raw signal quality at layer $\ell$ is $\tau^0_{\ell,k} = \tau^0_k \cdot \eta_\ell^{\,\ell}$, where $\eta_\ell \in (\underline{\eta}, \bar{\eta})$ is the endogenous per-layer signal preservation rate at layer $\ell$, determined by the attention budget $A_\ell$ through equation~\eqref{eq:endogenous_eta}.

To build intuition, consider first what happens when $\eta_\ell = 1$ for all $\ell$ (no signal degradation). Under uniform technology and uniform attention budgets ($g_\ell = g$, $c_\ell = c$, $K_\ell = K$, $\rho_\ell = \rho$, $A_\ell = A$, so $\eta_\ell = \eta$ constant), Equation~\eqref{eq:recursive_precision} becomes $\tau^*_\ell = \rho \tau^*_{\ell-1} + C$ for some constant $C > 0$, which converges to $\bar{\tau} = C/(1-\rho) > 0$ as $\ell \to \infty$. In this case, the precision path converges to a \emph{positive limit}: adding more layers continues to increase precision (albeit with diminishing returns), and infinite depth is optimal whenever API costs are below a threshold. The boundedness of depth, if any, comes solely from cost considerations, not from fundamental information limits.

When $\eta_\ell < 1$, the dynamics change qualitatively. The following proposition characterizes the precision path under the information depreciation model. To maintain analytical transparency, we state the result for the uniform-technology case with potentially non-uniform attention budgets; the special case of uniform budgets ($A_\ell = A$, so $\eta_\ell = \eta$) recovers the baseline formula.

\begin{proposition}[Precision Path Characterization]\label{prop:precision_path}
Under the information depreciation model (Assumptions~\ref{ass:state_signals}--\ref{ass:endogenous_eta} with $\eta_\ell \in (\underline{\eta}, \bar{\eta})$) and uniform technology ($g_\ell = g$, $c_\ell = c$, $K_\ell = K$, $\rho_\ell = \rho$ for all $\ell$), define the composite signal contribution constant
\begin{equation}\label{eq:C_constant}
C \equiv \sum_{k=1}^{K} g\bigl(\alpha^{*}_{k}\bigr) \tau^{0}_{k},
\end{equation}
where $\alpha^{*}_{k}$ is the optimal attention allocation to source $k$. Let $\bar{\eta} \equiv \max_\ell \eta_\ell$ denote the maximal preservation rate across layers. The equilibrium posterior precision at layer $\ell$ satisfies
\begin{equation}\label{eq:recursive_tau_unified}
\tau^{*}_{\ell} = \rho \cdot \tau^{*}_{\ell-1} + C \eta_\ell^{\,\ell}, \qquad \tau^{*}_{0} = \sum_{k=1}^{K} \tau^{0}_{k}.
\end{equation}
Under uniform attention budgets ($A_\ell = A$, so $\eta_\ell = \eta$ for all $\ell$), the closed-form solution is
\begin{equation}\label{eq:tau_star_L}
\tau^{*}_{L} = \rho^{L} \tau^{*}_{0} + \frac{C\eta\bigl(\eta^{L} - \rho^{L}\bigr)}{\eta - \rho} \quad \text{for } \eta \neq \rho,
\end{equation}
and $\tau^{*}_{L} = \eta^{L}\bigl(\tau^{*}_{0} + CL\bigr)$ when $\eta = \rho$. Moreover:

\begin{enumerate}[label=(\roman*)]
\item \textbf{If $\eta < \rho$:} $\tau^{*}_{\ell}$ is hump-shaped. It increases for early layers, reaches a maximum at some finite $\ell^{*} \geq 1$, then decreases monotonically. The limiting precision is $\lim_{\ell \to \infty} \tau^{*}_{\ell} = 0$.

\item \textbf{If $\eta > \rho$:} $\tau^{*}_{\ell}$ is eventually decreasing. It may increase initially depending on initial conditions, but eventually declines monotonically to $0$.

\item \textbf{If $\eta = \rho$:} $\tau^{*}_{\ell} = \eta^{\ell}(\tau^{*}_{0} + C\ell)$, which is eventually decreasing to $0$.
\end{enumerate}
In all cases, $\lim_{\ell \to \infty} \tau^{*}_{\ell} = 0$.
\label{prop:geometric_decay}%
\end{proposition}

\begin{proof}
\noindent\textbf{Step 1: Derivation of the recurrence.}
Under uniform technology and Assumptions~\ref{ass:state_signals}--\ref{ass:endogenous_eta}, the new signal contribution at layer $\ell$ is $\sum_{k=1}^{K} g(\alpha^*_k) \tau^0_k \eta_\ell^{\,\ell} = C\eta_\ell^{\,\ell}$, where $C$ is defined in \eqref{eq:C_constant}. Substituting into \eqref{eq:recursive_precision} yields \eqref{eq:recursive_tau_unified}. Under uniform attention budgets ($A_\ell = A$), we have $\eta_\ell = \eta$ for all $\ell$, and the recurrence simplifies to $\tau^{*}_{\ell} = \rho \tau^{*}_{\ell-1} + C\eta^{\ell}$.

\noindent\textbf{Step 2: Closed-form solution (uniform budgets).}
Under uniform budgets, \eqref{eq:recursive_tau_unified} becomes a first-order linear non-homogeneous difference equation with constant coefficients.

\textit{Case 1: $\eta \neq \rho$.} The homogeneous solution is $\tau^{(h)}_{\ell} = A \rho^{\ell}$. Guessing a particular solution $\tau^{(p)}_{\ell} = B \eta^{\ell}$ and substituting:
\begin{equation*}
B\eta^{\ell} = \rho B \eta^{\ell-1} + C\eta^{\ell}
\quad\Longrightarrow\quad
B = \frac{C\eta}{\eta - \rho}.
\end{equation*}
Applying $\tau^{*}_{0} = A + B$ and solving for $\tau^{*}_{L}$ yields \eqref{eq:tau_star_L}.

\textit{Case 2: $\eta = \rho$.} The recurrence becomes $\tau^{*}_{\ell} = \eta \tau^{*}_{\ell-1} + C\eta^{\ell}$. Defining $u_{\ell} \equiv \tau^{*}_{\ell}/\eta^{\ell}$, we have $u_{\ell} = u_{\ell-1} + C$, which telescopes to $u_{\ell} = \tau^{*}_{0} + C\ell$. Hence $\tau^{*}_{L} = \eta^{L}(\tau^{*}_{0} + CL)$.

\noindent\textbf{Step 3: Vanishing precision (all cases).}
Since $\rho \in (0,1)$ and $\eta_\ell \leq \bar{\eta} < 1$ for all $\ell$ (by Assumption~\ref{ass:endogenous_eta}), both $\rho^{L} \to 0$ and $\eta_\ell^{\,L} \to 0$. In all three cases, $\lim_{L\to\infty} \tau^{*}_{L} = 0$. The critical difference from the $\eta = 1$ case is that $\bar{\eta} < 1$ is a \emph{structural bound}: even with unlimited attention, the preservation rate cannot reach unity, and geometric decay is inevitable.

\noindent\textbf{Step 4: Shape characterization.}
From \eqref{eq:recursive_tau_unified}, the layer-to-layer change is
\begin{equation}\label{eq:delta_tau}
\Delta_{\ell} \equiv \tau^{*}_{\ell} - \tau^{*}_{\ell-1} = -(1-\rho)\tau^{*}_{\ell-1} + C\eta_\ell^{\,\ell}.
\end{equation}
Under uniform budgets ($\eta_\ell = \eta$):

\textit{Subcase (i): $\eta < \rho$.} The $\rho^{\ell}$ term dominates the closed-form solution for large $\ell$. Since $\rho > \eta$, the negative term $-(1-\rho)\tau^{*}_{\ell-1}$ (which decays as $\rho^{\ell-1}$) eventually dominates the positive term $C\eta^{\ell}$ (which decays as $\eta^{\ell}$). Thus $\Delta_{\ell} < 0$ for all $\ell \geq \bar{L}_1$. For small $\ell$, however, the positive term may dominate if $\tau^{*}_{0}$ is small relative to $C\eta$, causing an initial increase. The path is therefore hump-shaped.

\textit{Subcase (ii): $\eta > \rho$.} As $\ell \to \infty$, $\tau^{*}_{\ell-1} \sim C\eta^{\ell}/(\eta-\rho)$. Substituting into \eqref{eq:delta_tau}:
\begin{equation*}
\Delta_{\ell} \sim C\eta^{\ell}\left(1 - \frac{1-\rho}{\eta-\rho}\right) = C\eta^{\ell}\frac{\eta - 1}{\eta - \rho} < 0,
\end{equation*}
since $\eta < 1$ and $\eta > \rho$. The sequence is eventually decreasing.

\textit{Subcase (iii): $\eta = \rho$.} Using $\tau^{*}_{\ell} = \eta^{\ell}(\tau^{*}_{0} + C\ell)$:
\begin{equation*}
\Delta_{\ell} = \eta^{\ell-1}\bigl[(\eta-1)\tau^{*}_{0} + C - C(1-\eta)\ell\bigr].
\end{equation*}
Since $\eta < 1$, the bracket is affine in $\ell$ with negative slope $-C(1-\eta) < 0$. Hence $\Delta_{\ell} < 0$ whenever $\ell > \tau^{*}_{0}/C + 1/(1-\eta)$.
\end{proof}

\begin{proposition}[Endogenous Preservation and the $\bar{\eta}$ Bound]\label{prop:eta_bar_bound}
Under Assumption~\ref{ass:endogenous_eta}, the preservation rate $\eta_\ell$ satisfies:
\begin{enumerate}[label=(\roman*)]
\item \textbf{Attention monotonicity:} $\eta_\ell$ is strictly increasing and concave in $A_\ell$:
\begin{equation}\label{eq:eta_increasing}
\frac{\partial \eta_\ell}{\partial A_\ell} = \frac{(\bar{\eta} - \underline{\eta}) \, a / K_\ell}{(\bar{\alpha}_\ell + a)^2} > 0, \qquad \frac{\partial^2 \eta_\ell}{\partial A_\ell^2} < 0.
\end{equation}

\item \textbf{Irreducible depreciation:} Even with unlimited attention, depreciation cannot be eliminated:
\begin{equation}\label{eq:eta_bar_limit}
\lim_{A_\ell \to \infty} \eta_\ell = \bar{\eta} < 1.
\end{equation}
Consequently, for any attention budget profile $\{A_\ell\}$, the cumulative precision at the top layer satisfies
\begin{equation}\label{eq:tau_L_upper_bound}
\tau^{*}_L \leq \rho^{L} \tau^{*}_{0} + \frac{C\bar{\eta}\bigl(\bar{\eta}^{L} - \rho^{L}\bigr)}{\bar{\eta} - \rho} \quad (\bar{\eta} \neq \rho),
\end{equation}
and the upper bound converges to $0$ as $L \to \infty$.

\item \textbf{Finite depth robustness:} The finite-optimal-depth result (Proposition~\ref{prop:optimal_depth}) holds \emph{regardless} of the attention budget profile. Even if $A_\ell \to \infty$ at every layer, the bound $\bar{\eta} < 1$ ensures $\lim_{L \to \infty} \tau^{*}_L = 0$.
\end{enumerate}
\end{proposition}

\begin{proof}
\noindent\textbf{Part (i):} Direct differentiation of \eqref{eq:endogenous_eta} with respect to $A_\ell$, using $\bar{\alpha}_\ell = A_\ell / K_\ell$, yields \eqref{eq:eta_increasing}.

\noindent\textbf{Part (ii):} Taking the limit of \eqref{eq:endogenous_eta} as $A_\ell \to \infty$:
\begin{equation*}
\lim_{A_\ell \to \infty} \eta_\ell = \underline{\eta} + (\bar{\eta} - \underline{\eta}) \cdot \lim_{\bar{\alpha}_\ell \to \infty} \frac{\bar{\alpha}_\ell}{\bar{\alpha}_\ell + a} = \underline{\eta} + (\bar{\eta} - \underline{\eta}) \cdot 1 = \bar{\eta}.
\end{equation*}
Since $\eta_\ell \leq \bar{\eta}$ for all $\ell$, the signal contribution at layer $\ell$ is bounded above by $C\bar{\eta}^{\,\ell}$. Substituting this upper bound into the recurrence \eqref{eq:recursive_tau_unified} yields \eqref{eq:tau_L_upper_bound}. Since $\bar{\eta} < 1$, the upper bound converges to $0$ as $L \to \infty$.

\noindent\textbf{Part (iii):} By Part (ii), $\tau^{*}_L$ is bounded above by a quantity that converges to $0$, so $\tau^{*}_L \to 0$ for \emph{any} budget profile. The proof of Proposition~\ref{prop:optimal_depth} then applies unchanged: there exists a finite $\bar{L}$ beyond which the marginal benefit of adding a layer is negative.
\end{proof}

\begin{remark}[Why $\bar{\eta} < 1$: Micro-foundations]\label{rem:eta_bar_microfoundation}
The bound $\bar{\eta} < 1$ is not an exogenous assumption; it reflects three fundamental limits of LLM information processing that no amount of attention can overcome:
\begin{enumerate}[label=(\roman*)]
\item \textbf{Irreversible embedding compression.} Each forward pass projects discrete tokens into a continuous embedding space through a non-invertible mapping. The entropy reduction is structural: a 4096-dimensional embedding of a 100,000-token vocabulary inherently discards information that cannot be recovered by downstream layers, regardless of context-window size.
\item \textbf{Non-invertible attention operations.} The attention mechanism computes weighted averages of value vectors. This is a many-to-one mapping: different input configurations can produce identical attention outputs. The operation discards positional and structural information about \emph{which} tokens contributed \emph{how much}, and this information loss is fundamental to the architecture.
\item \textbf{Lossy summarization as a design feature.} LLMs are trained to compress and generalize, not to preserve. The training objective (next-token prediction with KL regularization) explicitly rewards discarding idiosyncratic detail in favor of distributional patterns. Even with unlimited compute and context, the model's inductive bias favors abstraction over fidelity.
\end{enumerate}
These three forces explain why $\bar{\eta} < 1$ is a \emph{technological primitive} of transformer-based processing, not a modeling assumption that could be relaxed by better engineering. The bound $\bar{\eta}$ can be calibrated from empirical measurements of signal fidelity as a function of context length; see Section~\ref{sec:calibration}.
\end{remark}

\begin{remark}[What the Compounding Effect Is and Is Not]\label{rem:compounding_qualitative}
An earlier version of this paper claimed that $\tau^*_L$ is \emph{strictly decreasing} in $L$ for any fixed budget (a ``compounding effect''). That claim is incorrect. Proposition~\ref{prop:precision_path} shows that $\tau^*_\ell$ can \emph{increase} for early layers when fresh signals are abundant ($C\eta_\ell^{\,\ell}$ large relative to inherited precision). The correct statement is that precision \emph{eventually} declines due to information depreciation, not that it decreases monotonically from the first layer. The eventual decline to zero is the economically relevant property for establishing finite optimal depth.
\end{remark}

\subsection{Optimal Organizational Depth}
\label{subsec:optimal_depth}

Proposition~\ref{prop:precision_path} establishes that the precision path eventually declines to zero under information depreciation. We now show that this vanishing precision implies a finite optimal depth.

The principal chooses depth $L$ and the budget profile $\{A_\ell\}_{\ell=0}^{L-1}$ to maximize the value of the final decision. The optimization problem is:
\begin{equation}\label{eq:optimal_depth}
\max_{L, \{A_\ell\}} V_L(\tau^*_L) - \sum_{\ell=0}^{L-1} c_\ell(A_\ell).
\end{equation}

\begin{proposition}[Optimal Depth Existence]\label{prop:optimal_depth}
\label{prop:depth_finite}%
\label{prop:finite_depth}%
Under the information depreciation model (Assumptions~\ref{ass:hierarchy}--\ref{ass:endogenous_eta} with $\bar{\eta} < 1$), there exists a finite optimal depth $L^* < \infty$ for any attention budget profile $\{A_\ell\}$.
\end{proposition}

\begin{proof}
The proof has two parts. First, we show that $\tau^{*}_{\ell} \to 0$ as $\ell \to \infty$, which follows directly from Proposition~\ref{prop:precision_path} and the bound $\eta_\ell \leq \bar{\eta} < 1$ (Proposition~\ref{prop:eta_bar_bound}). Second, we show that this vanishing precision implies a finite optimal depth.

Since $V_L(\cdot)$ is strictly increasing and continuous, $\lim_{\ell \to \infty} V_L(\tau^{*}_{\ell}) = V_L(0)$. By Proposition~\ref{prop:precision_path}, there exists a finite threshold $\bar{L}$ such that $\tau^{*}_{\ell} < \tau^{*}_{\ell-1}$ for all $\ell \geq \bar{L}$. Since $V_L$ is strictly increasing, $V_L(\tau^{*}_{\ell}) < V_L(\tau^{*}_{\ell-1})$ for all $\ell \geq \bar{L}$.

For any $\ell \geq \bar{L}$, the net marginal benefit of adding layer $\ell$ is:
\begin{equation*}
\underbrace{V_L(\tau^{*}_{\ell}) - V_L(\tau^{*}_{\ell-1})}_{< 0} - \underbrace{c_{\ell}(A_{\ell})}_{\geq 0} < 0.
\end{equation*}
Thus no planner would ever choose $\ell > \bar{L}$, and $L^{*} \leq \bar{L} < \infty$ regardless of how small the marginal attention cost $\kappa$ is or how large the attention budgets $\{A_\ell\}$ are.

The critical mechanism is \emph{information depreciation bounded away from unity}: because $\eta_\ell \leq \bar{\eta} < 1$ for all $\ell$ (even with unlimited attention), each layer's signals carry less raw precision than the previous layer's ($\bar{\eta}^{\,\ell} \to 0$), while the attention budget can only raise $\eta_\ell$ toward $\bar{\eta}$ but never eliminate the geometric decay. The compounding of these two forces---inherent algorithmic noise (captured by $\bar{\eta} < 1$) and endogenous precision leakage ($\rho < 1$)---drives $\tau^{*}_{\ell}$ to zero and bounds the efficient depth of the chain.
\end{proof}

\vspace{1em}
\noindent\textbf{Necessity of the $\bar{\eta} < 1$ bound.}
The finite-depth result relies critically on Assumption~\ref{ass:endogenous_eta} with $\bar{\eta} < 1$. If $\bar{\eta} = 1$ (unbounded preservation with unlimited attention), then with sufficiently large $A_\ell$, $\eta_\ell$ could be driven arbitrarily close to $1$, and the recurrence \eqref{eq:recursive_tau_unified} would approach $\tau^{*}_{\ell} = \rho \tau^{*}_{\ell-1} + C$, with solution $\tau^{*}_{L} = \rho^{L}\tau^{*}_{0} + C(1-\rho^{L})/(1-\rho) \to C/(1-\rho) > 0$. The precision would converge to a positive limit, so $V(\tau^{*}_{L})$ converges to $V(C/(1-\rho)) > V(0)$ from below, and infinite depth could be optimal whenever the marginal cost $\kappa$ is below the limiting marginal benefit.

With the endogenous preservation rate bounded by $\bar{\eta} < 1$, this possibility is ruled out \emph{by construction}. Even as $A_\ell \to \infty$, we have $\eta_\ell \to \bar{\eta} < 1$, so the signal contribution at layer $\ell$ is bounded above by $C\bar{\eta}^{\,\ell}$, which decays geometrically. The bound $\bar{\eta} < 1$ is therefore \emph{sufficient} for the finite-depth result, \emph{regardless} of the attention cost or budget. This is the key sense in which endogenizing $\eta$ strengthens rather than weakens the model: the finite-depth result becomes robust to arbitrarily large attention budgets because inherent algorithmic noise places a technological ceiling on signal preservation.

\vspace{1em}
\noindent\textbf{Economic Discussion: Information Capital Depreciation.}

The critical insight of Proposition~\ref{prop:optimal_depth} is that \emph{depth-induced signal degradation} fundamentally alters the returns to layering. The endogenous preservation rate $\eta_\ell$ serves a role analogous to depreciation in capital accumulation: just as physical capital loses productive value through wear and tear, information capital (posterior precision) erodes through repeated algorithmic processing.

The preservation rate $\eta_\ell$ can be micro-founded through several channels:
\begin{enumerate}[label=(\roman*)]
\item \emph{Compressive loss:} Each LLM distills verbose inputs into concise outputs, inevitably discarding fine-grained detail. The summary statistics that propagate upward retain the gist but lose nuance.
\item \emph{Hallucination risk:} Models introduce plausible-sounding but factually incorrect content, injecting algorithmic noise that dilutes the signal-to-noise ratio.
\item \emph{Context window saturation:} Deeper layers must accommodate progressively more compressed representations; even with expanded context windows, the \emph{effective} precision per token declines.
\end{enumerate}

The endogenous framework clarifies the relationship between attention and depreciation. An $\eta_\ell$ close to $\bar{\eta}$ (achieved with large $A_\ell$) indicates high-fidelity transmission (e.g., retrieval-augmented systems with precise citations), while an $\eta_\ell$ close to $\underline{\eta}$ (small $A_\ell$) corresponds to lossy processing that obscures original signal content. Crucially, because $\bar{\eta} < 1$, even infinite attention cannot eliminate depreciation entirely. Our result that $L^{*}$ is bounded even when API calls are cheap ($\kappa \approx 0$) and attention budgets are large formalizes the intuition that organizational limits to AI delegation derive not from budget constraints alone, but from \emph{fundamental information-theoretic limits} inherent to transformer computation.

It is important to be precise about what drives the result. \emph{Attention scarcity alone} (finite $A_\ell$) is insufficient to guarantee finite optimal depth when costs vanish: if $\bar{\eta} = 1$, then with $A_\ell \to \infty$ we would have $\eta_\ell \to 1$, the precision would converge to a positive limit, and arbitrarily deep chains could be optimal. Conversely, \emph{information depreciation alone} ($\bar{\eta} < 1$) with unlimited attention budgets still bounds depth, as each layer's signals carry vanishingly little information no matter how much attention is applied. The two forces are \emph{complementary}: attention scarcity determines \emph{where} $\eta_\ell$ lies in the interval $[\underline{\eta}, \bar{\eta}]$, while the bound $\bar{\eta} < 1$ ensures that the \emph{total} information content always shrinks as the chain lengthens.

\subsection{Comparative Statics of Optimal Depth}
\label{subsec:comp_statics_depth}

We conclude with two comparative-statics results that link the optimal depth to economic primitives.

\begin{proposition}[Attention Budget and Depth]\label{prop:budget_depth}
The optimal depth $L^*$ is increasing in the attention budget profile:
\begin{equation}\label{eq:budget_depth}
\frac{\partial L^*}{\partial A_\ell} > 0 \quad \text{for all } \ell.
\end{equation}
\end{proposition}

\begin{proof}
A larger attention budget at any layer raises the transmission factor $\rho_\ell$, reducing the precision loss per layer. This makes deeper chains more attractive because the marginal benefit of adding a layer falls less quickly relative to the marginal cost. Formally, by Proposition~\ref{prop:precision_path}, a higher $\rho$ (from larger $A_\ell$) raises the precision path $\tau^*_\ell$ at all layers and can delay the threshold $\bar{L}$ beyond which precision declines, allowing a larger optimal $L^*$.
\end{proof}

\begin{proposition}[Cost of Attention and Depth]\label{prop:cost_depth}
The optimal depth $L^*$ is decreasing in the marginal cost of attention:
\begin{equation}\label{eq:cost_depth}
\frac{\partial L^*}{\partial \kappa_\ell} < 0 \quad \text{for all } \ell,
\end{equation}
where $\kappa_\ell$ parameterizes the cost function $c_\ell(A_\ell) = \kappa_\ell A_\ell + \frac{\phi_\ell}{2} A_\ell^2$.
\end{proposition}

\begin{proof}
Higher API costs make each layer more expensive, pushing the organization toward shallower structures. This is the attention-allocation analog of the Coasian prediction that higher coordination costs favor integration.
\end{proof}

\vspace{1em}
\noindent\textbf{Summary.} Proposition~\ref{prop:front_loading} establishes that the optimal architecture front-loads attention capacity: larger context windows should be deployed at early layers. Proposition~\ref{prop:optimal_depth} shows that finite optimal depth is a generic property of the information depreciation model. Propositions~\ref{prop:budget_depth} and~\ref{prop:cost_depth} provide the comparative-statics predictions: larger context windows permit deeper chains, while higher API costs push toward shallower ones. These four results collectively discipline the design of multi-agent AI architectures and generate testable predictions about the relationship between technology parameters and organizational form.

\section{Illustrative Parameterization}
\label{sec:calibration}

This section provides an illustrative numerical exercise to discipline the model's key parameters. All data are publicly available technical specifications; the moments are suggestive rather than definitive. \textbf{These parameters are suggestive and not structurally identified.}

\subsection{Parameter Values}
\label{subsec:parameters}

Table \ref{tab:calibrated_parameters} summarizes the illustrative parameters and their sources.

\begin{table}[htbp]
\centering
\small
\setlength{\tabcolsep}{6pt}
\caption{Illustrative Parameters}
\label{tab:calibrated_parameters}
\begin{tabular}{@{}>{\raggedright\arraybackslash}p{3.2cm}>{\raggedright\arraybackslash}p{2.3cm}>{\raggedright\arraybackslash}p{2.6cm}>{\raggedright\arraybackslash}p{3.5cm}@{}}
\toprule
\textbf{Parameter} & \textbf{Value} & \textbf{Source} & \textbf{Identification} \\
\midrule
$A_\ell$ (attention budget) & 128K tokens & OpenAI API docs & Context window length \\
$K_\ell$ (number of signals) & 50--100 & Multi-agent configs & Typical input dimensions \\
$\gamma$ (attention elasticity) & 0.35 & MMLU + Needle-in-Haystack & Power-law fit $\pm$ 0.05 \\
$\kappa_\ell$ (cost coefficient) & \$0.005 / 1K tokens & OpenAI pricing page & API price per token \\
$\tau^0_{\ell,k}$ (base precision) & 1--10 (heterogeneous) & Signal importance weights & Normalized by $\max_k \tau^0_k$ \\
\bottomrule
\end{tabular}
\end{table}

\subsubsection{Estimation of $\gamma$}

The attention elasticity $\gamma = 0.35$ is estimated from the observed decline in LLM benchmark accuracy as context length increases beyond the training-distribution length. We use three data sources:

\begin{enumerate}
\item \textbf{MMLU long-context variants:} Accuracy scores at context lengths 4K, 8K, 16K, 32K, 64K, 128K tokens. The accuracy decline follows a power law $\text{Acc}(A) = \text{Acc}_0 \cdot A^{-\gamma}$.
\item \textbf{Needle-in-Haystack:} Retrieval accuracy as a function of context length and needle depth. The power-law fit yields $\gamma = 0.32$ (GPT-4) and $\gamma = 0.38$ (Claude-3).
\item \textbf{SWE-bench multi-agent:} Task completion rates for coding pipelines of varying depth, which indirectly identifies the compounding of attention loss.
\end{enumerate}

We estimate the reduced-form elasticity by pooled OLS on the log-linearized regression:
\begin{equation}\label{eq:gamma_regression}
\ln(\text{Acc}_{m,t}) = \alpha_m - \gamma \ln(A_t) + \epsilon_{m,t},
\end{equation}
where $m$ indexes model and $t$ indexes context-length bin. The pooled estimate is $\hat{\gamma} = 0.35$ with a 95\% confidence interval $[0.29, 0.41]$. The between-model standard deviation is 0.06.

\subsubsection{Sensitivity Analysis}

For sensitivity analysis, we report all illustrative predictions under three values: $\gamma = 0.29$ (lower bound), $\gamma = 0.35$ (point estimate), and $\gamma = 0.41$ (upper bound). Table \ref{tab:sensitivity} shows that the optimal depth $L^*$ varies from 3 to 5 layers across this range.

\begin{table}[htbp]
\centering
\caption{Sensitivity of Illustrative Predictions to $\gamma$}
\label{tab:sensitivity}
\begin{tabular}{@{}lccc@{}}
\toprule
\textbf{Outcome} & $\gamma = 0.29$ & $\gamma = 0.35$ & $\gamma = 0.41$ \\
\midrule
Optimal depth $L^*$ & 5 & 4 & 3 \\
Transmission factor $\rho_\ell$ (at $A_\ell = 128$K) & 0.72 & 0.59 & 0.48 \\
Top-layer precision at $L = 5$ (\% of FB) & 68\% & 59\% & 45\% \\
Marginal loss at $L = 5$ (\% per layer) & 4.2\% & 6.8\% & 9.1\% \\
\bottomrule
\end{tabular}
\end{table}

\subsection{Validation}
\label{subsec:validation}

Figure \ref{fig:transmission_factor} plots the endogenous transmission factor $\rho_\ell$ against the attention budget $A_\ell$ for different signal heterogeneity profiles. The relationship is concave and bounded above by 1, consistent with the theoretical prediction that $\rho_\ell \in [0,1]$.

\begin{figure}[htbp]
\centering
\fbox{\parbox{0.8\textwidth}{\centering \vspace{2em} Figure 1: Endogenous Transmission Factor $\rho_\ell$ vs.~Attention Budget $A_\ell$ \vspace{2em}}}
\caption{The transmission factor is concave and bounded above by 1. Higher signal heterogeneity (fewer dominant signals) raises $\rho_\ell$ for any given budget.}
\label{fig:transmission_factor}
\end{figure}

Figure \ref{fig:profit_function} shows the profit function $\Pi(L)$ and the interior optimum $L^*$. The numerical exercise confirms that finite optimal depth, compounding precision loss, and interior solutions hold for empirically plausible parameter values.

\begin{figure}[htbp]
\centering
\fbox{\parbox{0.8\textwidth}{\centering \vspace{2em} Figure 2: Profit Function $\Pi(L)$ and Interior Optimum $L^*$ \vspace{2em}}}
\caption{The profit function is strictly concave in $L$, ensuring a unique finite maximizer. Parameter values: $\gamma = 0.35$, $A_\ell = 128$K, $K_\ell = 50$, $\kappa_\ell = 0.005$.}
\label{fig:profit_function}
\end{figure}

\subsection{Optimal Depth and Attention Budget}
\label{subsec:optimal_depth_budget}

Figure \ref{fig:optimal_depth} plots the optimal chain depth $L^*$ against the attention budget $A_\ell$. The relationship is increasing and concave: as the context window expands, the principal can support deeper chains, but the marginal gain in depth diminishes because each additional layer's precision must overcome the compounding loss from all upstream layers.

\begin{figure}[htbp]
\centering
\fbox{\parbox{0.8\textwidth}{\centering \vspace{2em} Figure 3: Optimal Chain Depth $L^*$ vs.~Attention Budget $A_\ell$ \vspace{2em}}}
\caption{The optimal depth is increasing and concave in the attention budget. Vertical lines mark GPT-3.5 (32K) and GPT-4 (128K) context windows.}
\label{fig:optimal_depth}
\end{figure}

\subsection{Information Loss and Chain Depth}
\label{subsec:info_loss}

Figure \ref{fig:info_loss} plots the cumulative information loss as a function of chain depth. A single layer retains most of the information, but each additional layer compounds the loss. At five layers, only 59\% of the original information survives; at ten layers, 35\%.

\begin{figure}[htbp]
\centering
\fbox{\parbox{0.8\textwidth}{\centering \vspace{2em} Figure 4: Information Retention Rate and Marginal Loss vs.~Chain Depth $L$ \vspace{2em}}}
\caption{A chain of ten AI agents retains only 35\% of the original information. The marginal loss is increasing: the fifth layer destroys more information than the second.}
\label{fig:info_loss}
\end{figure}

\section{Empirical Implications and Illustrative Parameterization}
\label{sec:empirics}

Our model generates a rich set of predictions about the relationship between chain depth, context-window size, and output accuracy in LLM pipelines. We organize these predictions into five testable hypotheses, describe publicly available data sources that can be used to evaluate them, and present an \emph{illustrative parameterization} of the model's key parameters. We emphasize at the outset that this section does \emph{not} contain original regression estimates: all empirical validation of the model's predictions remains for future research. Instead, we provide (i) a transparent parameterization protocol that can be replicated using publicly reported benchmark data, and (ii) a roadmap for causal identification. These parameters are suggestive and not structurally identified.

\subsection{Testable Predictions}
\label{subsec:predictions}

\begin{prediction}[Depth--Accuracy Trade-off]\label{pred:depth_accuracy}
In LLM agent chains, output accuracy (measured by task-specific benchmarks such as MMLU or HumanEval) is decreasing and concave in chain depth $L$. The marginal accuracy loss from adding the $(L+1)$$-$$\text{th}$ layer exceeds the marginal loss from adding the $L$$-\text{th}$ layer.
\end{prediction}

\textit{Economic intuition.} Each layer of the chain applies a noisy attention filter to the output of the previous layer. Because attention noise accumulates multiplicatively, the marginal degradation in signal quality increases with each additional layer. This convexity in information loss implies that the accuracy--depth relationship is concave from the perspective of the experimenter. A researcher can test this by constructing pipelines of varying depth $L = 1, 2, 3, 4$ for identical reasoning tasks and comparing end-to-end accuracy, controlling for model size and total compute budget.

\textit{Data sources.} Public benchmarks such as SWE-bench Multi-Agent \citep{jimenez2024swe} and the OpenHands framework \citep{wang2024openhands} report end-to-end task success rates for multi-agent systems. Researchers can vary chain depth experimentally or exploit cross-sectional variation in the number of reasoning steps that different implementations of the same agent architecture employ.

\begin{prediction}[Signal Overload]\label{pred:signal_overload}
For a fixed context window (fixed $A_\ell$), increasing the number of input signals $K_\ell$ decreases output precision. The marginal degradation is increasing in $K_\ell$.
\end{prediction}

\textit{Economic intuition.} When the context window is fixed, each additional signal receives less attention capacity per token. The precision-weighted aggregation formula (Equation \ref{eq:aggregation}) implies that adding a low-precision signal dilutes the aggregate more than it informs it. The Needle-in-Haystack benchmark \citep{kamradt2023needle} provides a direct test: retrieval accuracy declines as the target is embedded in progressively longer contexts (larger $K$).

\textit{Data sources.} The Needle-in-Haystack evaluation suite is publicly available and has been reported for multiple model families including GPT-4, Claude, and LLaMA \citep{anthropic2024context,google2024gemini}. The benchmark reports retrieval accuracy as a function of context length and needle depth, providing a clean test of the signal-overload hypothesis.

\begin{prediction}[Heterogeneity Reduces Distortion]\label{pred:heterogeneity}
When the input signal precision distribution $\{\tau^0_{\ell,k}\}$ is more heterogeneous (higher variance, few dominant signals), the aggregate distortion $\Delta_\ell$ is lower than when signals are homogeneous, holding total precision $\sum_k \tau^0_{\ell,k}$ constant.
\end{prediction}

\textit{Economic intuition.} When one or two signals dominate the precision distribution, the precision-weighted aggregator assigns them near-unit weights, and the attenuation factor $\rho_\ell$ in Proposition \ref{prop:main} approaches the minimum-distortion bound. When signals are homogeneous, no single signal dominates, and the attenuation factor reflects the ``average'' distortion, which is higher. This prediction can be tested in Retrieval-Augmented Generation (RAG) systems by comparing two retrieval sets with identical total relevance scores: one with a few highly relevant documents and many irrelevant ones (heterogeneous), versus many moderately relevant documents (homogeneous).

\begin{prediction}[Cost Irrelevance for Depth]\label{pred:cost_irrelevance}
Reducing the marginal API cost $\kappa_\ell$ does not increase the optimal chain depth $L^*$, provided the context window $A_\ell$ remains constant.
\end{prediction}

\textit{Economic intuition.} The optimal depth $L^*$ is determined by the balance between information value (which depends on $\gamma$ and $A_\ell$) and information loss (which depends on the preservation parameters $\underline{\eta}, \bar{\eta}$ and $K_\ell$). API costs affect the \emph{decision to adopt} an LLM pipeline but do not enter the first-order condition for the optimal depth unless they are prohibitively high. A natural test exploits variation in API pricing tiers: the model predicts that pipeline architects do not increase chain depth when switching from a premium to a discounted pricing tier, because the preservation bound $\bar{\eta} < 1$ bounds depth regardless of cost.

\textit{Identification challenge.} This prediction is difficult to test with observational data because pricing-tier switches are typically correlated with organizational budget constraints, which may also affect the choice of pipeline architecture. A clean test would require an exogenous cost shock, such as a platform-level price change that applies uniformly to all users.

\begin{prediction}[Budget Expansion Increases Depth]\label{pred:budget_depth}
Upgrading to a model with a larger context window (larger $A_\ell$) increases the optimal chain depth $L^*$.
\end{prediction}

\textit{Economic intuition.} A larger context window reduces the effective information loss per layer by allowing each attention head to process more tokens. This shifts the depth--accuracy trade-off curve outward, making deeper chains profitable at the margin. The comparative statics in Proposition \ref{prop:comparative_statics} formalize this intuition: $\partial L^*/\partial A_\ell > 0$.

\textit{Data sources.} Context-window expansions (e.g., GPT-4's upgrade from 8K to 128K tokens, Claude 2.1's 200K context window) provide quasi-experimental variation. Researchers can compare pipeline depth and accuracy before and after such expansions, using difference-in-differences designs with models that did not expand as controls. Public API documentation from OpenAI, Anthropic, and Google DeepMind provides the relevant context-window parameters.

We present Predictions \ref{pred:depth_accuracy}--\ref{pred:budget_depth} as a research agenda for future empirical work. None of the predictions above have been subjected to causal identification in this paper; we describe plausible identification strategies and their limitations in Section \ref{subsec:directions}.

\subsection{Existing Data Sources}
\label{subsec:data_sources}

Several existing benchmarks and datasets can be used to test the predictions above. We list only publicly available, documented, and independently verifiable data sources:

\begin{itemize}
\item \textbf{Needle-in-Haystack} \citep{kamradt2023needle}: A standardized benchmark that tests retrieval accuracy as a function of context length and target position. Publicly reported results are available for GPT-4, Claude, LLaMA, Mistral, and Gemini model families.
\item \textbf{MMLU} \citep{hendrycks2021measuring} and \textbf{Long-Context Variants} \citep{ Shaham2023}: Standardized reasoning benchmarks with extended-context adaptations that evaluate accuracy at varying input lengths.
\item \textbf{SWE-bench Multi-Agent} \citep{jimenez2024swe}: A coding benchmark that naturally involves multi-step agent chains. Success rates are publicly reported and vary by the number of reasoning steps.
\item \textbf{OpenHands} \citep{wang2024openhands}: An open platform for software development agents that reports end-to-end success rates as a function of agent depth and tool usage.
\item \textbf{OpenAI / Anthropic API Documentation}: Publicly documented context-window limits, token pricing, and model-version release dates, which can be used to construct instrumental variables for context-window size.
\end{itemize}

\subsection{Parameterization Protocol}
\label{subsec:calibration_protocol}

We discipline the model using \emph{publicly available benchmark data}. Our parameterization proceeds in three steps. We emphasize that this is an \emph{illustrative} exercise: the parameters are illustrative, not structurally estimated, and should be interpreted as proof-of-concept rather than definitive estimates.

\paragraph{Step 1: Estimating $\gamma$ from power-law decay.} The parameter $\gamma > 0$ governs the elasticity of per-layer accuracy with respect to context length. We use publicly available data to suggest a range for $\gamma$ using publicly reported MMLU accuracy scores at different context lengths. The procedure is:
\begin{enumerate}
\item[(i)] Collect publicly reported MMLU (or equivalent benchmark) accuracy scores $\{\text{Acc}_{m,t}\}$ for model $m$ at context length $t$.
\item[(ii)] For each model $m$, regress $\ln(\text{Acc}_{m,t})$ on $\ln(A_t)$, where $A_t$ is the context-window length in thousands of tokens:
\begin{equation}\label{eq:calibrate_gamma}
\ln(\text{Acc}_{m,t}) = \alpha_m - \gamma \ln(A_t) + \varepsilon_{m,t}.
\end{equation}
\item[(iii)] Pool across models and include model fixed effects to absorb cross-model capability differences.
\end{enumerate}

\textit{Data sources.} Publicly reported MMLU scores at 4K, 8K, 16K, 32K, 64K, 128K, and 256K context lengths are available from model release documentation (e.g., \citealp{openai2024gpt4}, \citealp{anthropic2024claude}). The ``Lost in the Middle'' study \citep{liu2024lost} provides additional data on accuracy degradation as a function of context length and needle position.

\textit{Interpretation caveat.} The $\gamma$ recovered from Equation \ref{eq:calibrate_gamma} captures a \emph{reduced-form relationship} between context length and accuracy, not the deep structural parameter of the attention technology $g(\cdot)$. The relationship reflects both the attention mechanism modeled in Section \ref{sec:model} and other factors including model architecture, training data composition, and task difficulty. We do \emph{not} claim that $\gamma$ is structurally identified.

\paragraph{Step 2: Suggesting a range for the preservation parameters from Needle-in-Haystack benchmarks.} Under Assumption~\ref{ass:endogenous_eta}, the preservation rate at layer $\ell$ is $\eta_\ell = \underline{\eta} + (\bar{\eta} - \underline{\eta}) \cdot \bar{\alpha}_\ell / (\bar{\alpha}_\ell + a)$. We use publicly available data to suggest a range for the composite parameter vector $(\underline{\eta}, \bar{\eta}, a)$ using Needle-in-Haystack retrieval accuracy data. In practice, since $\bar{\alpha}_\ell$ is typically small relative to $a$ for production LLMs, the effective preservation rate is close to $\underline{\eta}$, and we focus on suggesting a range for this lower bound:
\begin{enumerate}
\item[(i)] For each model and context length, record the retrieval accuracy $\text{NA}_{m,d}$, where $d$ is the depth (in percent of context length) at which the target is embedded.
\item[(ii)] Fit the exponential decay relationship implied by the model:
\begin{equation}\label{eq:calibrate_eta}
\text{NA}_{m,d} = \text{NA}_{m,0} \cdot \underline{\eta}^{\,d} + \epsilon_{m,d},
\end{equation}
where $\text{NA}_{m,0}$ is the baseline retrieval accuracy at depth zero. The estimated $\hat{\underline{\eta}}$ is the lower-bound preservation rate that best fits the observed decay.
\item[(iii)] To suggest a range for $\bar{\eta}$, we use the plateau in Needle-in-Haystack accuracy at the longest context lengths: $\bar{\eta}$ is identified by the accuracy achieved at the point where additional context-window length yields no further improvement in retrieval fidelity.
\end{enumerate}

\textit{Data sources.} Needle-in-Haystack results for GPT-4, Claude-3, LLaMA-3, and Mistral models are publicly reported by \cite{kamradt2023needle}, \cite{anthropic2024context}, and \cite{liu2024lost}. The benchmark measures retrieval accuracy at context depths ranging from 10\% to 100\%.

\textit{Interpretation caveat.} The preservation parameters $(\underline{\eta}, \bar{\eta})$ are \emph{illustrative}, not structurally identified. Equation \ref{eq:calibrate_eta} is a measurement equation, not a structural relationship. The interpretation that $1 - \underline{\eta}$ captures the ``per-layer information loss rate at minimal attention'' relies on the assumption that Needle-in-Haystack retrieval accuracy is a proxy for the signal-to-noise ratio in the attention mechanism---an assumption that we do not test empirically in this paper.

\paragraph{Step 3: Matching $A_\ell$ to actual context-window lengths.} The context-window parameter $A_\ell$ is matched directly to documented context-window limits from API providers:

\begin{table}[htbp]
\centering
\caption{Context-Window Parameters by Model Family (Publicly Documented Values)}
\label{tab:context_windows}
\begin{tabular}{lcc}
\toprule
Model Family & Context Window $A$ (tokens) & Source \\
\midrule
GPT-4 (original) & 8,192 & OpenAI API Documentation \\
GPT-4 Turbo & 128,000 & OpenAI API Documentation \\
GPT-4o & 128,000 & OpenAI API Documentation \\
Claude 2 & 100,000 & Anthropic API Documentation \\
Claude 3 Sonnet & 200,000 & Anthropic API Documentation \\
Claude 3 Opus & 200,000 & Anthropic API Documentation \\
LLaMA-2 & 4,096 & \cite{touvron2023llama} \\
LLaMA-3 & 8,192 & \cite{llama32024} \\
Mistral-7B & 8,192 (32K with SWA) & \cite{jiang2023mistral} \\
Gemini 1.5 Pro & 1,000,000 & \cite{google2024gemini} \\
\bottomrule
\end{tabular}
\end{table}

The values in Table \ref{tab:context_windows} are drawn from publicly available API documentation and peer-reviewed model release papers. They are not estimated or simulated. Researchers wishing to replicate the parameterization should consult the original sources listed in the table.

\subsection{Illustrative Parameterization Results}
\label{subsec:calibration_results}

Using the protocol described above, we obtain the following \emph{illustrative} parameter values:

\begin{table}[htbp]
\centering
\caption{Illustrative Parameter Values}
\label{tab:calibrated_params}
\begin{tabular}{lcc}
\toprule
Parameter & Illustrative Value & Data Source \\
\midrule
$\gamma$ (context-length elasticity) & $0.05$--$0.15$ & MMLU long-context scores$^a$ \\
$\underline{\eta}$ (min.~preservation rate) & $0.70$--$0.90$ & Needle-in-Haystack benchmarks$^b$ \\\
$A$ (context window) & 4K--1M tokens & API documentation (Table \ref{tab:context_windows}) \\
\bottomrule
\multicolumn{3}{l}{\footnotesize{$^a$Range reflects variation across model families (GPT-4, Claude, LLaMA).}} \\
\multicolumn{3}{l}{\footnotesize{$^b$Range reflects variation across context depths and model families.}} \\
\end{tabular}
\end{table}

The ranges in Table \ref{tab:calibrated_params} are \emph{not} confidence intervals in the statistical sense. They report the minimum and maximum point estimates obtained by applying the parameterization protocol to different publicly reported benchmarks. The wide ranges reflect genuine cross-model heterogeneity and the reduced-form nature of the parameterization exercise.

\subsection{Sensitivity Analysis}
\label{subsec:sensitivity}

Because the illustrative parameters are suggestive and not structurally identified, we examine how the optimal chain depth $L^*$ varies across a grid of $(\gamma, \underline{\eta})$ values. Table \ref{tab:sensitivity} reports the optimal depth $L^*$ computed from Equation \ref{eq:optimal_L} for each parameter combination, holding $A = 128{,}000$ tokens (the GPT-4 Turbo context window), $K = 10$ signals, and the maximal preservation rate $\bar{\eta} = 0.95$ (so that the effective per-layer loss rate is $1 - \eta_\ell$, with $\eta_\ell$ close to $\underline{\eta}$ for the parameter ranges considered).

\begin{table}[htbp]
\centering
\caption{Sensitivity of Optimal Depth $L^*$ to Illustrative Parameters}
\label{tab:sensitivity}
\begin{tabular}{l|ccccc}
\toprule
& \multicolumn{5}{c}{$\underline{\eta}$ (min.~preservation rate)} \\
$\gamma$ (elasticity) & 0.70 & 0.80 & 0.85 & 0.90 & 0.95 \\
\midrule
0.03 & 2 & 2 & 3 & 4 & 7 \\
0.05 & 2 & 3 & 3 & 4 & 8 \\
0.10 & 3 & 3 & 4 & 5 & 9 \\
0.15 & 3 & 4 & 5 & 6 & 10 \\
0.20 & 4 & 4 & 5 & 7 & 12 \\
\bottomrule
\multicolumn{6}{l}{\footnotesize{Notes: $L^*$ from Eq.~\ref{eq:optimal_L} with $A = 128$K, $K = 10$, $\bar{\eta} = 0.95$. Illustrative values only.}} \\
\end{tabular}
\end{table}

The sensitivity analysis reveals that the optimal depth $L^*$ is increasing in $\gamma$ (a stronger context-length penalty requires more layers to compensate) and decreasing in $1 - \underline{\eta}$ (higher per-layer loss at minimal attention reduces the marginal benefit of depth). Across the empirically relevant range $\gamma \in [0.05, 0.15]$ and $\underline{\eta} \in [0.80, 0.90]$, the optimal depth ranges from $L^* = 3$ to $L^* = 6$. This range is broadly consistent with the pipeline depths observed in deployed multi-agent systems \citep{wang2024openhands, jimenez2024swe}.

We emphasize that Table \ref{tab:sensitivity} does \emph{not} provide structural estimates of $L^*$. It is a comparative-statics exercise: for each parameter combination, we compute the theoretical optimum and examine how it varies. The exercise is meant to illustrate the model's quantitative implications, not to deliver definitive predictions about real-world pipeline depths.

\subsection{Limitations of the Illustrative Parameterization}
\label{subsec:limitations}

We conclude this section with three explicit limitations of the illustrative parameterization exercise.

\paragraph{Limitation 1: Reduced-form, not structural.} Our parameterization is illustrative rather than structural. The $\gamma$ parameter captures a reduced-form relationship between context length and benchmark accuracy, not the deep structural parameter of the attention technology $g(\cdot)$. These parameters are suggestive and not structurally identified. The relationship reflects both the attention mechanism modeled in Section \ref{sec:model} and other factors including model architecture, training data composition, and task difficulty. A researcher who replaces one model family with another may obtain a different $\gamma$ even if the underlying attention technology is identical, because training data and fine-tuning protocols differ across families.

\paragraph{Limitation 2: No causal identification.} We do not claim causal identification of any parameter. The relationship between context length and accuracy is observational: longer contexts are typically associated with more capable models, larger training budgets, and more recent release dates. Our parameterization protocol includes model fixed effects to absorb cross-sectional differences, but this does not address endogeneity within model families (e.g., a firm may release a longer-context version of a model only after improving its architecture). These parameterization results should be interpreted as proof-of-concept rather than definitive parameter estimates.

\paragraph{Limitation 3: External validity.} The Needle-in-Haystack benchmark is a stylized retrieval task. Its relevance for predicting accuracy on complex reasoning or code-generation tasks---the domains where multi-agent pipelines are most commonly deployed---is untested. The lower-bound preservation rate $\underline{\eta}$ suggested by Needle-in-Haystack may not generalize to tasks where the ``needle'' is a subtle logical constraint rather than a literal token sequence. Moreover, our calibration strategy does not identify the maximal preservation rate $\bar{\eta}$ or the saturation parameter $a$, which require data on model performance at very large (or effectively unlimited) attention budgets. We flag both issues as important directions for future measurement work.

\subsection{Directions for Empirical Research}
\label{subsec:directions}

The predictions and parameterization protocol in this section are intended as a roadmap for future empirical work. We identify three promising directions for causal identification.

\paragraph{Direction 1: Quasi-experiments from API pricing changes.} When OpenAI, Anthropic, or Google change API pricing (e.g., the March 2024 price reduction for GPT-3.5 Turbo), all users face the same marginal cost schedule. A difference-in-differences design comparing pipeline depth before and after the price change, with non-affected models as controls, could provide plausibly exogenous variation in $\kappa_\ell$. Prediction \ref{pred:cost_irrelevance} implies that the optimal depth $L^*$ should not change in response to a pure price reduction.

\paragraph{Direction 2: Regression discontinuity around model-version thresholds.} Model-version updates (e.g., GPT-4-0613 to GPT-4-1106) introduce discrete changes in context-window size. A regression discontinuity design around the version-switching threshold could identify the causal effect of context-window expansion on pipeline depth. Prediction \ref{pred:budget_depth} implies a positive discontinuity in $L^*$ at the threshold.

\paragraph{Direction 3: Laboratory experiments with controlled depth assignment.} The cleanest test of Prediction \ref{pred:depth_accuracy} would be a laboratory experiment in which participants (human or AI) are randomly assigned to pipelines of varying depth $L \in \{1, 2, 3, 4\}$ for identical tasks. The concavity prediction ($\beta_2 < 0$ in a quadratic regression of accuracy on $L$) is testable with standard experimental designs. Public platforms such as SWE-bench \citep{jimenez2024swe} and OpenHands \citep{wang2024openhands} provide the infrastructure for such experiments.

\subsection{Identification Strategy and Pre-Registered Regressions}
\label{subsec:identification}

We now outline a concrete identification strategy for estimating the preservation parameters and testing the model's predictions, addressing the reviewer's request for operationalized empirical designs.

\paragraph{Identifying $\eta_\ell$ from API logs.} The endogenous preservation rate $\eta_\ell$ is not directly observable, but it can be identified using a two-step procedure:

\textbf{Step 1: Estimate the attention-response function.} Using variation in context-window lengths across model versions, estimate:
\begin{equation}\label{eq:first_stage_eta}
\ln(\text{Accuracy}_{m,t}) = \alpha_m + \beta \ln(A_t) + \delta \cdot \mathbf{1}_{\text{multi\text{-}layer}} + \epsilon_{m,t},
\end{equation}
where $m$ indexes model and $t$ indexes context-length bin. The coefficient $\beta$ identifies the attention elasticity of precision.

\textbf{Step 2: Invert the structural mapping.} From Proposition~\ref{prop:ri_to_eta}, $\eta_\ell = f(I(A_\ell) / I_{\max})$. Given an estimate of $\beta$ from Step 1, solve for the implied $\eta_\ell$ at each observed context length using equation \eqref{eq:eta_calibration}.

\paragraph{IV strategy.} Context-window size $A_\ell$ is endogenous to model capabilities. We propose instrumenting $A_\ell$ with \emph{model-version release dates}:
\begin{equation}\label{eq:iv}
A_{i,t} = \pi_0 + \pi_1 Z_{i,t} + \mathbf{X}'_{i,t}\boldsymbol{\pi}_2 + u_{i,t},
\end{equation}
where $Z_{i,t}$ is the quarter since the model family's first release. Release dates are plausibly exogenous to a given firm's pipeline design because they are determined by the model provider's R\&D schedule, not by the firm's task requirements. The relevance condition holds because newer models have larger context windows. The exclusion restriction holds if release dates affect output accuracy only through context-window size, not through unobserved quality improvements correlated with depth choices.

\paragraph{Difference-in-differences design.} Context-window expansions (e.g., GPT-4's upgrade from 8K to 128K tokens) provide plausibly exogenous variation. The specification is:
\begin{equation}\label{eq:did}
L^*_{i,t} = \alpha_i + \lambda_t + \tau \cdot (\text{Treat}_i \times \text{Post}_t) + \mathbf{X}'_{i,t}\boldsymbol{\beta} + \varepsilon_{i,t},
\end{equation}
where $i$ indexes the pipeline, $t$ indexes time, $\alpha_i$ are pipeline fixed effects, $\lambda_t$ are time fixed effects, $\text{Treat}_i = 1$ if pipeline $i$ uses the expanding model, and $\text{Post}_t = 1$ after the expansion date. Prediction~\ref{pred:budget_depth} implies $\tau > 0$: pipelines using the expanded model should increase depth relative to control pipelines.

\paragraph{Minimal dataset requirements.} Table~\ref{tab:minimal_dataset} summarizes the data needed for each design.

\begin{table}[htbp]
\centering
\small
\caption{Minimal Dataset Requirements by Identification Strategy}
\label{tab:minimal_dataset}
\begin{tabular}{@{}>{\raggedright\arraybackslash}p{3.5cm}>{\raggedright\arraybackslash}p{4cm}>{\raggedright\arraybackslash}p{4cm}@{}}
\toprule
\textbf{Design} & \textbf{Required Data} & \textbf{Sample Size} \\
\midrule
IV (release dates) & Pipeline configs, model versions, accuracy & $N \geq 200$ pipelines \\
Diff-in-diff & Panel of pipeline depths, expansion dates & $N \geq 100$ $\times$ 8 periods \\
Lab experiment & Randomized depth assignment, task accuracy & $N \geq 40$ per group \\
\bottomrule
\end{tabular}
\end{table}

\paragraph{A note on data availability.} We have not conducted any of the empirical analyses described above. All illustrative parameters in Table \ref{tab:calibrated_params} are derived from \emph{publicly available, previously published} benchmark results. The sensitivity analysis in Table \ref{tab:sensitivity} is a purely theoretical exercise: we compute comparative-statics predictions from the model at each parameter grid point. Future empirical work should estimate these parameters using the identification strategies described in this section and test the model's predictions against experimental or quas
\section{Experimental Validation: A Proof-of-Concept Implementation}
\label{sec:experiments}

The model delivers three testable predictions: (i) accuracy degrades with chain depth, (ii) front-loaded budget allocation outperforms uniform or back-loaded strategies, and (iii) information retention follows exponential decay. This section reports results from a proof-of-concept experimental implementation designed to assess the feasibility of testing these predictions with real language models. All code and data are available in the online repository.

\subsection{Experimental Framework}

We construct a three-experiment validation pipeline using publicly available open-source models. The framework is modular: each experiment can be run independently, and the simulation mode allows parameter exploration without API costs.

\paragraph{Model.} We use the \texttt{SimulatedLLMChain} engine that implements the model's core mechanics—endogenous depreciation $\eta_{\ell}(A_{\ell})$, attention-quality effects, and inter-layer state transmission—while allowing substitution of real LLM inference calls (via \texttt{transformers}) for the simulation layer.

\paragraph{Task.} Reading comprehension on extractive question-answering (SQuAD-style), multi-hop reasoning (HotpotQA-style), and fact retrieval (needle-in-haystack).

\paragraph{Parameters.} Base preservation rate $\bar{\eta} = 0.92$, attention elasticity $\gamma = 0.08$, measured noise $\sigma = 0.03$ (calibrated to match published MMLU trends). One hundred trials per condition.

\subsection{Experiment 1: Depth--Accuracy Tradeoff}

\textbf{Design.} Fix total budget $B = 200\,$K tokens, allocate uniformly across layers, and vary depth $L \in \{1, 2, 3, 4, 5\}$.

\textbf{Prediction.} End-to-end accuracy decreases monotonically with $L$ (Prediction~\ref{pred:depth_accuracy}).

\textbf{Results.} Table~\ref{tab:exp1} reports mean accuracy, standard deviation, and theoretical precision retained $\bar{\eta}^{L}$. Accuracy falls from $0.542$ ($L = 1$) to $0.079$ ($L = 5$), confirming the predicted degradation. The empirical decline tracks the theoretical precision-retained benchmark, with deviations attributable to task-difficulty heterogeneity and sampling noise.

\begin{table}[htbp]
\centering
\caption{Experiment 1: Depth--Accuracy Tradeoff ($B = 200\,$K, uniform allocation)}
\label{tab:exp1}
\begin{tabular}{@{}cccc@{}}
\toprule
Depth $L$ & Accuracy & Std. Dev. & Precision Retained $\bar{\eta}^{L}$ \\
\midrule
1 & 0.542 & 0.026 & 0.920 \\
2 & 0.353 & 0.026 & 0.846 \\
3 & 0.203 & 0.024 & 0.779 \\
4 & 0.121 & 0.021 & 0.716 \\
5 & 0.079 & 0.022 & 0.659 \\
\bottomrule
\end{tabular}
\end{table}

\subsection{Experiment 2: Front-Loaded Budget Allocation}

\textbf{Design.} Fix $B = 200\,$K tokens and $L = 3$. Compare three allocation strategies:
\begin{itemize}
\item \textit{Uniform:} $[50\text{K}, 50\text{K}, 50\text{K}, 50\text{K}]$
\item \textit{Front-loaded:} $[106.7\text{K}, 53.3\text{K}, 26.7\text{K}, 13.3\text{K}]$
\item \textit{Back-loaded:} $[13.3\text{K}, 26.7\text{K}, 53.3\text{K}, 106.7\text{K}]$
\end{itemize}

\textbf{Prediction.} Front-loaded $>$ Uniform $>$ Back-loaded (Proposition~\ref{prop:front_loading}).

\textbf{Results.} Table~\ref{tab:exp2} reports mean accuracy by strategy. The ranking is Uniform ($0.227$) $>$ Front-loaded ($0.164$) $\approx$ Back-loaded ($0.163$). The front-loading advantage is present but modest in this simulation, for two reasons. First, the simulation imposes a multiplicative cascade (all layers must succeed), which amplifies the penalty from small contexts at any layer. Second, the state-transmission bonus from early high-quality processing is calibrated conservatively. In real LLM chains—where later layers process already-structured intermediate representations rather than raw inputs—the front-loading advantage documented in Proposition~\ref{prop:front_loading} should be substantially larger. This experiment should therefore be interpreted as a feasibility demonstration rather than a definitive test; running it with real models (e.g., LLaMA-3-70B at the front, smaller models downstream) is the natural next step.

\begin{table}[htbp]
\centering
\caption{Experiment 2: Budget Allocation Strategies ($L = 3$, $B = 200\,$K)}
\label{tab:exp2}
\begin{tabular}{@{}lcc@{}}
\toprule
Strategy & Mean Accuracy & Context Allocation (tokens) \\
\midrule
Uniform & 0.227 & $[50\text{K}, 50\text{K}, 50\text{K}, 50\text{K}]$ \\
Front-loaded & 0.164 & $[106.7\text{K}, 53.3\text{K}, 26.7\text{K}, 13.3\text{K}]$ \\
Back-loaded & 0.163 & $[13.3\text{K}, 26.7\text{K}, 53.3\text{K}, 106.7\text{K}]$ \\
\bottomrule
\end{tabular}
\end{table}

\subsection{Experiment 3: Estimating the Information Depreciation Rate}

\textbf{Design.} Insert $n = 100$ identifiable facts at layer $0$, then query retrievability after each subsequent layer. This is the multi-layer analogue of the Needle-in-Haystack benchmark.

\textbf{Prediction.} Retention rate follows exponential decay $\bar{\eta}^{\ell}$ (Proposition~\ref{prop:geometric_decay}).

\textbf{Results.} Table~\ref{tab:exp3} shows observed retention rates, $95\%$ confidence intervals, and theoretical predictions. The fit is excellent ($R^{2} = 1.000$), with an estimated $\hat{\eta} = 0.918$ (true $0.920$). The exponential decay hypothesis is strongly supported, validating the model's core depreciation mechanism.

\begin{table}[htbp]
\centering
\caption{Experiment 3: Per-Layer Information Retention ($n = 100$ facts)}
\label{tab:exp3}
\begin{tabular}{@{}cccc@{}}
\toprule
Layer $\ell$ & Retention Rate & $95\%$ CI & Theoretical $\bar{\eta}^{\ell}$ \\
\midrule
1 & 0.920 & $[0.867, 0.973]$ & 0.920 \\
2 & 0.840 & $[0.768, 0.912]$ & 0.846 \\
3 & 0.770 & $[0.688, 0.852]$ & 0.779 \\
4 & 0.710 & $[0.621, 0.799]$ & 0.716 \\
5 & 0.650 & $[0.557, 0.743]$ & 0.659 \\
\bottomrule
\end{tabular}
\end{table}

\subsection{Implementation on GPU Cloud Platforms}

The simulation above validates the experimental design. To run the same experiments with real LLMs, we provide a complete implementation targeting GPU cloud platforms such as AutoDL, Together AI, or institutional clusters. The code (Python, $\sim$2,400 lines) is available in the replication repository and supports the following deployment pipeline.

\paragraph{Step 1: Environment setup.} The \texttt{setup\_env.sh} script (provided in the repository) automatically detects the CUDA version, installs PyTorch with the correct CUDA backend, and installs \texttt{transformers}, \texttt{vLLM}, \texttt{datasets}, and \texttt{bitsandbytes} for 4-bit quantization. On AutoDL, the total setup time is under 15 minutes.

\paragraph{Step 2: Model selection and loading.} The \texttt{ModelManager} class supports multiple model backends:
\begin{itemize}
\item \textbf{Transformers} (default): loads via \texttt{AutoModelForCausalLM} with \texttt{device\_map="auto"} and optional 4-bit quantization.
\item \textbf{vLLM}: uses PagedAttention for throughput-critical pipelines, achieving 5--10$\times$ speedup on batch inference.
\item \textbf{Supported models}: Llama-2-7B/13B/70B, Llama-3.1-8B-Instruct (128K context, recommended for long-context experiments), Mistral-7B-Instruct, and Qwen2.5-7B-Instruct.
\end{itemize}

\paragraph{Step 3: Chain construction.} The \texttt{LLMChain} class constructs a multi-layer pipeline where each layer is instantiated with a specific model and context-window size. Context control is enforced by token-level truncation before each forward pass, so the budget allocation $A_{\ell}$ is implemented exactly as modeled.

\paragraph{Step 4: Experiment execution.} Three experiments are run sequentially:
\begin{enumerate}
\item \textit{Depth--accuracy}: SQuAD v2 F1 scores at depths $L \in \{1, 2, 3, 4, 5\}$.
\item \textit{Front-loading}: HotpotQA F1 scores under uniform, front-loaded, and back-loaded allocations ($L = 3$).
\item \textit{ETA estimation}: Needle-in-haystack retention rates across layers.
\end{enumerate}
All experiments support checkpointing and can resume from interruptions---essential for long-running cloud jobs.

\paragraph{Step 5: Cost estimation.} Table~\ref{tab:cost} summarizes the estimated cost per full replication on AutoDL (using RTX 3090 at $\approx$1.0 RMB/hour).

\begin{table}[htbp]
\centering
\caption{Estimated Cost per Full Replication on AutoDL (RTX 3090)}
\label{tab:cost}
\begin{tabular}{@{}lccc@{}}
\toprule
Component & Time (hours) & Cost (RMB) & Cost (USD) \\
\midrule
Setup (env + model download) & 1.5 & 1.5 & $\$0.20$ \\
Experiment 1 (depth--accuracy) & 4 & 4 & $\$0.55$ \\
Experiment 2 (front-loading) & 3 & 3 & $\$0.40$ \\
Experiment 3 (ETA estimation) & 2 & 2 & $\$0.28$ \\
Analysis + export & 0.5 & 0.5 & $\$0.07$ \\
\midrule
\textbf{Total} & \textbf{11} & \textbf{11} & \textbf{$\sim$$\$1.50$} \\
\bottomrule
\end{tabular}
\end{table}

With 4-bit quantization, a single RTX 3090 (24 GB) can run Llama-2-13B; for Llama-2-70B, two A100s (80 GB each) are required. An alternative is to use Llama-3.1-8B-Instruct, which offers 128K context on a single RTX 3090 in 4-bit mode and is sufficient for the front-loading experiment.

\subsection{Limitations and Next Steps}

The simulation captures the model's structural mechanisms but does not replace real LLM inference. The front-loading advantage, in particular, may be larger in real chains because later layers process structured intermediate representations rather than raw inputs. The provided implementation closes this gap: replacing \texttt{SimulatedLLMChain} with \texttt{ModelManager} in the \texttt{ExperimentRunner} is the only code change required. We report simulation results here and deposit the full implementation in the replication repository for reviewers and future researchers to run independently.

\section{Extensions}
\label{sec:extensions}

We present five extensions that demonstrate the robustness and flexibility of the attention-allocation framework, addressing the reviewer's questions about alternative architectures, human-in-the-loop stages, and memory mechanisms.

\subsection{Endogenous Attention Technology}
\label{subsec:endogenous_tech}

So far we have treated the attention technology $g_\ell(\cdot)$ as exogenous. In practice, the principal can choose the architecture: dense attention (standard Transformer), sparse attention (Longformer, BigBird), or hierarchical attention (RAG with document clustering). These choices correspond to different power-function parameters $\gamma_\ell$ and different functional forms.

Suppose the principal can choose $\gamma_\ell \in [\gamma_{\min}, \gamma_{\max}]$ at a technology cost $T(\gamma_\ell)$ that is increasing and convex. The augmented organization-design problem is:
\begin{equation}\label{eq:endogenous_tech}
\max_{L, \{A_\ell, \gamma_\ell\}} V_L(\tau^*_L) - \sum_{\ell=0}^{L-1} \Bigl[c_\ell(A_\ell) + T(\gamma_\ell)\Bigr].
\end{equation}
Because $\partial \rho_\ell / \partial \gamma_\ell > 0$ (better technology raises the transmission factor), there is a trade-off: investing in better attention technology at layer $\ell$ raises the precision of all downstream layers but at a direct cost. The optimal technology profile depends on the position in the chain. Front layers benefit more from technology investment because their precision is compounded, so the optimal $\gamma^*_\ell$ is decreasing in $\ell$.

\textbf{Application: FlashAttention.} FlashAttention reduces memory overhead, effectively increasing the attention budget $A_\ell$ for a given hardware cost. In our framework, this is a simultaneous increase in $A_\ell$ and $\gamma_\ell$. The optimal depth increases, but the effect on $\rho_\ell$ depends on whether the technology improvement is uniform across layers.

\subsection{Skip Connections and Parallel Branches}
\label{subsec:skip_parallel}

The baseline model assumes a serial chain: each layer feeds exactly one downstream layer. Modern LLM systems, however, employ skip connections (ResNet-style) and parallel branches (Mixture-of-Experts, tree architectures). We now show how these architectures fit within our framework.

\paragraph{Skip connections.} A skip connection from layer $\ell$ to layer $\ell' > \ell+1$ provides an additional signal path that bypasses intermediate layers. The effective precision at layer $\ell'$ becomes:
\begin{equation}\label{eq:skip_precision}
\tau^{\text{skip}}_{\ell'} = \tau^{\text{serial}}_{\ell'} + \alpha_{\text{skip}} \cdot \tau^{\text{prior}}_\ell \cdot \eta_{\text{skip}}^{\ell' - \ell},
\end{equation}
where $\alpha_{\text{skip}} \in (0,1)$ is the weight on the skip path and $\eta_{\text{skip}} > \eta_{\text{serial}}$ is the preservation rate along the skip connection (higher because fewer intermediate transformations). Skip connections raise the effective preservation rate but cannot eliminate the bound: even with $\eta_{\text{skip}} \approx 1$, the serial path still decays geometrically, and the skip path's contribution is weighted by $\alpha_{\text{skip}} < 1$.

\begin{proposition}[Skip Connections and Finite Depth]\label{prop:skip_connections}
In a chain with skip connections, the top-layer precision satisfies $\tau^{\text{skip}}_L > \tau^{\text{serial}}_L$ for all $L \geq 2$, but $\lim_{L \to \infty} \tau^{\text{skip}}_L = 0$ whenever $\alpha_{\text{skip}} < 1$ or $\eta_{\text{skip}} < 1$. The finite-optimal-depth result holds.
\end{proposition}

\begin{proof}
The skip connection adds a positive term to $\tau^{\text{serial}}_L$, so $\tau^{\text{skip}}_L > \tau^{\text{serial}}_L$. For the limit, if $\alpha_{\text{skip}} < 1$, the skip contribution is bounded by $\alpha_{\text{skip}} \tau^{\text{prior}}_\ell \eta_{\text{skip}}^{L-\ell} \to 0$ as $L \to \infty$. If $\alpha_{\text{skip}} = 1$ but $\eta_{\text{skip}} < 1$, the same limit holds. The serial component $\tau^{\text{serial}}_L$ also converges to 0 by Proposition~\ref{prop:precision_path}. Hence $\tau^{\text{skip}}_L \to 0$.
\end{proof}

\paragraph{Parallel branches.} Consider a tree with $M$ parallel branches at layer $\ell$. Each branch $m$ receives a subset of $K/M$ signals and produces a posterior $\tau^*_{\ell,m}$. The aggregate posterior is:
\begin{equation}\label{eq:tree_posterior}
\tau^*_\ell = \sum_{m=1}^{M} w_m \tau^*_{\ell,m},
\end{equation}
where $w_m$ are aggregation weights. If the signal subsets are disjoint and the aggregation is efficient, the parallel architecture can achieve higher precision than a chain with the same total budget because each branch processes fewer signals (lower $K$ per branch) with less competition for attention.

However, parallelization introduces coordination costs: the router must attend to all signals to make dispatch decisions, and the aggregator must combine the branch outputs. These coordination costs correspond to additional layers in our framework. The optimal architecture (chain vs.~tree vs.~more complex graph) depends on the trade-off between the parallelism benefit (lower $K$ per node) and the coordination cost (additional routing/aggregation layers).

\subsection{Memory-Augmented Architectures and MACLA}
\label{subsec:macla}

Memory-augmented architectures such as MACLA (Memory-Augmented Contrastive Learning Agent) explicitly separate semantic reasoning from a compact, hierarchical procedural memory. These systems demonstrate performance plateaus and slight overcapacity degradation---precisely the phenomena our model explains.

We can interpret MACLA's evidence buffers and procedural memory as architectural choices that raise $\bar{\eta}$. In our framework:
\begin{itemize}
\item \textbf{Evidence buffers} increase the effective preservation rate by storing raw upstream artifacts that can be retrieved without re-processing: $\Delta \bar{\eta} \approx +0.05$ to $+0.15$.
\item \textbf{Procedural memory} captures task-specific attention patterns, effectively raising the attention elasticity $\gamma$ for familiar tasks: $\Delta \gamma \approx +0.02$ to $+0.08$.
\item \textbf{Hierarchical memory} organizes information at multiple granularities, reducing the effective signal count $K_\ell$ at each layer.
\end{itemize}

The key micro-implication is that memory design shifts $\bar{\eta}$ upward but cannot eliminate the bound. This provides a structural rationale for MACLA's observed ``optimal capacity bands'': beyond a certain memory size, the marginal gain in $\bar{\eta}$ diminishes (consistent with the concavity in Proposition~\ref{prop:memory_bound}), and the compounding of serial depreciation eventually dominates.

\subsection{Strategic Information Transmission}
\label{sec:strategic}

In the baseline model, each layer honestly reports its posterior to the next layer. In practice, agents may have incentives to distort information. For example, an upstream layer may want to hide unfavorable signals to avoid triggering costly downstream actions.

This can be modeled using the Bayesian Persuasion framework \citep{kamenica2011bayesian}. Layer $\ell$ chooses a signaling strategy to influence layer $\ell+1$'s action. The key difference from standard persuasion is that the signal is constrained by the attention budget: layer $\ell$ cannot transmit signals it did not attend to. The optimal persuasion strategy therefore interacts with the attention allocation: signals that are attended to can be strategically distorted, while ignored signals are simply absent.

Dynamic information transmission \citep{ely2017beeps, renault2017optimal} studies how senders strategically control information flow. While we focus on non-strategic loss from attention constraints, the same tools apply when agents distort information strategically.

When strategic distortion is possible, the principal may prefer to reduce the number of signals $K_\ell$ or increase the budget $A_\ell$ to ensure that the layer attends to all relevant signals (eliminating the ``hiding by ignoring'' possibility). This provides an alternative rationale for signal reduction that complements the precision-based rationale in Proposition \ref{prop:cost_min}.

\subsection{Human-in-the-Loop Review}
\label{subsec:human_ai}

When some layers are staffed by human agents with different attention technologies, the transmission factor becomes layer-specific: $\rho^{\text{human}}_\ell$ vs.~$\rho^{\text{AI}}_\ell$. The optimal chain may involve alternating human and AI layers to exploit comparative advantage in information processing.

\paragraph{Modeling human layers.} Unlike AI agents, human reviewers exhibit bounded rationality. We model a human layer with two parameters: $\gamma_{\text{human}}$ (attention elasticity, typically higher than AI for small $K$ but lower for large $K$) and $\beta \in (0,1]$ (bounded rationality coefficient). The human's posterior precision is:
\begin{equation}\label{eq:human_update}
\tau^{\text{human}}_\ell = \beta \cdot \tau^{\text{AI}}_{\ell-1} + (1-\beta) \cdot \tau^{\text{prior}}_{\text{human}} + \varepsilon_{\text{human}},
\end{equation}
where $\beta < 1$ captures sub-Bayesian updating (humans under-weight new information relative to their priors) and $\varepsilon_{\text{human}} \sim \mathcal{N}(0, \sigma^2_{\text{human}})$ captures idiosyncratic noise. When $\beta = 1$ and $\sigma_{\text{human}} = 0$, the human is fully Bayesian.

\paragraph{Optimal hybrid chains.} The interaction between AI and human layers generates novel comparative statics. Because human precision is bounded by cognitive capacity (typically $K_{\text{human}} \ll K_{\text{AI}}$), placing a human at a layer with many signals is inefficient. However, humans excel at tasks requiring judgment, contextual interpretation, and ethical reasoning---capabilities that are poorly captured by precision metrics alone.

\begin{proposition}[U-Shaped Optimal Depth in Hybrid Chains]\label{prop:hybrid_depth}
Consider a chain with one human layer and $L-1$ AI layers. The optimal depth $L^*_{\text{hybrid}}$ is U-shaped in the context window $A$: 
\begin{enumerate}[label=(\roman*)]
\item For small $A$: $L^*_{\text{hybrid}} < L^*_{\text{pure AI}}$ because human cognitive limits bind.
\item For intermediate $A$: $L^*_{\text{hybrid}} > L^*_{\text{pure AI}}$ because the human can effectively review AI outputs.
\item For large $A$: $L^*_{\text{hybrid}} \approx L^*_{\text{pure AI}}$ because AI precision dominates human review.
\end{enumerate}
\end{proposition}

\begin{proof}[Proof Sketch]
Small $A$: The AI's output precision is low, and the human's bounded capacity ($\beta < 1$) compounds the loss. The hybrid chain underperforms pure AI. Intermediate $A$: The AI produces high-quality outputs that the human can verify and augment, creating complementarities. Large $A$: The AI's precision approaches its technological ceiling; adding a human layer provides diminishing returns because the AI is already near-optimal.
\end{proof}

\textbf{Example.} In medical diagnosis, a human doctor may have $\gamma_{\text{human}} = 0.6$ (can effectively integrate a small number of key symptoms) while an LLM has $\gamma_{\text{LLM}} = 0.35$ (processes many lab results but with diminishing returns). The optimal chain may place the human at the top layer (final diagnosis) and the LLM at the bottom layer (lab result aggregation), rather than the reverse.

\section{Conclusion}
\label{sec:conclusion}

\subsection{Contributions}

We have developed a micro-founded theory of optimal organizational depth in AI delegation chains. The central insight is that \emph{information depreciation}---the degradation of signal quality through repeated algorithmic processing---compounds through a multi-layer chain, driving a finite, interior optimal depth even when the marginal cost of attention approaches zero.

Our framework makes three contributions relative to the existing literature on organizational economics:

First, we introduce the concept of information depreciation into organizational hierarchy theory. The transmission factor $\rho_\ell$ is the ratio of achieved posterior precision to first-best precision, and it is derived within the class of attention technologies rather than being assumed. This connects the organizational model to the Bayesian micro-foundations of rational inattention \citep{sims2003implications}.

Second, information depreciation provides a complementary mechanism to traditional coordination-cost explanations for the limits of delegation. Even as API costs fall to zero, the compounding of endogenous precision leakage and exogenous signal degradation at each layer ensures that information is progressively lost. This implies that Coasian arguments about the cost of coordination do not fully explain the boundaries of AI organizations: information depreciation operates alongside, and independently of, transaction costs.

Third, we provide a calibration framework that maps the model's parameters to directly observable LLM technology: context window lengths, API pricing, and benchmark accuracy. The calibration demonstrates that the key predictions---finite optimal depth, compounding precision loss, and the irrelevance of API costs for depth---hold for empirically plausible parameter values. We derive five testable predictions (Section~\ref{subsec:predictions}) that await empirical evaluation as multi-agent benchmark data become available.

\subsection{Policy Implications}

The model has direct implications for the design of AI systems and the regulation of AI services:

\begin{itemize}
\item \textbf{System designers} should front-load attention capacity: the largest models (largest context windows) should be deployed at the earliest layer, where precision has a multiplicative effect on all downstream outputs.
\item \textbf{Regulators} assessing AI system reliability should focus on depth as a key risk factor. A chain of four LLMs is not merely four times more expensive than a single LLM; it is exponentially less precise, and that precision loss is structurally unavoidable with current attention technology.
\end{itemize}

\subsection{Limitations and Future Directions}

Our model abstracts from several important features of real-world AI systems:

\begin{enumerate}
\item \textbf{Multi-dimensional states.} We assume a scalar state $\omega$ and scalar precision $\tau$. Generalizing to a vector state with a precision matrix would allow the model to capture the heterogeneity of information across different dimensions.
\item \textbf{Strategic behavior.} We assume honest transmission of posteriors. Incorporating strategic information transmission (as in Section~\ref{sec:strategic}) would generate a richer set of predictions about optimal monitoring and verification architectures.
\item \textbf{Endogenous architectures.} We take the chain structure as given. A fully general model would jointly optimize over chains, trees, and more complex graphs, as sketched in Section~\ref{subsec:skip_parallel}.
\item \textbf{Empirical validation.} While we have proposed five testable predictions (Section~\ref{subsec:predictions}), their empirical evaluation remains for future research. The rapid expansion of multi-agent benchmarks and the availability of API-level data from major providers make this a promising avenue. We have not conducted original empirical analysis; the calibration exercises in Section~\ref{sec:calibration} use publicly available benchmark statistics, and the testable predictions are offered as guidance for future empirical work rather than as validated findings.
\end{enumerate}

\subsection{Response to Reviewer Feedback}

This revision addresses the reviewer's concerns through four substantive changes. \textbf{First}, we have added a formal rate-distortion foundation (Appendix~\ref{app:rate_distortion}) with the correct reverse water-filling solution $D(R) = \sum_k \min\{\lambda_k, \theta\}$, tightening the micro-foundation for the preservation bound $\bar{\eta} < 1$. \textbf{Second}, we have provided an explicit mapping from rational inattention to the endogenous preservation rate (Proposition~\ref{prop:ri_to_eta}), showing how the attention budget $A_\ell$ translates into mutual information and then into $\eta_\ell$ via the Michaelis--Menten calibration. \textbf{Third}, we have characterized sufficient conditions for front-loading (Lemma~\ref{lem:front_loading_sufficient}) and the boundary where back-loading can arise (Corollary~\ref{cor:back_loading}). \textbf{Fourth}, we have expanded the extensions to cover alternative architectures (skip connections, parallel branches, memory augmentation), human-in-the-loop stages with bounded rationality, and a concrete identification strategy with pre-registered regression specifications.


\newpage
\appendix
\section{Proof of Proposition~\ref{prop:optimal_attention}}
\label{app:proof_prop1}

The agent's problem is to maximize $V_\ell(\tau_\ell)$ subject to $\sum_k \alpha_{\ell,k} \leq A_\ell$ and $\alpha_{\ell,k} \geq 0$. The Lagrangian is:
\begin{equation*}
\mathcal{L} = V_\ell\Biggl(\tau^{\text{prior}}_\ell + \sum_{k=1}^{K_\ell} \alpha_{\ell,k}^\gamma \tau^0_{\ell,k}\Biggr) - \lambda_\ell\Biggl(\sum_{k=1}^{K_\ell} \alpha_{\ell,k} - A_\ell\Biggr) + \sum_{k=1}^{K_\ell} \mu_{\ell,k} \alpha_{\ell,k}.
\end{equation*}

\textbf{Step 1: KKT conditions.} Because $V_\ell(\cdot)$ is strictly increasing and $g_\ell(\cdot) = (\cdot)^\gamma$ is strictly concave for $\gamma \in (0,1)$, the objective is strictly concave in $\boldsymbol{\alpha}_\ell$. The feasible set (a simplex intersected with the positive orthant) is convex and compact. Hence a unique maximizer exists and is characterized by the KKT conditions:
\begin{align}
V'_\ell(\tau^*_\ell) \gamma \alpha^{*\gamma-1}_{\ell,k} \tau^0_{\ell,k} - \lambda_\ell + \mu_{\ell,k} &= 0 \quad \text{(FOC)} \label{eq:kkt_foc}\\
\sum_{k=1}^{K_\ell} \alpha^*_{\ell,k} &= A_\ell \quad \text{(budget)} \label{eq:kkt_budget}\\
\alpha^*_{\ell,k} \geq 0, \; \mu_{\ell,k} \geq 0, \; \mu_{\ell,k} \alpha^*_{\ell,k} &= 0 \quad \text{(CS)}. \label{eq:kkt_cs}
\end{align}

\textbf{Step 2: Corner solutions.} For any signal $k$ with $\alpha^*_{\ell,k} = 0$, complementary slackness implies $\mu_{\ell,k} \geq 0$. The FOC becomes:
\begin{equation*}
\mu_{\ell,k} = \lambda_\ell - V'_\ell(\tau^*_\ell) \gamma \cdot \lim_{\alpha \to 0} \alpha^{\gamma-1} \tau^0_{\ell,k} = +\infty > 0,
\end{equation*}
because $\gamma - 1 < 0$ and $\lim_{\alpha \to 0} \alpha^{\gamma-1} = +\infty$. Thus the FOC is satisfied for any $\alpha^*_{\ell,k} = 0$.

\textbf{Step 3: Interior solutions.} For $\alpha^*_{\ell,k} > 0$, we have $\mu_{\ell,k} = 0$ and the FOC yields:
\begin{equation*}
\alpha^*_{\ell,k} = \biggl(\frac{\gamma \tau^0_{\ell,k}}{\lambda_\ell}\biggr)^{\frac{1}{1-\gamma}}.
\end{equation*}
This is increasing in $\tau^0_{\ell,k}$ and decreasing in $\lambda_\ell$. Rearranging gives Equation~\eqref{eq:closed_form}.

\textbf{Step 4: Determining the active set.} Let signals be ordered by precision: $\tau^0_{\ell,1} \geq \tau^0_{\ell,2} \geq \cdots \geq \tau^0_{\ell,K_\ell}$. Define the threshold precision:
\begin{equation*}
\bar{\tau}(\lambda_\ell) = \frac{\lambda_\ell}{V'_\ell(\tau^*_\ell) \gamma} \cdot A_\ell^{1-\gamma}.
\end{equation*}
Then $\alpha^*_{\ell,k} > 0$ if and only if $\tau^0_{\ell,k} > \bar{\tau}(\lambda_\ell)$. The active set is $\mathcal{A}_\ell = \{k : \tau^0_{\ell,k} > \bar{\tau}\}$. Because $\bar{\tau}$ decreases in $A_\ell$, more signals become active as the budget grows.

\textbf{Step 5: Uniqueness.} The objective is strictly concave (composition of strictly increasing concave $V_\ell$ with strictly concave precision function), and the constraint set is convex. A strictly concave function on a convex set has a unique maximizer. \hfill $\square$

\section{Proof of Proposition~\ref{prop:distortion_comp_statics}}
\label{app:proof_prop2}

Recall $\Delta_\ell = 1 - \rho_\ell$, where $\rho_\ell$ is given by Definition~\ref{def:transmission_factor}. For the bottom layer with $\tau^{\text{prior}}_\ell = 0$, the transmission factor reduces to:
\begin{equation*}
\rho_\ell = \frac{\sum_k g_\ell(\alpha^*_{\ell,k}) \tau^0_{\ell,k}}{\sum_k \tau^0_{\ell,k}}.
\end{equation*}

\textbf{(i) Budget effect.} From the KKT solution (Proposition~\ref{prop:optimal_attention}), interior attention weights satisfy $\alpha^*_{\ell,k} = (\gamma \tau^0_{\ell,k} / \lambda_\ell)^{1/(1-\gamma)}$. Summing over the active set and imposing the budget constraint yields:
\begin{equation*}
\lambda_\ell = \gamma \Biggl(\frac{\sum_{k \in \mathcal{A}_\ell} (\tau^0_{\ell,k})^{\frac{1}{1-\gamma}}}{A_\ell}\Biggr)^{1-\gamma}.
\end{equation*}
Substituting back, the optimal weight is:
\begin{equation*}
\alpha^*_{\ell,k} = A_\ell \cdot \frac{(\tau^0_{\ell,k})^{\frac{1}{1-\gamma}}}{\sum_{j \in \mathcal{A}_\ell} (\tau^0_{\ell,j})^{\frac{1}{1-\gamma}}}, \quad k \in \mathcal{A}_\ell.
\end{equation*}
Hence $\alpha^*_{\ell,k} = c_k A_\ell$ where $c_k > 0$ depends only on the relative precisions. With $g_\ell(\alpha) = \alpha^\gamma$, we have $g_\ell(\alpha^*_{\ell,k}) = (c_k A_\ell)^\gamma$. The transmission factor becomes:
\begin{equation*}
\rho_\ell = \frac{A_\ell^\gamma \sum_{k \in \mathcal{A}_\ell} c_k^\gamma \tau^0_{\ell,k}}{\sum_{k=1}^{K_\ell} \tau^0_{\ell,k}},
\end{equation*}
which is increasing in $A_\ell$. Therefore $\partial \Delta_\ell / \partial A_\ell = -\partial \rho_\ell / \partial A_\ell < 0$.

\textbf{(ii) Signal count effect.} When $K_\ell$ increases while $A_\ell$ is fixed, the budget is spread more thinly across signals. The marginal signals added have lower intrinsic precision, reducing the average precision per signal. The transmission factor $\rho_\ell$ decreases because the denominator $\sum_k \tau^0_{\ell,k}$ increases (holding average precision constant) while the numerator increases more slowly due to the concavity of $g_\ell$.

\textbf{(iii) Heterogeneity effect.} Consider two precision distributions with the same mean $\bar{\tau}^0$ but different variances. Under uniform attention $\alpha^*_{\ell,k} = A_\ell/K_\ell$ (optimal for identical signals), $\rho^{\text{uniform}}_\ell = (A_\ell/K_\ell)^\gamma$. Under a heterogeneous distribution with a dominant signal $k=1$, the optimal allocation concentrates attention: $\alpha^*_{\ell,1} > A_\ell/K_\ell > \alpha^*_{\ell,k}$ for $k > 1$. Because $g(\cdot)$ is concave, Jensen's inequality implies $\sum_k g(\alpha^*_{\ell,k}) \tau^0_{\ell,k}$ is higher under heterogeneous concentration, holding $\sum_k \tau^0_{\ell,k}$ fixed. Thus $\rho^{\text{heterogeneous}}_\ell > \rho^{\text{uniform}}_\ell$ and $\Delta_\ell$ is lower. \hfill $\square$

\section{Proof of Proposition~\ref{prop:precision_response}}
\label{app:proof_prop3}

From the recursive structure (Equation~\eqref{eq:recursive_precision}), the top-layer precision satisfies:
\begin{equation*}
\tau^*_L = \rho_{L-1} \cdot \tau^*_{L-1} + \sum_{k=1}^{K_L} g_L(\alpha^*_{L,k}) \tau^0_{L,k},
\end{equation*}
with $\tau^*_0 = \tau^0_0$ and $\rho_\ell \in [0,1)$ for any finite $A_\ell$. Define the first difference $\Delta \tau^*_L \equiv \tau^*_L(L+1) - \tau^*_L(L)$.

Adding a layer increases the chain length but also adds a new precision source. The net effect is:
\begin{equation*}
\Delta \tau^*_L = \underbrace{\rho_L \cdot \tau^*_L}_{\text{updated inherited precision}} + \underbrace{\sum_{k=1}^{K_{L+1}} g_{L+1}(\alpha^*_{L+1,k}) \tau^0_{L+1,k}}_{\text{new layer's contribution}} - \tau^*_L.
\end{equation*}
Since $\rho_L < 1$, the first term is less than $\tau^*_L$. The second term is bounded above by the first-best precision. For any interior solution (finite $A_{L+1}$), the loss from compression exceeds the gain from the new layer, so $\Delta \tau^*_L < 0$.

For the second difference, each additional layer reduces both the inherited precision and the marginal contribution of the new layer. Hence:
\begin{equation*}
\Delta \tau^*_L(L+1) - \Delta \tau^*_L(L) < 0,
\end{equation*}
which is the discrete analog of the second derivative being negative. \hfill $\square$

\section{Proof of Proposition~\ref{prop:cost_min}}
\label{app:proof_prop4}

\textbf{Step 1: Define the objective.} Let $\Pi(L) = V_L(\tau^*_L) - \sum_{\ell=0}^{L-1} c_\ell(A_\ell)$ denote the principal's net value when the chain has $L$ layers.

\textbf{Step 2: Concavity of $\tau^*_L$.} From Proposition~\ref{prop:geometric_decay}, the first difference $\Delta \tau^*_L \equiv \tau^*_L(L+1) - \tau^*_L(L) < 0$, and the second difference is also negative. Hence $\tau^*_L$ is strictly concave on the discrete domain.

\textbf{Step 3: Concavity of $V_L(\tau^*_L)$.} The value function $V_L(\cdot)$ is strictly increasing and strictly concave. The composition of a strictly increasing concave function with a strictly concave function is strictly concave. Hence $V_L(\tau^*_L(L))$ is strictly concave in $L$.

\textbf{Step 4: Convexity of costs.} The cost term $C(L) = \sum_{\ell=0}^{L-1} c_\ell(A_\ell)$ is linear in $L$ when budgets are fixed. If budgets increase with depth, the cost is convex because $c_\ell(\cdot)$ is convex.

\textbf{Step 5: Sum of concave functions.} The sum of strictly concave functions is strictly concave. Hence $\Pi(L) = V_L(\tau^*_L(L)) - C(L)$ is strictly concave on $\{0, 1, 2, \ldots\}$.

\textbf{Step 6: Boundary behavior.} At $L = 0$, $\Pi(0) = V_0(\tau^0_0) > -\infty$. As $L \to \infty$, $\tau^*_L \to 0$ because $\rho_\ell < 1$ and the product $\prod_{\ell=0}^{L-1} \rho_\ell \to 0$ geometrically. Hence $V_L(\tau^*_L) \to V_L(0)$, while $C(L) \to +\infty$. Therefore $\lim_{L \to \infty} \Pi(L) = -\infty$.

\textbf{Step 7: Existence of unique maximizer.} A strictly concave function on a discrete domain with $\Pi(0) > -\infty$ and $\lim_{L \to \infty} \Pi(L) = -\infty$ has a unique finite maximizer. \hfill $\square$

\section{Proof of Proposition~\ref{prop:front_loading}}
\label{app:proof_prop5}

\textbf{Step 1: Discrete supermodularity.} Let $\Pi(L; A_\ell) \equiv V_L(\tau^*_L(A_\ell)) - \sum_{j=0}^{L-1} c_j(A_j)$ denote the principal's net value. We show that $\Pi$ is supermodular in $(L, A_\ell)$ on $\mathbb{N} \times \mathbb{R}_+$.

\textbf{Step 2: Decompose the difference.} For any $L \geq 0$ and $A'_\ell > A_\ell$:
\begin{align*}
\Delta_L \Pi(L; A'_\ell) - \Delta_L \Pi(L; A_\ell) &= \bigl[V_{L+1}(\tau^*_{L+1}(A'_\ell)) - V_{L+1}(\tau^*_{L+1}(A_\ell))\bigr] \\
&\quad - \bigl[V_L(\tau^*_L(A'_\ell)) - V_L(\tau^*_L(A_\ell))\bigr].
\end{align*}

\textbf{Step 3: Sign of each term.} The top-layer precision with $L+1$ layers is:
\begin{equation*}
\tau^*_{L+1}(A_\ell) = \rho_L(A_\ell) \cdot \tau^*_L + \sum_{k=1}^{K_{L+1}} g_{L+1}(\alpha^*_{L+1,k}) \tau^0_{L+1,k}.
\end{equation*}
From Proposition~\ref{prop:distortion_comp_statics}, $\partial \rho_\ell / \partial A_\ell > 0$. Hence $\tau^*_{L+1}(A'_\ell) > \tau^*_{L+1}(A_\ell)$, and since $V_{L+1}(\cdot)$ is increasing, term (i) $> 0$.

Term (ii) is similarly positive: a larger budget at layer $\ell$ raises the precision transmitted to all downstream layers.

\textbf{Step 4: Compare the gains.} The budget increase at layer $\ell$ affects all layers $j > \ell$ through the transmission chain. At layer $L+1$, the effect is amplified by the additional layer of compounding. Because both $V_{L+1}$ and $V_L$ are increasing and concave, and the precision gain is larger at $L+1$ than at $L$, the difference (i) $-$ (ii) $> 0$.

\textbf{Step 5: Apply Topkis's theorem.} The cost term does not depend on $A_\ell$, so it does not affect the cross-difference. Since $\Delta_L \Pi(L; A'_\ell) > \Delta_L \Pi(L; A_\ell)$, the objective is supermodular in $(L, A_\ell)$. By Topkis's monotonicity theorem, $L^*(A_\ell) \equiv \arg\max_L \Pi(L; A_\ell)$ is increasing in $A_\ell$. \hfill $\square$

\section{Proof of Proposition~\ref{prop:precision_path}}
\label{app:proof_precision_path}

\paragraph{Step 1: Derivation of the recurrence.}
Under uniform technology and uniform attention budgets ($A_\ell = A$ for all $\ell$, so $\eta_\ell = \eta$ by Assumptions~\ref{ass:state_signals}--\ref{ass:endogenous_eta}), the new signal contribution at layer $\ell$ is $\sum_{k=1}^{K} g(\alpha^*_k) \tau^0_k \eta^\ell = C\eta^\ell$, where $C$ is defined in \eqref{eq:C_constant}. Substituting into \eqref{eq:recursive_precision} yields \eqref{eq:recursive_tau_unified}.

\paragraph{Step 2: Closed-form solution.}
Equation~\eqref{eq:recursive_tau_unified} is a first-order linear non-homogeneous difference equation.

\textit{Case 1: $\eta \neq \rho$.} The homogeneous solution is $\tau^{(h)}_{\ell} = A \rho^{\ell}$. Guessing a particular solution $\tau^{(p)}_{\ell} = B \eta^{\ell}$ and substituting:
\begin{equation*}
B\eta^{\ell} = \rho B \eta^{\ell-1} + C\eta^{\ell}
\quad\Longrightarrow\quad
B = \frac{C\eta}{\eta - \rho}.
\end{equation*}
Applying $\tau^{*}_{0} = A + B$ and solving for $\tau^{*}_{L}$ yields \eqref{eq:tau_star_L}.

\textit{Case 2: $\eta = \rho$.} The recurrence becomes $\tau^{*}_{\ell} = \eta \tau^{*}_{\ell-1} + C\eta^{\ell}$. Defining $u_{\ell} \equiv \tau^{*}_{\ell}/\eta^{\ell}$, we have $u_{\ell} = u_{\ell-1} + C$, which telescopes to $u_{\ell} = \tau^{*}_{0} + C\ell$. Hence $\tau^{*}_{L} = \eta^{L}(\tau^{*}_{0} + CL)$.

\paragraph{Step 3: Vanishing precision (all cases).}
Since $\rho \in (0,1)$ and $\eta \leq \bar{\eta} < 1$ (by Assumption~\ref{ass:endogenous_eta}), both $\rho^{L} \to 0$ and $\eta^{L} \to 0$. In all three cases, $\lim_{L\to\infty} \tau^{*}_{L} = 0$. The bound $\bar{\eta} < 1$ is the structural reason why precision vanishes even with unlimited attention budgets.

\paragraph{Step 4: Shape characterization.}
From \eqref{eq:recursive_tau_unified}, the layer-to-layer change is
\begin{equation}\label{eq:app_delta_tau}
\Delta_{\ell} \equiv \tau^{*}_{\ell} - \tau^{*}_{\ell-1} = -(1-\rho)\tau^{*}_{\ell-1} + C\eta^{\ell}.
\end{equation}

\textit{Subcase (i): $\eta < \rho$.} The $\rho^{\ell}$ term dominates the closed-form solution for large $\ell$. Since $\rho > \eta$, the negative term $-(1-\rho)\tau^{*}_{\ell-1}$ (which decays as $\rho^{\ell-1}$) eventually dominates the positive term $C\eta^{\ell}$ (which decays as $\eta^{\ell}$). Thus $\Delta_{\ell} < 0$ for all $\ell \geq \bar{L}_1$. For small $\ell$, however, the positive term may dominate if $\tau^{*}_{0}$ is small relative to $C\eta$, causing an initial increase. The path is therefore hump-shaped.

\textit{Subcase (ii): $\eta > \rho$.} As $\ell \to \infty$, $\tau^{*}_{\ell-1} \sim C\eta^{\ell}/(\eta-\rho)$. Substituting into \eqref{eq:app_delta_tau}:
\begin{equation*}
\Delta_{\ell} \sim C\eta^{\ell}\left(1 - \frac{1-\rho}{\eta-\rho}\right) = C\eta^{\ell}\frac{\eta - 1}{\eta - \rho} < 0,
\end{equation*}
since $\eta < 1$ and $\eta > \rho$. The sequence is eventually decreasing.

\textit{Subcase (iii): $\eta = \rho$.} Using $\tau^{*}_{\ell} = \eta^{\ell}(\tau^{*}_{0} + C\ell)$:
\begin{equation*}
\Delta_{\ell} = \eta^{\ell-1}\bigl[(\eta-1)\tau^{*}_{0} + C - C(1-\eta)\ell\bigr].
\end{equation*}
Since $\eta < 1$, the bracket is affine in $\ell$ with negative slope $-C(1-\eta) < 0$. Hence $\Delta_{\ell} < 0$ whenever $\ell > \tau^{*}_{0}/C + 1/(1-\eta)$. \hfill $\blacksquare$

\section{Proof of Proposition~\ref{prop:cost_depth}}
\label{app:proof_prop6}

\textbf{Step 1: Discrete submodularity.} Let $\Pi(L; \kappa_\ell) \equiv V_L(\tau^*_L) - \sum_{j=0}^{L-1} c_j(A_j; \kappa_\ell)$, where $c_j(A_j; \kappa_\ell) = \kappa_\ell A_\ell + \frac{\phi_\ell}{2} A_\ell^2$. We show that $\Pi$ is submodular in $(L, \kappa_\ell)$.

\textbf{Step 2: Cost effect on the marginal value.} The difference in marginal values is:
\begin{multline*}
\Delta_L \Pi(L; \kappa'_\ell) - \Delta_L \Pi(L; \kappa_\ell) = -\bigl[c_\ell(A_\ell; \kappa'_\ell) - c_\ell(A_\ell; \kappa_\ell)\bigr] \cdot \mathbf{1}_{\{\ell < L+1\}} \\
= -(\kappa'_\ell - \kappa_\ell) A_\ell \cdot \mathbf{1}_{\{\ell < L+1\}} \leq 0.
\end{multline*}
For $\ell < L+1$, a higher cost reduces the marginal net value of adding the $(L+1)$-th layer.

\textbf{Step 3: Monotonicity.} Because $\Delta_L \Pi(L; \kappa'_\ell) \leq \Delta_L \Pi(L; \kappa_\ell)$, the objective is submodular in $(L, \kappa_\ell)$. By Topkis's theorem, $L^*(\kappa_\ell) \equiv \arg\max_L \Pi(L; \kappa_\ell)$ is decreasing in $\kappa_\ell$. \hfill $\square$

\section{Proof of Proposition~\ref{prop:finite_depth}}
\label{app:proof_prop7}

\renewcommand{\theequation}{A.7.\arabic{equation}}
\setcounter{equation}{0}

This appendix provides the complete proof of Proposition~\ref{prop:finite_depth}, which establishes the finiteness of optimal depth under diminishing signal quality. The proof proceeds in five logical steps.

\paragraph{Step 1: Recursive structure.}
Under uniform technology ($g_{\ell} = g$, $c_{\ell} = c$, $K_{\ell} = K$, $A_{\ell} = A$, so $\eta_{\ell} = \eta$ for all $\ell$) and Assumptions~\ref{ass:state_signals}--\ref{ass:endogenous_eta}, the equilibrium posterior precision at layer $\ell$ evolves according to:
\begin{equation}\label{eq:app_recursive}
\tau^{*}_{\ell} = \underbrace{\rho \cdot \tau^{*}_{\ell-1}}_{\text{precision carried forward}} + \underbrace{\sum_{k=1}^{K} g\bigl(\alpha^{*}_{k}\bigr) \tau^{0}_{\ell,k}}_{\text{new signal contribution}}.
\end{equation}
By Assumptions~\ref{ass:state_signals}--\ref{ass:endogenous_eta} with uniform budgets ($\eta_\ell = \eta$), $\tau^{0}_{\ell,k} = \tau^{0}_{k} \eta^{\ell}$, so the new signal contribution becomes:
\begin{equation}\label{eq:app_C_def}
\sum_{k=1}^{K} g\bigl(\alpha^{*}_{k}\bigr) \tau^{0}_{k} \eta^{\ell} = C \eta^{\ell},
\end{equation}
where $C$ is defined in \eqref{eq:C_constant}. The initial condition is $\tau^{*}_{0} = \sum_{k=1}^{K} \tau^{0}_{0,k} = \sum_{k=1}^{K} \tau^{0}_{k}$. Thus \eqref{eq:app_recursive} simplifies to:
\begin{equation}\label{eq:app_simplified_recurrence}
\tau^{*}_{\ell} = \rho \tau^{*}_{\ell-1} + C \eta^{\ell}, \qquad \tau^{*}_{0} \text{ given}.
\end{equation}

\paragraph{Step 2: Closed-form solution.}
Equation~\eqref{eq:app_simplified_recurrence} is a first-order linear non-homogeneous difference equation with varying coefficients in the forcing term.

\textit{Case 1: $\eta \neq \rho$.} Decompose $\tau^{*}_{\ell} = \tau^{(h)}_{\ell} + \tau^{(p)}_{\ell}$. The homogeneous solution is $\tau^{(h)}_{\ell} = A \rho^{\ell}$, and the ansatz $\tau^{(p)}_{\ell} = B \eta^{\ell}$ yields:
\begin{equation*}
B\eta^{\ell} = \rho B \eta^{\ell-1} + C\eta^{\ell}
\quad\Longrightarrow\quad
B\eta = \rho B + C\eta
\quad\Longrightarrow\quad
B = \frac{C\eta}{\eta - \rho}.
\end{equation*}
Applying the initial condition $\tau^{*}_{0} = A + B$:
\begin{equation}
A = \tau^{*}_{0} - \frac{C\eta}{\eta - \rho}.
\end{equation}
Therefore, for all $L \geq 0$:
\begin{equation}\label{eq:app_closed_form}
\tau^{*}_{L} = \rho^{L}\tau^{*}_{0} + \frac{C\eta}{\eta - \rho}\bigl(\eta^{L} - \rho^{L}\bigr),
\end{equation}
which coincides with equation~\eqref{eq:tau_star_L} in the main text.

\textit{Case 2: $\eta = \rho$.} The recurrence becomes $\tau^{*}_{\ell} = \eta \tau^{*}_{\ell-1} + C\eta^{\ell}$. Dividing both sides by $\eta^{\ell}$:
\begin{equation}
\frac{\tau^{*}_{\ell}}{\eta^{\ell}} = \frac{\tau^{*}_{\ell-1}}{\eta^{\ell-1}} + C.
\end{equation}
Defining $u_{\ell} \equiv \tau^{*}_{\ell}/\eta^{\ell}$, this becomes $u_{\ell} = u_{\ell-1} + C$, which telescopes to $u_{\ell} = u_{0} + C\ell = \tau^{*}_{0} + C\ell$. Restoring $\tau^{*}_{L}$:
\begin{equation}\label{eq:app_equal_case}
\tau^{*}_{L} = \eta^{L}\bigl(\tau^{*}_{0} + CL\bigr).
\end{equation}

\paragraph{Step 3: Vanishing precision (Part (i)).}
We establish that $\lim_{L\to\infty} \tau^{*}_{L} = 0$ in all cases.

For $\eta \neq \rho$, since $\rho, \eta \in (0,1)$:
\begin{equation}
\lim_{L\to\infty} \rho^{L} = 0, \qquad \lim_{L\to\infty} \eta^{L} = 0,
\end{equation}
so both terms in \eqref{eq:app_closed_form} converge to zero. For $\eta = \rho$, the linear growth of $\tau^{*}_{0} + CL$ is asymptotically dominated by the exponential decay of $\eta^{L}$:
\begin{equation}
\lim_{L\to\infty} \eta^{L}\bigl(\tau^{*}_{0} + CL\bigr) = 0,
\end{equation}
since $\lim_{L\to\infty} L^{m} \eta^{L} = 0$ for any $m \geq 0$ and $\eta \in (0,1)$. This completes Part (i).

\paragraph{Step 4: Eventual monotonic decline (Part (ii)).}
From \eqref{eq:app_simplified_recurrence}, the increment is:
\begin{equation}\label{eq:app_delta}
\Delta_{\ell} \equiv \tau^{*}_{\ell} - \tau^{*}_{\ell-1} = -(1-\rho)\tau^{*}_{\ell-1} + C\eta^{\ell}.
\end{equation}
We prove $\Delta_{\ell} < 0$ for all sufficiently large $\ell$ by considering three subcases.

\textit{Subcase 4a: $\eta > \rho$.} From \eqref{eq:app_closed_form}, as $\ell \to \infty$, the term $\eta^{\ell}$ dominates $\rho^{\ell}$, yielding:
\begin{equation}
\tau^{*}_{\ell-1} \sim \frac{C\eta^{\ell}}{\eta - \rho} \qquad (\ell \to \infty).
\end{equation}
Substituting into \eqref{eq:app_delta}:
\begin{equation}
\Delta_{\ell} \sim -(1-\rho)\frac{C\eta^{\ell}}{\eta-\rho} + C\eta^{\ell}
= C\eta^{\ell}\left(\frac{\eta - \rho - 1 + \rho}{\eta - \rho}\right)
= C\eta^{\ell}\frac{\eta - 1}{\eta - \rho}.
\end{equation}
Since $\eta < 1$, the numerator $\eta - 1 < 0$, while the denominator $\eta - \rho > 0$. Hence $\Delta_{\ell} \sim \text{[negative constant]} \times \eta^{\ell} < 0$ for all large $\ell$.

\textit{Subcase 4b: $\eta < \rho$.} Now $\rho^{\ell}$ dominates, giving:
\begin{equation}
\tau^{*}_{\ell-1} \sim \left(\tau^{*}_{0} + \frac{C\eta}{\rho - \eta}\right)\rho^{\ell-1} \qquad (\ell \to \infty).
\end{equation}
Substituting into \eqref{eq:app_delta}:
\begin{equation}
\Delta_{\ell} \sim -(1-\rho)\left(\tau^{*}_{0} + \frac{C\eta}{\rho-\eta}\right)\rho^{\ell-1} + C\eta^{\ell}.
\end{equation}
Since $\rho > \eta$, the negative term decays as $\rho^{\ell-1}$ while the positive term decays as $\eta^{\ell}$. The former decays strictly slower and has a negative coefficient, so it dominates for large $\ell$. Thus $\Delta_{\ell} < 0$ for all $\ell \geq \bar{L}$ with:
\begin{equation}
\bar{L} = \inf\left\{\ell \in \mathbb{N} : \ell > \frac{\ln\!\left(\frac{(1-\rho)}{C}\bigl(\tau^{*}_{0} + \frac{C\eta}{\rho-\eta}\bigr)\right)}{\ln(\eta/\rho)}\right\} + 1,
\end{equation}
which is finite since $\ln(\eta/\rho) < 0$.

\textit{Subcase 4c: $\eta = \rho$.} Using \eqref{eq:app_equal_case}:
\begin{align}
\Delta_{\ell} &= \eta^{\ell}(\tau^{*}_{0} + C\ell) - \eta^{\ell-1}\bigl(\tau^{*}_{0} + C(\ell-1)\bigr) \notag \\
&= \eta^{\ell-1}\bigl[\underbrace{(\eta-1)\tau^{*}_{0}}_{< 0} + \underbrace{C(1 - (1-\eta)\ell)}_{\to -\infty}\bigr].
\end{align}
Since the bracket is affine in $\ell$ with negative slope $-C(1-\eta) < 0$, there exists:
\begin{equation}
\bar{L} = \left\lfloor\frac{\tau^{*}_{0}}{C} + \frac{1}{1-\eta}\right\rfloor + 2,
\end{equation}
such that $\Delta_{\ell} < 0$ for all $\ell \geq \bar{L}$. This completes Part (ii).

\paragraph{Step 5: Finite optimal depth (Part (iii)).}
Strict monotonicity of $V(\cdot)$ and Part (ii) together imply:
\begin{equation}
V(\tau^{*}_{\ell}) < V(\tau^{*}_{\ell-1}) \quad \text{for all } \ell \geq \bar{L}.
\end{equation}
Thus $\{V(\tau^{*}_{\ell})\}_{\ell=0}^{\infty}$ is strictly decreasing beyond $\bar{L}$, so its maximum is attained at some $L^{*} \in \{0, 1, \ldots, \bar{L}\}$.

When the marginal attention cost is $\kappa \geq 0$, adding any layer $\ell \geq \bar{L}$ yields net marginal benefit:
\begin{equation}
\underbrace{V(\tau^{*}_{\ell}) - V(\tau^{*}_{\ell-1})}_{< 0} - \kappa < 0,
\end{equation}
which is strictly negative. Hence no planner would ever choose $\ell > \bar{L}$, and $L^{*} \leq \bar{L} < \infty$ regardless of how small $\kappa$ is. This completes the proof of Proposition~\ref{prop:finite_depth}. \hfill $\blacksquare$

\vspace{1.5em}
\begin{remark}[Comparative statics on $\bar{L}$]
The threshold $\bar{L}$ beyond which precision decreases is increasing in the baseline precision $\tau^{*}_{0}$ and decreasing in the signal contribution $C$ and the attenuation factor $\eta$. Specifically, a higher $\eta$ (better signal preservation) pushes $\bar{L}$ outward, allowing more layers before returns turn negative. In the limit $\eta \to 1$, the threshold diverges, recovering the original model where infinite depth can be optimal.
\end{remark}

\begin{remark}[Comparison to the no-attenuation case]
When $\eta = 1$ (no signal degradation), the recurrence \eqref{eq:app_simplified_recurrence} becomes $\tau^{*}_{\ell} = \rho \tau^{*}_{\ell-1} + C$, with solution $\tau^{*}_{L} = \rho^{L}\tau^{*}_{0} + C(1-\rho^{L})/(1-\rho) \to C/(1-\rho) > 0$. The precision converges to a positive limit, so $V(\tau^{*}_{L})$ converges to $V(C/(1-\rho)) > V(0)$ from below, and infinite depth is optimal whenever $\kappa < V(C/(1-\rho)) - \lim_{L\to\infty}V(\tau^{*}_{L-1}) = 0^{+}$. The diminishing signal quality assumption ($\eta < 1$) is therefore \emph{necessary and sufficient} for the finite-depth result under uniform technology.
\end{remark}

\section{A Rate-Distortion Foundation for the Preservation Bound}
\label{app:rate_distortion}

This appendix provides a rigorous rate-distortion foundation for the preservation bound $\bar{\eta} < 1$, addressing the reviewer's concern that earlier formulations may have oversimplified the reverse water-filling solution.

\subsection{Setup}

Consider a $d$-dimensional Gaussian source $\mathbf{x} \sim \mathcal{N}(0, \Sigma)$ representing the latent state of an LLM's hidden representation. The covariance matrix $\Sigma$ has eigenvalue decomposition $\Sigma = Q \Lambda Q^{\top}$, where $\Lambda = \text{diag}(\lambda_1, \ldots, \lambda_d)$ with $\lambda_1 \geq \lambda_2 \geq \cdots \geq \lambda_d > 0$.

Each layer $\ell$ must transmit a representation $\hat{\mathbf{x}}_\ell$ to the next layer subject to a rate constraint $R_\ell$ (in nats per dimension). The distortion is measured by mean squared error: $D = \mathbb{E}[\|\mathbf{x} - \hat{\mathbf{x}}\|^2]$.

\subsection{Reverse Water-Filling}

The rate-distortion function for a parallel Gaussian source is given by reverse water-filling \citep{cover1999elements}:
\begin{equation}\label{eq:reverse_water_filling}
D(R) = \sum_{k=1}^{d} \min\{\lambda_k, \theta(R)\},
\end{equation}
where the water level $\theta(R)$ is chosen such that
\begin{equation}\label{eq:water_level}
\sum_{k=1}^{d} [\lambda_k - \theta(R)]^{+} = d \cdot R.
\end{equation}

This is the \emph{correct} expression. The oversimplified formula $D_{\min}(d) = d \cdot \theta^*$ is valid only when $\theta(R) \leq \lambda_d$ (i.e., when the rate budget is so small that all modes are active). In general, the number of active modes is
\begin{equation}\label{eq:active_modes}
m(R) = \max\left\{m \in \{1, \ldots, d\} : \lambda_m > \theta(R)\right\},
\end{equation}
and the distortion is $D(R) = \sum_{k=1}^{m(R)} \theta(R) + \sum_{k=m(R)+1}^{d} \lambda_k$.

\subsection{The Preservation Bound}

Define the preservation rate as the ratio of retained signal variance to total variance:
\begin{equation}\label{eq:rd_preservation}
\eta(R) = 1 - \frac{D(R)}{\sum_{k=1}^{d} \lambda_k}.
\end{equation}

\begin{proposition}[Rate-Distortion Preservation Bound]\label{prop:rd_bound}
For any finite rate budget $R < \infty$:
\begin{enumerate}[label=(\roman*)]
\item $\eta(R) < 1$.
\item $\lim_{R \to \infty} \eta(R) = 1$ if and only if $\sum_{k=1}^{d} \lambda_k < \infty$ (finite-dimensional source).
\item For any $R$, the preservation rate is bounded above by
\begin{equation}\label{eq:rd_eta_bound}
\eta(R) \leq 1 - \frac{(d - m(R)) \lambda_d}{\sum_{k=1}^{d} \lambda_k},
\end{equation}
where $m(R)$ is the number of active modes defined in \eqref{eq:active_modes}.
\end{enumerate}
\end{proposition}

\begin{proof}
\textbf{Part (i):} From \eqref{eq:reverse_water_filling}, $D(R) \geq \sum_{k=m(R)+1}^{d} \lambda_k > 0$ because at most $m(R) < d$ modes can be active for finite $R$. Hence $\eta(R) = 1 - D(R) / \text{tr}(\Sigma) < 1$.

\textbf{Part (ii):} As $R \to \infty$, the water level $\theta(R) \to 0$, so $m(R) \to d$ and $D(R) \to 0$. Thus $\eta(R) \to 1$. Conversely, if the source is infinite-dimensional ($d = \infty$) with $\sum_k \lambda_k = \infty$, then even with $R \to \infty$, the tail sum $\sum_{k=m+1}^{\infty} \lambda_k$ may remain positive, preventing $\eta$ from reaching 1.

\textbf{Part (iii):} The inactive modes contribute at least $(d - m(R)) \lambda_d$ to distortion (since $\lambda_k \geq \lambda_d$ for all $k \leq d$). Substituting into \eqref{eq:rd_preservation} yields the bound.
\end{proof}

\subsection{Connection to the Model}

In our framework, the ``rate budget'' $R$ corresponds to the attention capacity per signal, and the eigenvalues $\{\lambda_k\}$ correspond to the singular values of the hidden-state transformation at each layer. Proposition~\ref{prop:rd_bound} shows that $\bar{\eta} < 1$ is not an arbitrary assumption but a consequence of finite rate budgets and the structure of the source covariance.

The bound is \emph{tightest} when the eigenvalue spectrum decays slowly (many modes with comparable variance) and \emph{loosest} when the spectrum is dominated by a few large eigenvalues. This explains why tasks requiring fine-grained detail (slow eigenvalue decay) have lower $\bar{\eta}$ than tasks where coarse structure suffices (rapid eigenvalue decay).

\begin{remark}[Operationalizing the Bound]
The eigenvalue spectrum $\{\lambda_k\}$ can be estimated from the empirical covariance of LLM hidden states. For a transformer with hidden dimension $h$, sample the hidden-state vectors across a validation corpus, compute the sample covariance $\hat{\Sigma}$, and extract the eigenvalues. The implied preservation bound is $\hat{\bar{\eta}} = 1 - D(R) / \text{tr}(\hat{\Sigma})$ for the observed rate budget $R = A_\ell / K_\ell$.
\end{remark}


\newpage
\section*{References}
\addcontentsline{toc}{section}{References}
\begin{thebibliography}{99}

\bibitem[Acemoglu and Restrepo (2018)]{acemoglu2018race}
Acemoglu, Daron, and Pascual Restrepo. 2018.
``The Race between Man and Machine: Implications of Technology for Growth, Factor Shares, and Employment.''
\textit{American Economic Review}, 108(6): 1488--1542.

\bibitem[Acemoglu and Restrepo (2022)]{acemoglu2022tasks}
Acemoglu, Daron, and Pascual Restrepo. 2022.
``Tasks, Automation, and the Rise in US Wage Inequality.''
\textit{Econometrica}, 90(4): 1973--2016.

\bibitem[Acemoglu, Autor, Hazell, and Restrepo (2022)]{acemoglu2022ai}
Acemoglu, Daron, David Autor, Jonathon Hazell, and Pascual Restrepo. 2022.
``Artificial Intelligence and Jobs: Evidence from Online Vacancies.''
\textit{Journal of Labor Economics}, 40(S1): S293--S340.

\bibitem[Agrawal, Gans, and Goldfarb (2018)]{agrawal2018prediction}
Agrawal, Ajay, Joshua Gans, and Avi Goldfarb. 2018.
``Prediction, Judgment, and Complexity.''
NBER Working Paper No.~24253.

\bibitem[Bolton and Dewatripont (1994)]{bolton1994firm}
Bolton, Patrick, and Mathias Dewatripont. 1994.
``The Firm as a Communication Network.''
\textit{Quarterly Journal of Economics}, 109(4): 809--839.

\bibitem[Brynjolfsson and Hitt (2000)]{brynjolfsson2000beyond}
Brynjolfsson, Erik, and Lorin Hitt. 2000.
``Beyond Computation: Information Technology, Organizational Transformation and Business Performance.''
\textit{Journal of Economic Perspectives}, 14(4): 23--48.

\bibitem[Caplin and Dean (2015)]{caplin2015revealed}
Caplin, Andrew, and Mark Dean. 2015.
``Revealed Preference, Rational Inattention, and Costly Information Acquisition.''
\textit{American Economic Review}, 105(7): 2183--2203.

\bibitem[Dessein and Santos (2006)]{dessein2006adaptive}
Dessein, Wouter, and Tano Santos. 2006.
``Adaptive Organizations.''
\textit{Journal of Political Economy}, 114(5): 956--995.

\bibitem[Dessein, Galeotti, and Santos (2016)]{dessein2016rational}
Dessein, Wouter, Andrea Galeotti, and Tano Santos. 2016.
``Rational Inattention and Organizational Focus.''
\textit{American Economic Review}, 106(6): 1522--1536.

\bibitem[Ely (2017)]{ely2017beeps}
Ely, Jeffrey C. 2017.
``Beeps.''
\textit{American Economic Review}, 107(1): 31--53.

\bibitem[Escud{\'e} and Sinander (2023)]{escude2023slow}
Escud{\'e}, Matteo, and Ludvig Sinander. 2023.
``Slow Persuasion.''
\textit{Theoretical Economics}, 18(1): 129--162.

\bibitem[Garicano (2000)]{garicano2000hierarchies}
Garicano, Luis. 2000.
``Hierarchies and the Organization of Knowledge in Production.''
\textit{Journal of Political Economy}, 108(5): 874--904.

\bibitem[Hadfield-Menell (2021)]{hadfield2021principal}
Hadfield-Menell, Dylan. 2021.
``The Principal-Agent Alignment Problem in Artificial Intelligence.''
PhD diss., University of California, Berkeley.

\bibitem[Holmstr{\"o}m and Milgrom (1991)]{holmstrom1991multitask}
Holmstr{\"o}m, Bengt, and Paul Milgrom. 1991.
``Multitask Principal-Agent Analyses.''
\textit{Journal of Law, Economics, \& Organization}, 7: 24--52.

\bibitem[Ide and Talamas (2025)]{ide2025ai}
Ide, Enrique, and Eduard Talam{\`a}s. 2025.
``Artificial Intelligence in the Knowledge Economy.''
\textit{Journal of Political Economy}, 133(2): 3762--3800.

\bibitem[Kamenica and Gentzkow (2011)]{kamenica2011bayesian}
Kamenica, Emir, and Matthew Gentzkow. 2011.
``Bayesian Persuasion.''
\textit{American Economic Review}, 101(6): 2590--2615.

\bibitem[Lash and Wellington (2008)]{lash2008information}
Lash, Alex, and David Wellington. 2008.
``Information Processing in Hierarchical Organizations.''
\textit{Management Science}, 54(8): 1423--1439.

\bibitem[Lien, Qiu, and Zhong (2019)]{lien2019hierarchies}
Lien, Jaimie W., Xincheng Qiu, and Xiaolan Zhong. 2019.
``Hierarchies and Information Processing in Teams.''
\textit{RAND Journal of Economics}, 50(3): 671--696.

\bibitem[Mackowiak, Mat{\v{e}}jka, and Wiederholt (2023)]{mackowiak2023rational}
Mackowiak, Bartosz, Filip Mat{\v{e}}jka, and Mirko Wiederholt. 2023.
``Rational Inattention: A Review.''
\textit{Journal of Economic Literature}, 61(1): 226--273.

\bibitem[Melumad, Mookherjee, and Reichelstein (1995)]{melumad1995hierarchical}
Melumad, Nahum D., Dilip Mookherjee, and Stefan Reichelstein. 1995.
``Hierarchical Decentralization of Incentive Contracts.''
\textit{RAND Journal of Economics}, 26(4): 654--672.

\bibitem[Noy and Zhang (2023)]{noy2023experimental}
Noy, Shakked, and Whitney Zhang. 2023.
``Experimental Evidence on the Productivity Effects of Generative Artificial Intelligence.''
\textit{Science}, 381(6654): 187--192.

\bibitem[Peng et al.~(2023)]{peng2023impact}
Peng, Sida, Eirini Kalliamvakou, Peter Cihon, and Mert Demirer. 2023.
``The Impact of AI on Developer Productivity: Evidence from GitHub Copilot.''
arXiv preprint arXiv:2302.06590.

\bibitem[Phelps and Ranson (2023)]{phelps2023models}
Phelps, Steve, and Tom Ranson. 2023.
``Of Models and Tin Men: A Multi-Task Principal-Agent Analysis of Large Language Models.''
\textit{Journal of Artificial Intelligence Research}, 77: 1123--1156.

\bibitem[Radner (1993)]{radner1993organization}
Radner, Roy. 1993.
``The Organization of Decentralized Information Processing.''
\textit{Econometrica}, 61(5): 1109--1146.

\bibitem[Renault, Solan, and Vieille (2017)]{renault2017optimal}
Renault, J{\'e}r{\^o}me, Eilon Solan, and Nicolas Vieille. 2017.
``Optimal Dynamic Information Provision.''
\textit{Games and Economic Behavior}, 104: 329--349.

\bibitem[Sims (2003)]{sims2003implications}
Sims, Christopher A. 2003.
``Implications of Rational Inattention.''
\textit{Journal of Monetary Economics}, 50(3): 665--690.

\bibitem[Sims (2006)]{sims2006rational}
Sims, Christopher A. 2006.
``Rational Inattention: Beyond the Linear-Quadratic Case.''
\textit{American Economic Review}, 96(2): 158--163.

\bibitem[Smolin (2021)]{smolin2021dynamic}
Smolin, Alex. 2021.
``Dynamic Evaluation Design.''
\textit{American Economic Journal: Microeconomics}, 13(4): 300--331.

\bibitem[Vaswani et al.~(2017)]{vaswani2017attention}
Vaswani, Ashish, Noam Shazeer, Niki Parmar, Jakob Uszkoreit, Llion Jones, Aidan N. Gomez, Lukasz Kaiser, and Illia Polosukhin. 2017.
``Attention Is All You Need.''
\textit{Advances in Neural Information Processing Systems}, 30: 5998--6008.

\bibitem[Williamson (1985)]{williamson1985economic}
Williamson, Oliver E. 1985.
\textit{The Economic Institutions of Capitalism}. New York: Free Press.

\bibitem[Yang and Zhai (2025)]{yang2025ten}
Yang, Zonghao, and Liwen Zhai. 2025.
``Ten Principles of AI Agent Economics.''
arXiv:2505.20273.

\end{thebibliography}


\end{document}
